\documentclass[11pt,a4paper]{article}
\usepackage[T1]{fontenc}
\usepackage[utf8]{inputenc}
\usepackage{lmodern}
\usepackage[a4paper,margin=25mm,headheight=15pt]{geometry}
\usepackage{amsmath,amssymb,amsthm,mathtools}
\usepackage{microtype,booktabs,longtable,array,tabularx}
\usepackage{xcolor,enumitem,fancyhdr,titlesec}
\usepackage{tocloft}
\setlength{\cftbeforesecskip}{3pt}
\usepackage{listings,xurl,hyperref,bookmark}
\definecolor{ink}{HTML}{17324D}
\definecolor{accent}{HTML}{245D83}
\definecolor{pale}{HTML}{F2F5F8}
\hypersetup{colorlinks=true,linkcolor=accent,citecolor=accent,urlcolor=accent,
 pdftitle={Specialization-Safe Radical Solvers},
 pdfauthor={Research prepared for Vladimir Reshetnikov},
 pdfsubject={Exact bad characteristics, subset resolvents, Fourier-Kummer charts, and a global obstruction; Part II: optimal Kummer atlases}}
\titleformat{\section}{\Large\bfseries\color{ink}}{\thesection}{0.65em}{}
\titleformat{\subsection}{\large\bfseries\color{ink}}{\thesubsection}{0.65em}{}
\pagestyle{fancy}\fancyhf{}
\fancyhead[L]{\small Specialization-safe radical solvers}
\fancyhead[R]{\small Research article}
\fancyfoot[C]{\thepage}
\setlength{\parskip}{3.5pt}
\setlength{\parindent}{1.25em}
\setlength{\emergencystretch}{2em}
\allowdisplaybreaks
\setlist{itemsep=3pt,topsep=5pt}
\lstset{basicstyle=\ttfamily\small,backgroundcolor=\color{pale},
 frame=single,rulecolor=\color{black!20},breaklines=true,
 columns=fullflexible,keepspaces=true,showstringspaces=false}
\newtheorem{theorem}{Theorem}[section]
\newtheorem{lemma}[theorem]{Lemma}
\newtheorem{proposition}[theorem]{Proposition}
\newtheorem{corollary}[theorem]{Corollary}
\theoremstyle{definition}
\newtheorem{definition}[theorem]{Definition}
\newtheorem{example}[theorem]{Example}
\newtheorem{question}[theorem]{Research question}
\theoremstyle{remark}
\newtheorem{remark}[theorem]{Remark}
\newcommand{\Q}{\mathbb Q}
\newcommand{\Z}{\mathbb Z}
\newcommand{\F}{\mathbb F}
\newcommand{\C}{\mathbb C}
\newcommand{\Gal}{\operatorname{Gal}}
\newcommand{\Res}{\operatorname{Res}}
\newcommand{\Spec}{\operatorname{Spec}}
\newcommand{\Pic}{\operatorname{Pic}}
\newcommand{\AGL}{\operatorname{AGL}}
\newcommand{\supp}{\operatorname{supp}}
\newcommand{\charac}{\operatorname{char}}
\newcommand{\Norm}{\operatorname{N}}
\newcommand{\disc}{\operatorname{disc}}
\newcommand{\Orb}{\operatorname{Orb}}
\newcommand{\inv}{^{-1}}
\newcommand{\code}[1]{{\small\nolinkurl{#1}}}
\newcommand{\snapshot}{\nolinkurl{e21766d04c2b8a9b2cdba0cd43e563024b3bd1b9}}
\newcolumntype{L}[1]{>{\raggedright\arraybackslash}p{#1}}
% [write] Unnumbered notes added when this manuscript was written into the
% ProveIt research-report collection (batch 36).  They use no theorem counter,
% so every theorem, equation and table number of the delivered text is kept.
\newenvironment{writenote}[1][]{\par\smallskip\begingroup\small
 \noindent\textbf{[write]}\if\relax\detokenize{#1}\relax\else~\textbf{#1.}\fi\ }%
 {\par\endgroup\smallskip}
\newcommand{\repo}[1]{\code{#1}}
% [write] Batch 52: package and macros of Part II (manuscript 09), as delivered.
\usepackage{needspace}
\newcommand{\A}{\mathbb A}
\newcommand{\Pj}{\mathbb P}
\newcommand{\OO}{\mathcal O}
\newcommand{\LL}{\mathcal L}
\newcommand{\EE}{\mathcal E}
\newcommand{\CH}{\operatorname{CH}}
\newcommand{\Tot}{\operatorname{Tot}}
\newcommand{\Conf}{\operatorname{Conf}}
\newcommand{\im}{\operatorname{im}}
\newcommand{\codim}{\operatorname{codim}}
\newcommand{\ord}{\operatorname{ord}}
\newcommand{\rk}{\operatorname{rk}}
\newcommand{\Sym}{\operatorname{Sym}}
\newcommand{\atlas}{\kappa}
\newcommand{\gen}{\gamma}
\newcommand{\alggen}{\alpha}
\newcolumntype{Y}{>{\raggedright\arraybackslash}X}

\begin{document}
\begin{center}
{\small\sffamily\color{accent} PROVEIT RESEARCH COMPANION\par}
\vspace{0.45cm}
{\LARGE\bfseries\color{ink} Specialization-Safe Radical Solvers\par}
\vspace{0.25cm}
{\large Exact Bad Characteristics for Subset-Sum Resolvents,\par
 a Fourier--Kummer Atlas, and a Global Obstruction\par}
\vspace{0.5cm}
{Research prepared for Vladimir Reshetnikov\par}
{September 2026\par}
\vspace{0.3cm}
{\footnotesize ProveIt research-report collection, continuing the formal
project \code{Algebra/PolynomialFormulas}. Batch~36, manuscript~04
(archive \code{ProveIt_Radical_Solvers_Research}); pinned to ProveIt
\code{e21766d04}. AI-assisted, unrefereed, not formalized.
Notes marked \textbf{[write]} were added when the manuscript was written into
the collection; see Section~\ref{rss:sec:provenance}.\par
\smallskip
Part~II (added 29 September 2026), \emph{Optimal Kummer atlases for prime
cyclic configurations}: batch~52, manuscript~09 (archive
\code{ProveIt_Optimal_Kummer_Atlases}); pinned to ProveIt \code{1ee53d57d};
see Section~\ref{ka:sec:provenance}.\par}
\end{center}
\begin{abstract}
A separable polynomial can have an inseparable auxiliary resolvent; a correct
radical identity can become unusable on a nonempty exceptional locus; and
separately chosen radical branches need not reconstruct one compatible root
tuple. We study these three problems in the prime-degree setting of the
ProveIt polynomial-formula project.

For equal-cardinality subset sums of the roots of an irreducible prime-degree
polynomial, we give an exact, finite cyclotomic-resultant criterion for the
characteristics in which collisions are possible. It yields the complete
exceptional sets $\{5\}$ for quintic pair sums, $\{2,7\}$ for septic pair sums,
and $\{2,7,29\}$ for septic triple sums. In particular, over $\F_{29}(t)$ the
separable irreducible polynomial $X^7-t$ has triple-sum resolvent
\[
 (Y^7-12t)^2(Y^7-t)(Y^7-17t)(Y^7-28t).
\]
In characteristic $29$, collisions occur precisely for translated pure
septics; in characteristic two they characterize quotients of additive
degree-eight polynomials, up to translation. The explicit example supplies
a concrete obstruction to an unrestricted distinct-factor
reading of a published septic factor-pattern assertion. A polynomial
separating-invariant construction repairs the specialization problem over
any sufficiently large field; for septic triples, 1191 distinct parameter
values always suffice.

We then prove a branch-coherent Fourier--Kummer reconstruction using one
$p$th root on each nonvanishing Fourier chart. Every nonempty Fourier support
occurs for an irreducible rational polynomial with full affine Galois group.
Finally, the cyclic quotient of the ordered distinct-root configuration space
has Picard group $\Z/p\Z$, and its universal cyclic cover has no global regular
Kummer generator. Thus local formulas are not merely an artifact of a poor
choice of coordinates. Complete arguments, finite integer certificates,
reproducible exact checks, a repository integration plan, and further
research questions are included.
\end{abstract}

\medskip
\noindent\textbf{Status.} This is a mathematical research article with exact
computational certificates, not a newly kernel-checked Lean development.
The basic Fourier, Galois, resultant, and descent tools are classical. The
specific classifications, examples, and synthesis below are proved here;
priority over all existing literature has not been established. No claim to
have solved the general problem of finding new radical formulas in every
degree is made.

\clearpage
\begingroup\small
\tableofcontents
\endgroup
\clearpage

\phantomsection\addcontentsline{toc}{part}{Part I.\texorpdfstring{\quad}{ } Specialization-safe radical solvers}
\section{Research target, provenance, and contribution boundary}
\label{rss:sec:boundary}

\subsection{Why this problem arises in ProveIt}

The repository snapshot used for the investigation is commit
\begin{center}\snapshot.\end{center}
The root README identifies the audit of Lazard's solvable-quintic algorithm
as an active direction. The polynomial-formula report distinguishes raw
invariant equations from root-origin or Fourier evidence, and explicitly
records exceptional denominators and the need for coherent radical choices
\cite{repoReadme,repoReport,repoCrosswalk}. These are mathematical issues, not
only proof-assistant engineering issues.

The inspected source \code{LazardGeneralFourier.lean} already proves
arbitrary-degree character orthogonality, cyclic-shift covariance, and
Fourier inversion. It deliberately stops before branch-sensitive recovery
of modes from invariants \cite{repoFourier}. Accordingly, ordinary discrete
Fourier inversion is not presented below as a new result. Our reconstruction
uses that existing interface and adds a pivot chart, invariant ratios,
and a proof that every choice of the one remaining radical changes only the
cyclic ordering of the roots.

\begin{writenote}[Repository context the manuscript did not cite]
Checked against the tree, which is unchanged in
\code{Algebra/PolynomialFormulas} between the pin and batch~36.
\begin{itemize}
\item \emph{The Rocq twin.} Beside the Lean module
\repo{Algebra/PolynomialFormulas/Lean/PolynomialFormulas/LazardGeneralFourier.lean}
(\code{character_orthogonality}, l.~49; \code{dft_shift}, l.~57;
\code{invDFT_dft}, l.~98; the scope boundary quoted above is its docstring,
l.~18--21) the project has
\repo{Algebra/PolynomialFormulas/Coq/LazardGeneralFourier.v}, proving the
same arbitrary-degree core over any MathComp field:
\code{lazard_general_character_orthogonality} (l.~96),
\code{lazard_general_dft_shift} (l.~128) and
\code{lazard_general_invDFT_dft} (l.~207), with the coordinatewise
\code{lazard_general_invDFT_dft_coordinate} (l.~179)
\cite{repoFourierRocq}. The claim crosswalk
(\repo{Algebra/PolynomialFormulas/LazardPaperClaimCrosswalk.md}, l.~199) records both focused modules
as kernel-green with their public aggregate/audit pending, and states the
same boundary: arbitrary-degree branch-sensitive invariant recovery is
deliberately outside the general Fourier modules, while the quintic
invariant recovery is covered by a separate five-coordinate chain.
\item \emph{Lazard's own one-radical scheme.} The pivot chart and invariant
ratios are not new in the quintic case. The project report
\repo{Algebra/PolynomialFormulas/LazardQuinticFormalization.tex}
(l.~606--612) records Lazard's
invariants $S_1=s_1^5$, $S_2=s_2s_1^3$, $S_3=s_3s_1^2$, $S_4=s_4s_1$ and his
reconstruction ``once $P_1^5=S_1$ is chosen with $P_1\neq0$'', with
$P_2=S_2/P_1^3$, $P_3=S_3/P_1^2$, $P_4=S_4/P_1$ \cite{repoReport}. This is
the case $p=5$, pivot $j=1$ of Theorem~\ref{rss:thm:chart}: with
$\zeta=\omega^{-1}$ one has $u_k=s_k$, $a_1=S_1$, $c_{1,k}a_1=S_k$ and
$\rho=P_1$ (dictionary in Section~\ref{rss:sec:notation}). What the chart
theorem adds is the arbitrary odd prime $p$, the arbitrary pivot $j$, the
statement that \emph{every} branch of the one radical yields the cyclic
shift $x_{i+rj^{-1}}$, the transition rules on overlaps, and
Corollary~\ref{rss:cor:generator}. A name search of the Lean quintic modules
at the write (for \code{shift} and \code{rotate}) found no theorem that every
fifth-root choice gives a cyclic shift; that search is not a proof of
absence.
\end{itemize}
\end{writenote}

Three precise research problems guide the article. First, determine exactly
when a standard subset-sum resolvent can fail to separate subsets despite
separability and irreducibility of the original polynomial. Second, replace
an unsafe invariant by one with a proved specialization guarantee. Third,
understand whether one can eliminate exceptional radical charts globally,
or whether there is a structural obstruction. These are the questions
formulated and answered here; they are not attributed to an existing issue
or to a named open-problem list.

\subsection{Relation to the published formulas}

Lazard's 2004 chapter gives explicit formulas for solvable quintics
\cite{Lazard2004}. Faggal and Lazard's 2014 article also treats septics, using
sums of three roots and their factorization patterns; its specialized septic
construction assumes characteristic different from $2,3,7$
\cite{FaggalLazard2014}. Thus the existence of a radical treatment of solvable
septics is not a new claim of this article.

Our characteristic-$29$ example concerns the distinction between an orbit
product and a squarefree resolvent with distinct conjugate values. It does
not, on its own, disprove an abstract radical-solvability theorem. In
particular, we do not infer that every later identity in that article fails.
Section~\ref{rss:sec:published} states the precise implication and its limits.

\subsection{What is established, and by what kind of evidence}

\begin{center}
\small
\begin{tabularx}{\textwidth}{@{}L{0.26\textwidth}XL{0.22\textwidth}@{}}
\toprule
Result & Mathematical content & Evidence here \\
\midrule
Exact bad-characteristic criterion & Necessary and sufficient cyclotomic norm
condition, with irreducible witnesses & General proof \\
Small-degree classification & Complete exceptional sets for $(p,k)=(5,2),(7,2),(7,3)$
& Proof plus finite integer tables \\
Characteristic-$29$ witness & Explicit repeated-factor triple-sum resolvent
& Direct field calculation \\
Septic collision rigidity & Complete root-configuration classification in
characteristics $2$ and $29$ & General proof and finite norm certificate \\
Separating repair & Characteristic-independent subset descriptors and a
finite parameter bound & General proof \\
One-radical charts & Coherent reconstruction and transition functions
& General proof \\
Support realization & Every nonempty support with full affine Galois group
& Explicit rational construction \\
Global obstruction & $\Pic(A^{C_p})\cong\Z/p\Z$ and no eigenunit
& Algebraic descent and unit calculation \\
Lean integration & Proposed declarations and dependencies
& Plan only; no new Lean build \\
\bottomrule
\end{tabularx}
\end{center}

The universal theorems do not depend on numerical experiments. The three
small-degree classifications use a fully explicit finite calculation: all
850 unordered pairs are checked by exact resultants, and the 29 affine-orbit
representatives are checked independently by integer determinants. The
provided standard-library auditor checks the finite certificates without
SymPy. This is still an ordinary computational proof boundary, not a
proof-assistant kernel boundary.

\subsection{[write] Conventions and notation}
\label{rss:sec:notation}

No symbol of the delivered manuscript was renamed and no normalization was
changed. Several letters carry more than one meaning, inside the article or
between the article and the formal project it continues; each is
unambiguous in context, and the table fixes the readings.

\begin{center}
\small
\begin{tabularx}{\textwidth}{@{}L{0.14\textwidth}>{\raggedright\arraybackslash}X>{\raggedright\arraybackslash}X@{}}
\toprule
Symbol & Meaning here & Not to be confused with\\
\midrule
$u_j$, $\zeta$ &
Fourier mode $u_j=\sum_i\zeta^{-ij}x_i$ (negative exponent),
\eqref{rss:eq:fourier} &
Lazard's $s_k=\sum_j\omega^{kj}x_j$ (positive exponent), in the project
report and in \code{LazardGeneralFourier.lean}, whose character is called
$\psi$. \emph{True:} with $\zeta=\omega^{-1}$, $u_k=s_k$.
\emph{False reading:} $u_k=s_k$ for $\zeta=\omega$; then $u_k=s_{p-k}$.\\
$a_j$, $c_{j,k}$, $\rho$ &
Pivot power $u_j^p$, invariant ratios $u_k/u_j^{e_j(k)}$, a root of
$T^p-a_j$ (Theorem~\ref{rss:thm:chart}) &
For $p=5$, $j=1$: $a_1=S_1$, $c_{1,k}a_1=S_k$, $\rho=P_1$ in Lazard's
notation \cite{repoReport}; $c_{1,k}$ itself is \emph{not} $S_k$.\\
$e_j(k)$ &
The integer in $\{1,\dots,p-1\}$ with $je_j(k)\equiv k$ (Section~\ref{rss:sec:atlas}) &
$e_r$, the $r$th elementary symmetric function (Sections~\ref{rss:sec:witnesses}, \ref{rss:sec:repair}).\\
$\mathcal B_{p,k}$ &
Set of bad primes (Section~\ref{rss:sec:criterion}) &
$B_{n,k}=(k-1)\binom N2$, an integer bound \eqref{rss:eq:bound}.\\
$A$, $B$ &
Subsets of indices (Sections~\ref{rss:sec:subset}--\ref{rss:sec:repair}) &
$B$: the fixed field $E^{\langle\sigma\rangle}$ (Section~\ref{rss:sec:atlas}),
$B_0$, $B$ in Section~\ref{rss:sec:groups}; $A$, $B=A^{C_p}$: rings in
Section~\ref{rss:sec:global}.\\
$E$ &
A field (Sections~\ref{rss:sec:atlas}, \ref{rss:sec:groups}); $E(A)$: the
descriptor vector (Section~\ref{rss:sec:repair}) &
Lazard's invariant denominator $E$ in the project report.\\
$k$ &
Subset size (Sections~\ref{rss:sec:subset}--\ref{rss:sec:repair}); mode index
in $c_{j,k}$; algebraically closed field (Section~\ref{rss:sec:global}) & ---\\
$N$ &
Norm $N_{A,B}$ \eqref{rss:eq:norm}; $N=\binom nk$ (Section~\ref{rss:sec:repair});
a minimal normal subgroup (Lemma~\ref{rss:lem:affine}) & ---\\
$f$, $g$, $m$ &
The polynomial; but $f=e_h(k)$, $g=e_j(k)$ and $m=(ef-g)/p$ in
\eqref{rss:eq:transitioncoeff} &
$m_{ij}$, entries of the multiplication matrix $M_q$.\\
$T$ &
Polynomial variable (Sections~\ref{rss:sec:criterion}, \ref{rss:sec:tables}),
Kummer variable in $T^p-a_j$ &
Lazard's quadratic-stage quantity $T$ in the project report.\\
$\psi$ &
The $\F_2$-linear map in the proof of Theorem~\ref{rss:thm:2rigidity} &
The additive character $\psi$ of \code{LazardGeneralFourier.lean}.\\
$L$ &
Splitting field; $L_W$ additive root polynomial; $L_r$ character
eigenspace modules (Section~\ref{rss:sec:global}) & ---\\
$\Delta$ &
Discriminant $\prod_{i<j}(x_i-x_j)^2$ \eqref{rss:eq:configuration} & ---\\
$c$ &
$c_r(A)$ (Theorem~\ref{rss:thm:29}); $c_{j,k}$; $c_k$ (Theorem~\ref{rss:thm:support});
coefficient vector (Section~\ref{rss:sec:criterion}); a constant in \eqref{rss:eq:units} & ---\\
\bottomrule
\end{tabularx}
\end{center}

\subsection{[write] Provenance of this report}
\label{rss:sec:provenance}

This report is built from \textbf{one} manuscript: batch~36, manuscript~04 of
ProveIt's incoming-reports intake, delivered as the archive
\code{ProveIt_Radical_Solvers_Research} in commit \code{e13affd32}, pinned to
ProveIt commit \code{e21766d04c2b8a9b2cdba0cd43e563024b3bd1b9}, and placed in
commit \code{1a1396d4d}. It is filed in ProveIt's research-report collection
and continues the formal project \code{Algebra/PolynomialFormulas}; it shares
no theorem with any other report. Every statement, proof, table, example,
remark, question and limitation of the delivered text is printed. No merge
choice arose. (This paragraph describes Part~I. Part~II, added on
29 September 2026, is built from a second manuscript, batch~52
manuscript~09; its provenance is Section~\ref{ka:sec:provenance}.)

The write changed only the following. Every label carries the prefix
\code{rss:}. The title-page provenance line, Sections~\ref{rss:sec:notation},
\ref{rss:sec:provenance} and~\ref{rss:sec:nonclaims}, the unnumbered
\textbf{[write]} notes, and the bibliography entries \cite{repoFourierRocq},
\cite{repoSextic}, \cite{repoKummer} and \cite{repoCounterexamples} were
added. The notes use no theorem counter, so every theorem, equation and
table number is that of the delivered 29-page PDF. The delivered text names
the delivery layout: its files \code{certificates/}\emph{X} are shipped as
\code{data/}\emph{X}, its two scripts as \code{code/verify_results.py} and
\code{code/audit_certificates.py}, and its \code{requirements.txt} as
\code{data/requirements.txt}; the delivered PDF and README are not shipped.
Appendix~\ref{rss:app:repro} gives commands that work with the shipped names.

\subsection{[write] Formal status and what is not claimed}
\label{rss:sec:nonclaims}

\begin{itemize}
\item \textbf{Not formalized.} No statement of this report is formalized in
Lean or Rocq, and no \code{rss:} label has a Lean or Rocq counterpart. That
the report continues the \code{Algebra/PolynomialFormulas} development
confers no formal status on it. The project's own nearby declarations are named where they are used:
the arbitrary-degree Fourier core in Lean and Rocq (Section~\ref{rss:sec:boundary}),
sextic descriptor separation (Section~\ref{rss:sec:repair}), the cyclic Kummer
generator (Section~\ref{rss:sec:atlas}) and two positive-characteristic
counterexamples (Section~\ref{rss:sec:criterion}). None of them states a
theorem of this report.
\item \textbf{Integration plan.} Section~\ref{rss:sec:formalization} is a
plan. None of it is formalized, and no Lean or Rocq code is delivered or
shipped; its declaration names are proposals.
\item \textbf{Evidence.} The general theorems rest on written proofs. The
finite tables rest on exact computations by an ordinary Python program and an
independent standard-library auditor, not a proof-assistant kernel. The
Fourier checks cover coefficient-one supports only. Nothing computational
checks the Picard-group and eigenunit section (Section~\ref{rss:sec:global}),
the transition rule \eqref{rss:eq:transitioncoeff}, the Galois groups
$\AGL_1(\F_p)$, or the rigidity theorems beyond their finite character-pattern
counts.
\item \textbf{The source's own non-claims} stand as printed: the Status
paragraph after the abstract (not a kernel-checked development; classical
tools are not new; priority over the literature not established; no general
radical formulas in every degree); the limits of Section~\ref{rss:sec:published}
(distinct-factor separation only, not a complete audit of the published
algorithm); the assumed invariant data of Theorem~\ref{rss:thm:chart}
(Section~\ref{rss:sec:atlas}); relative minimality only in
Corollary~\ref{rss:cor:generator}; minimality of the coordinate-pivot cover
only in Corollary~\ref{rss:cor:charts}; what Theorem~\ref{rss:thm:eigenunit}
does not rule out; and the status of the research questions
(Section~\ref{rss:sec:questions}).
\item \textbf{Review.} The report is AI-assisted and unrefereed. At the intake
the proofs were read, the headline computations were repeated by a separate
script, and both delivered suites were rerun on a copy; this is not an
independent proof review. Parts of Section~\ref{rss:sec:criterion} (the
relation-module viewpoint and characteristic-zero separation) may be close
to the classical literature on linear relations among conjugates; this was
not checked, and no priority is claimed.
\item \textbf{Part~II} (added 29 September 2026) answers Research
question~\ref{ka:q:rss-minimal-covers}. It is likewise AI-assisted,
unrefereed and not formalized; its own status and non-claims are
Section~\ref{ka:sec:nonclaims}.
\end{itemize}

\section{Subset-sum resolvents and the correct notion of separation}
\label{rss:sec:subset}

Let $p$ be an odd prime and let $f\in K[X]$ be monic, separable, and
irreducible of degree $p$. In its splitting field choose a numbering
$x_0,\ldots,x_{p-1}$ of its roots. For $1\leq k\leq p-1$ put
\[
 \mathcal S_{p,k}=\{A\subseteq\{0,\ldots,p-1\}:|A|=k\},\qquad
 s_A=\sum_{i\in A}x_i,
\]
and define the \emph{orbit product}
\begin{equation}
 R_{p,k,f}(Y)=\prod_{A\in\mathcal S_{p,k}}(Y-s_A).
 \label{rss:eq:resolvent}
\end{equation}
Galois automorphisms permute the subsets, so this monic polynomial belongs
to $K[Y]$. Its degree, with multiplicity, is always $\binom pk$.

\begin{definition}
The $k$-subset sums of $f$ \emph{separate} if $s_A\neq s_B$ for every two
distinct members of $\mathcal S_{p,k}$. Equivalently,
$R_{p,k,f}$ is separable, or $\gcd(R_{p,k,f},R'_{p,k,f})=1$.
\end{definition}

This is stronger than separability of $f$. The latter distinguishes the
individual $x_i$, not all sums made from them. It is also important not to
call \eqref{rss:eq:resolvent} the minimal polynomial of a single value without
qualification. If the Galois group has several orbits on subsets, then even
a separating orbit product is reducible. If distinct subsets have the same
value, the product has repeated factors.

\begin{lemma}[Orbit-to-factor principle]\label{rss:lem:orbitfactor}
Let a finite Galois group $G$ act on a finite set $\mathcal S$, and let
$A\mapsto v_A$ be a $G$-equivariant map into the Galois extension. If this
map is injective, the irreducible factors over the fixed field of
$\prod_{A\in\mathcal S}(Y-v_A)$ are indexed by the $G$-orbits on
$\mathcal S$, and their degrees equal the orbit sizes.
\end{lemma}
\begin{proof}
The conjugates of $v_A$ are exactly $v_{gA}$ for $g\in G$. Injectivity
identifies its stabilizer with that of $A$. Its minimal polynomial is
therefore the product over the orbit of $A$, with distinct roots. Different
orbits have disjoint value sets, and multiplication gives the assertion.
\end{proof}

Without injectivity, two different subset orbits can produce the same
irreducible factor. This is exactly the mechanism in the characteristic-$29$
example, not a loss of separability of the original degree-seven equation.

Complementation gives a useful symmetry:
\[
 s_{A^c}=\operatorname{Tr}(x_0)-s_A,
 \qquad \operatorname{Tr}(x_0)=\sum_{i=0}^{p-1}x_i.
\]
Consequently separation for $k$ is equivalent to separation for $p-k$.
The cases $k=1,p-1$ always separate. The interesting cases are
$2\leq k\leq p-2$.

\section{An exact cyclotomic criterion for bad characteristics}
\label{rss:sec:criterion}

For distinct $A,B\in\mathcal S_{p,k}$ define
\begin{equation}
 q_{A,B}(T)=\sum_{i\in A}T^i-\sum_{i\in B}T^i,
 \qquad \Phi_p(T)=1+T+\cdots+T^{p-1},
 \label{rss:eq:q}
\end{equation}
and the positive integer
\begin{equation}
 N_{A,B}=\left|\Res_T(\Phi_p,q_{A,B})\right|.
 \label{rss:eq:norm}
\end{equation}
Equivalently, $N_{A,B}$ is the absolute norm of $q_{A,B}(\zeta)$ from
$\Q(\zeta)$, for a primitive $p$th root of unity $\zeta$.

\begin{lemma}\label{rss:lem:normnonzero}
Every integer $N_{A,B}$ in \eqref{rss:eq:norm} is nonzero. It is divisible by $p$.
\end{lemma}
\begin{proof}
The polynomial $q=q_{A,B}$ is nonzero, has degree at most $p-1$, and
$q(1)=0$. If $q$ were divisible by $\Phi_p$ over $\Q$, it would equal
$c\Phi_p$ for some constant $c$, and $0=q(1)=cp$ would force $c=0$.
Irreducibility of $\Phi_p$ over $\Q$ therefore gives a nonzero resultant.
Modulo $p$, $\Phi_p=(T-1)^{p-1}$ and $q(1)=0$, so the reductions have a
common root. Their resultant is zero modulo $p$.
\end{proof}

\begin{definition}
For $2\leq k\leq p-2$, let $\mathcal B_{p,k}$ be the set of prime divisors
of the integers $N_{A,B}$, as $\{A,B\}$ runs over all unordered pairs of
distinct $k$-subsets. It is a finite, effectively computable set.
\end{definition}

\begin{theorem}[Exact universal separation criterion]\label{rss:thm:bad}
Let $p$ be an odd prime and $2\leq k\leq p-2$.
\begin{enumerate}[label=(\roman*)]
\item In characteristic zero, the $k$-subset sums of every separable
irreducible degree-$p$ polynomial separate.
\item In characteristic $\ell\neq p$, the following are equivalent:
\begin{enumerate}[label=(\alph*)]
\item for every field $K$ of characteristic $\ell$ and every separable
irreducible degree-$p$ polynomial over $K$, the $k$-subset sums separate;
\item $\ell\notin\mathcal B_{p,k}$.
\end{enumerate}
\item In characteristic $p$, universal separation fails.
\end{enumerate}
Thus $\mathcal B_{p,k}$ is exactly the set of characteristics in which a
counterexample can exist.
\end{theorem}

\subsection{Proof of the safe direction}

Let $L/K$ be the splitting field of $f$, and $G=\Gal(L/K)$. Irreducibility
makes the root action transitive, so $p$ divides $|G|$. Cauchy's theorem
supplies an element of order $p$. A nonidentity order-$p$ permutation on
$p$ letters is a $p$-cycle. We may consequently number the roots so that
\[
 \sigma(x_i)=x_{i+1},\qquad i\pmod p.
\]
All pairs of subsets occur in the definition of $\mathcal B_{p,k}$, so
changing the numbering does not change the criterion.

Suppose that $s_A=s_B$, and write $q=q_{A,B}$. Then
\[
 q(\sigma)x_0=0.
\]
In characteristic zero, Lemma~\ref{rss:lem:normnonzero} gives
$\gcd(q,\Phi_p)=1$ in $\Q[T]$. In characteristic
$\ell\notin\mathcal B_{p,k}$, the resultant is nonzero modulo $\ell$, so
the same coprimality holds in $\F_\ell[T]$. In either case, over the prime
field there are polynomials $a,b$ with
\[
 a(T)q(T)+b(T)\Phi_p(T)=1.
\]
Apply the corresponding operator identity to $(\sigma-1)x_0$. Both terms
on the left vanish: the first because $q(\sigma)x_0=0$, and the second
because
\[
 \Phi_p(\sigma)(\sigma-1)=\sigma^p-1=0.
\]
Hence $(\sigma-1)x_0=0$, contradicting $x_1\neq x_0$. This proves (i) and
the implication (b)$\Rightarrow$(a) of (ii).

Notice that no irreducibility assertion about $\Phi_p$ over $K$ was used.
In particular, the safe direction remains valid when $K$ already contains
all $p$th roots of unity.

\subsection{Proof that every bad prime really occurs}

Suppose $\ell\neq p$ divides $N_{A,B}$. The reductions of $\Phi_p$ and
$q_{A,B}$ have a common root $\zeta$ in $\overline{\F}_\ell$. Since
$\ell\neq p$, this is a primitive $p$th root of unity. Put
\[
 K=\F_\ell(\zeta)(t),\qquad f(X)=X^p-t,
\]
where $t$ is transcendental. The polynomial is Eisenstein at $t$ in
$\F_\ell(\zeta)[t][X]$, hence irreducible over $K$. Its derivative is
$pX^{p-1}$, so it is separable. If $\alpha^p=t$, its roots are
$x_i=\alpha\zeta^i$, and
\[
 s_A-s_B=\alpha q_{A,B}(\zeta)=0.
\]
Moreover, $K(\alpha)/K$ is already its splitting field and has cyclic
Galois group of order $p$. Thus the unsafe characteristic is witnessed even
by a very simple solvable polynomial. This proves the converse in (ii).

For (iii), use the Artin--Schreier polynomial
\[
 f(X)=X^p-X-t\in\F_p(t)[X].
\]
It is separable. To see its irreducibility, first note that
$t\neq b^p-b$ for any $b\in\F_p(t)$: a pole of $b^p-b$ has order divisible
by $p$, whereas $t$ has a pole of order one at infinity. If $\alpha$ is a
root, all its roots $\alpha+i$ lie in $\F_p(t)(\alpha)$; the Galois group
embeds into the additive group $\F_p$ and therefore has order one or $p$.
The absence of a root in the base field rules out order one.

For a $k$-subset $A$, its sum is
\[
 k\alpha+\sum_{i\in A}i.
\]
There are at most $p$ such values, but $\binom pk>p$ in the stated range of
$k$. Some distinct subsets have equal sums. This finishes the proof of
Theorem~\ref{rss:thm:bad}. \hfill$\square$

\begin{remark}[The exact logical scope]
A prime in $\mathcal B_{p,k}$ does not mean that every polynomial in that
characteristic has a collision. It means that a collision is possible and
that a characteristic-only universal guarantee is false. A particular
resolvent can still be tested for separability by a gcd computation.
\end{remark}

\begin{writenote}[Related positive-characteristic material in the project]
Two formal developments of \code{Algebra/PolynomialFormulas} concern
neighbouring phenomena at the prime~$3$ \cite{repoCounterexamples}; neither
states a theorem of this section.
\begin{itemize}
\item \code{LazardAlternatingResolventCounterexample} (Lean and Rocq; claim
crosswalk l.~202) refutes the assertion that the alternating resolvent
$X^2-\Delta$ of an irreducible polynomial is always separable: over
$\F_3(t)$ the cubic $X^3-t$ is irreducible but \emph{inseparable}
(derivative zero), its discriminant vanishes, and its alternating resolvent
is $X^2$. There the source polynomial itself is inseparable. The
characteristic-$29$ phenomenon of Theorem~\ref{rss:thm:29} is different: the
septic $X^7-t$ is separable, and only the auxiliary resolvent acquires
repeated factors.
\item \code{LazardArtinSchreierRadicalCounterexample} (Lean and Rocq;
crosswalk l.~200) uses the Artin--Schreier cubic $X^3-X-t$, the case $p=3$ of
the witness $X^p-X-t$ in part~(iii) above, to show that a cyclic, hence
solvable, Galois group in characteristic~$3$ need not give solvability by
radicals. It concerns radical towers, not subset-sum collisions.
\end{itemize}
\end{writenote}

\subsection{A general relation-space interpretation}

The same proof applies to a prescribed coefficient vector
$c=(c_0,\ldots,c_{p-1})\in\Z^p$ with $\sum c_i=0$ and
$q_c(T)=\sum c_iT^i\neq0$. Outside the prime divisors of
$\Res(\Phi_p,q_c)$, the relation
$\sum c_i x_i=0$ is impossible for a separable irreducible degree-$p$
polynomial. Conversely, every non-$p$ prime divisor gives a pure-binomial
witness for that relation after adjoining the necessary roots of unity.
Thus the subset-sum theorem is one instance of an exact integral
relation-obstruction method, not merely a collection of examples.

For a fixed characteristic different from $p$, relations among the
cyclically ordered roots form a module over
$\F_\ell[T]/(T^p-1)$. The cyclotomic component
$\F_\ell[T]/(\Phi_p)$ detects the nonconstant part. Our resultant test asks
whether a prescribed small coefficient vector can vanish on this component.
This viewpoint suggests extensions to cyclic codes and more general
weighted resolvents, discussed in Section~\ref{rss:sec:questions}.

\section{Complete classifications in degrees five and seven}
\label{rss:sec:tables}

\begin{theorem}[Small-degree exceptional sets]\label{rss:thm:small}
The exact exceptional characteristics are
\[
 \boxed{\mathcal B_{5,2}=\{5\}},\qquad
 \boxed{\mathcal B_{7,2}=\{2,7\}},\qquad
 \boxed{\mathcal B_{7,3}=\{2,7,29\}}.
\]
The complementary subset sizes have the same exceptional sets.
\end{theorem}

We give the finite certificate behind this theorem explicitly. It is useful
both for independent review and for a future formalization.

\subsection{Reduction by the affine group}

The affine group $\AGL_1(\F_p)$ acts on subset labels by
$i\mapsto ai+b$, where $a\neq0$. For a primitive $p$th root $\zeta$,
\[
 q_{aA+b,aB+b}(\zeta)=\zeta^b q_{A,B}(\zeta^a).
\]
The cyclotomic field automorphism $\zeta\mapsto\zeta^a$ preserves absolute
norms, and $\zeta^b$ has norm one. Interchanging $A$ and $B$ changes only a
sign. Therefore $N_{A,B}$ is constant on affine orbits of unordered pairs.

For each pair, choose the lexicographically smallest representative of its
affine orbit. In the following tables, $012$ denotes $\{0,1,2\}$, and
similarly for other strings. The orbit sizes add to the number of all
unordered pairs; disjointness is checked by the displayed normalization rule.
Each norm can be computed as an integer determinant without algebraic
number arithmetic.

\subsection{The determinant certificate}

Write
\[
 q(T)T^j\equiv\sum_{i=0}^{p-2}m_{ij}T^i\pmod{\Phi_p(T)},
 \qquad 0\leq j\leq p-2.
\]
Multiplication by $q$ in $\Z[T]/(\Phi_p)$ has matrix $M_q=(m_{ij})$.
Its determinant is $\Res(\Phi_p,q)$. Reduction uses only
\[
 T^p=1,\qquad T^{p-1}=-(1+T+\cdots+T^{p-2}).
\]
These identities supply a particularly small exact certificate.

For the exceptional septic pair $A=012$, $B=346$, the matrix is
\begin{equation}
 M_q=\begin{pmatrix}
 2&-1&1&0&-2&0\\
 2&1&0&1&-2&-2\\
 2&1&2&0&-1&-2\\
 0&1&2&2&-2&-1\\
 0&-1&2&2&0&-2\\
 1&-1&0&2&0&0
 \end{pmatrix},\qquad |\det M_q|=203=7\cdot29.
 \label{rss:eq:matrix203}
\end{equation}
The integer matrices for every other row are included in
\code{certificates/cyclotomic_norm_certificates.json}.\footnote{[write] Shipped
in this report as \code{data/cyclotomic_norm_certificates.json}.}

\begin{center}
\begin{tabular}{@{}ccrr@{}}
\toprule
$A$ & $B$ & Orbit size & $N_{A,B}$\\
\midrule
01&02&20&5\\
01&04&10&5\\
01&23&10&5\\
01&24&5&25\\
\midrule
\multicolumn{2}{l}{Total, quintic pairs}&45&\\
\bottomrule
\end{tabular}
\qquad
\begin{tabular}{@{}ccrr@{}}
\toprule
$A$ & $B$ & Orbit size & $N_{A,B}$\\
\midrule
01&02&42&7\\
01&03&42&7\\
01&06&21&7\\
01&23&21&7\\
01&24&42&56\\
01&26&21&49\\
01&34&21&7\\
\midrule
\multicolumn{2}{l}{Total, septic pairs}&210&\\
\bottomrule
\end{tabular}
\end{center}

\begin{center}
\begin{longtable}{@{}rccrr@{}}
\caption{All affine orbits of unordered pairs of septic triples.}\label{rss:tab:triples}\\
\toprule
Row&$A$&$B$&Orbit size&$N_{A,B}$\\
\midrule\endfirsthead
\toprule Row&$A$&$B$&Orbit size&$N_{A,B}$\\\midrule\endhead
1&012&013&42&7\\
2&012&014&42&7\\
3&012&015&42&7\\
4&012&016&21&7\\
5&012&023&42&7\\
6&012&034&42&7\\
7&012&035&42&56\\
8&012&045&42&7\\
9&012&046&42&49\\
10&012&056&21&7\\
11&012&134&42&56\\
12&012&136&21&49\\
13&012&345&21&7\\
14&012&346&42&203\\
15&013&015&21&7\\
16&013&026&42&56\\
17&013&046&21&7\\
18&013&245&7&343\\
\midrule
\multicolumn{3}{l}{Total}&595&\\
\bottomrule
\end{longtable}
\end{center}

\begin{proof}[Proof of Theorem~\ref{rss:thm:small}]
The tables exhaust the pair sets of sizes
$\binom{10}{2}=45$, $\binom{21}{2}=210$, and $\binom{35}{2}=595$.
The norms occurring in the three cases are, respectively,
\[
 \{5,25\},\qquad \{7,49,56\},\qquad
 \{7,49,56,203,343\}.
\]
Their prime divisors are exactly the displayed exceptional sets.
Theorem~\ref{rss:thm:bad} supplies both the universal safe conclusion and the
existence of counterexamples in every exceptional characteristic.
\end{proof}

As an additional checksum on the finite enumeration, the complete norm
multiplicities are
\begin{align*}
 (5,2):&\quad 5^{[40]},\ 25^{[5]},\\
 (7,2):&\quad 7^{[147]},\ 49^{[21]},\ 56^{[42]},\\
 (7,3):&\quad 7^{[357]},\ 49^{[63]},\ 56^{[126]},\ 203^{[42]},\ 343^{[7]}.
\end{align*}
Here $n^{[m]}$ means that norm $n$ occurs $m$ times, not an exponentiation.
The independent auditor regenerates the orbits and multiplication matrices,
checks every determinant, and verifies that their union is exactly the full
pair set. No inference from a sample of pairs is being made.

\section{Explicit exceptional-characteristic witnesses}
\label{rss:sec:witnesses}

\subsection{The characteristic-29 septic}

Let $K=\F_{29}(t)$, where $t$ is transcendental, and let $\alpha^7=t$.
Take $\zeta=23\in\F_{29}$. Its successive powers are
\[
 (1,23,7,16,20,25,24),\qquad \zeta^7=1.
\]
They are distinct, so $\zeta$ is primitive. The polynomial
\[
 f(X)=X^7-t
\]
is irreducible by Eisenstein, is separable because its derivative is
$7X^6$, and has splitting field $K(\alpha)$ with cyclic group $C_7$.
Its roots are $x_i=\alpha\zeta^i$.

The following equality is entirely explicit:
\[
 1+23+7\equiv16+20+24\equiv2\pmod{29}.
\]
Thus
\begin{equation}
 x_0+x_1+x_2=x_3+x_4+x_6.
 \label{rss:eq:collision29}
\end{equation}
The roots in the two triples are all different. This is not a collision
caused by an overlap or by a repeated root of $f$.

\begin{theorem}[Exact repeated-factor resolvent]\label{rss:thm:29}
Over $\F_{29}(t)$,
\begin{align}
 R_{7,3,X^7-t}(Y)
 &= (Y^7-12t)^2(Y^7-t)(Y^7-17t)(Y^7-28t)\notag\\
 &= (Y^7-12t)(Y^{28}-t^4).
 \label{rss:eq:29factor}
\end{align}
Every displayed degree-seven binomial is irreducible. The orbit product
has degree $35$, but its squarefree part has degree $28$.
\end{theorem}
\begin{proof}
There are five translation orbits on the $35$ triples. For a representative
$A$, set
\[
 c_r(A)=e_r(\zeta^i:i\in A),\qquad r=1,2,3.
\]
The representative values and their seventh powers are as follows; the
second and third coordinates will also be used for the repair below.
\begin{center}
\begin{tabular}{@{}crrrr@{}}
\toprule
$A$&$c_1$&$c_2$&$c_3$&$c_1^7$\\
\midrule
012&2&17&16&12\\
013&11&1&20&12\\
014&15&10&25&17\\
015&20&14&24&1\\
024&28&22&24&28\\
\bottomrule
\end{tabular}
\end{center}
Translation by $b$ multiplies the sum $c_1\alpha$ by $\zeta^b$.
The product for one orbit is therefore $Y^7-c_1^7t$. Multiplying the five
orbit products gives the first line of \eqref{rss:eq:29factor}. The four
constants $1,12,17,28$ are the fourth roots of unity in $\F_{29}$, which
gives the second line. Each nonzero binomial $Y^7-ct$ is Eisenstein at $t$,
hence irreducible. The four constants are distinct, so the asserted
squarefree degree follows.
\end{proof}

\subsection{Characteristic two: the Fano-plane degeneration}

Let $K=\F_8(t)$, let $\zeta$ generate $\F_8^\times$, and again let
$\alpha^7=t$. The polynomial $X^7-t$ is separable and irreducible over $K$.
Its seven nonzero root coefficients are all of $\F_8^\times$.

A triple of distinct nonzero field elements has sum zero precisely when it
is $\{a,b,a+b\}$. There are $7\cdot6/6=7$ such triples, the lines of the
Fano plane. Multiplication by $\F_8^\times$ permutes the other triples and
acts transitively on their nonzero sums. The remaining $28$ triples are
therefore distributed equally, four per nonzero sum. Hence
\begin{equation}
 R_{7,3,X^7-t}(Y)=Y^7(Y^7-t)^4\quad\text{over }\F_8(t).
 \label{rss:eq:char2}
\end{equation}
The auxiliary product has only eight distinct roots, even though the
original polynomial has seven distinct roots and a cyclic Galois group.

For septic pair sums in characteristic two, the same field gives
$\zeta^a+\zeta^b=\zeta^c+\zeta^d$ for suitable distinct pairs, in agreement
with $2\in\mathcal B_{7,2}$. For instance, in
$\F_2[z]/(z^3+z+1)$, the pairs of field elements $\{1,z\}$ and
$\{z^2,1+z+z^2\}$ have equal sum.

\subsection{Characteristic seven: an Artin--Schreier degeneration}

Over $\F_7(t)$, let $f=X^7-X-t$ and choose a root $\alpha$.
The irreducibility argument in Section~\ref{rss:sec:criterion} applies, and the
roots are $\alpha+i$ for $i\in\F_7$. A triple sum is $3\alpha+c$.
Translation of a triple by $b$ changes $c$ by $3b$, so each of the seven
values occurs exactly five times. Moreover,
\[
 (3\alpha)^7-3\alpha=3t.
\]
It follows that
\begin{equation}
 R_{7,3,X^7-X-t}(Y)=(Y^7-Y-3t)^5.
 \label{rss:eq:char7}
\end{equation}
Likewise,
\begin{equation}
 R_{5,2,X^5-X-t}(Y)=(Y^5-Y-2t)^2
 \quad\text{over }\F_5(t).
 \label{rss:eq:char5}
\end{equation}
These identities make the exceptional prime-degree characteristic
concrete, rather than merely noting that Fourier inversion cannot divide
by $p$.

\section{The exact implication for published septic factor patterns}
\label{rss:sec:published}

Proposition 7 and the subsequent specialized discussion in
Faggal--Lazard describe a cyclic septic case using five degree-seven
factors; Section 5.2 permits characteristic $29$ under its stated
exclusions \cite{FaggalLazard2014}. Equation~\eqref{rss:eq:29factor} shows that
one cannot universally interpret those factors as five \emph{distinct}
irreducible factors attached to five distinguishable subset orbits.
Two orbit products coincide.

If factors are counted with multiplicity, our example still has five
irreducible factors in that weaker sense. The failure is therefore stated
precisely: the unqualified separation needed to identify different subset
orbits from different scalar factors is false. Also, the degree-$35$
orbit product is not a minimal polynomial of any one triple sum; in this
example each individual triple sum has degree seven.

The sharp repair for this particular separation step is now known:
\[
 \charac K\notin\{2,7,29\}.
\]
This assertion follows for \emph{every} irreducible separable septic from
Theorem~\ref{rss:thm:small}, without a solvability assumption. When retaining
the separate characteristic restrictions of a radical construction, one
must take their union with these restrictions. Conversely, an
input-specific squarefreeness check is sufficient for the orbit-to-factor
principle even in a bad characteristic.

This is not a complete audit of all matching predicates or denominator
conditions used in the published algorithm. Additional root relations can
have different cyclotomic norm obstructions. Excluding $29$ repairs the
triple-sum separation claim; it does not by itself certify every subsequent
formula. The next section offers a repair that avoids relying on a
characteristic-only exclusion at all.

\section{A uniform separating-invariant repair}
\label{rss:sec:repair}

Let $x_1,\ldots,x_n$ be pairwise distinct elements of a field extension of
$K$. Fix $1\leq k<n$, and set $N=\binom nk$. For each $k$-subset $A$ let
\[
 E(A)=(e_1(x_A),\ldots,e_k(x_A)),
\]
where $e_r$ is the $r$th elementary symmetric function of the elements
whose indices lie in $A$.

\begin{lemma}[Subset descriptors]\label{rss:lem:descriptor}
The map $A\mapsto E(A)$ is injective in every characteristic.
\end{lemma}
\begin{proof}
Its coordinates determine the monic polynomial
\[
 \prod_{i\in A}(Z-x_i)
 =Z^k-e_1(x_A)Z^{k-1}+\cdots+(-1)^k e_k(x_A).
\]
Equal coordinate vectors give equal polynomials and therefore equal sets of
roots. Pairwise distinctness of the original $x_i$ then gives equal index
subsets.
\end{proof}

Introduce a parameter $u$ and compress this vector to the scalar
\begin{equation}
 h_A(u)=\sum_{r=1}^k u^{r-1}e_r(x_A).
 \label{rss:eq:descriptor}
\end{equation}
For different subsets $A,B$, the difference $h_A-h_B$ is a nonzero
polynomial in $u$ of degree at most $k-1$.

\begin{theorem}[Finite separating family]\label{rss:thm:repair}
At most
\begin{equation}
 B_{n,k}=(k-1)\binom{N}{2},\qquad N=\binom nk,
 \label{rss:eq:bound}
\end{equation}
values of $u$ in an algebraic closure of the field containing the $x_i$
can cause two of the $h_A(u)$ to
coincide. Consequently, any $B_{n,k}+1$ distinct values in $K$ include one
for which all the $h_A(u)$ are distinct.

If the $x_i$ are the roots of a separable polynomial over $K$, then
\begin{equation}
 H_u(Y)=\prod_{|A|=k}(Y-h_A(u))\in K[Y]
 \label{rss:eq:repaired}
\end{equation}
for every $u\in K$. At a separating parameter, its irreducible factors
recover the exact Galois orbit sizes on $k$-subsets.
\end{theorem}
\begin{proof}
Each of the $\binom N2$ nonzero difference polynomials has at most $k-1$
roots. Taking their union gives \eqref{rss:eq:bound}, with no restriction on
characteristic. The product \eqref{rss:eq:repaired} is invariant under every
permutation of the $x_i$, so its coefficients are symmetric polynomials in
the roots and hence belong to $K$. At a separating parameter,
Lemma~\ref{rss:lem:orbitfactor} applies.
\end{proof}

\begin{writenote}[Existing formal counterparts]
The project already formalizes this strategy for sextic block systems in
\repo{Algebra/PolynomialFormulas/Lean/PolynomialFormulas/SexticSeparatingInvariants.lean}
\cite{repoSextic}, whose docstring starts from the same observation that
scalar pair and triple invariants can collide after specialization.
\code{rootPolynomial_injective} (l.~115) is the argument of
Lemma~\ref{rss:lem:descriptor} for six roots and arbitrary index subsets,
over any commutative domain (no characteristic hypothesis).
\code{exists_nat_eval_injective} (l.~1048) separates a finite injective
family of multivariate polynomials by one natural-number evaluation point,
via the combinatorial Nullstellensatz, over a \emph{characteristic-zero}
domain, and \code{exists_pairDescriptor_separating_evaluation} and
\code{exists_tripleDescriptor_separating_evaluation} (l.~1074, 1083) apply it
to the sextic block descriptors. The resulting ``proved terminating unbounded
search'' is described in \repo{Algebra/PolynomialFormulas/Decidability.md}
(l.~284--290). These are analogues of the declarations
\code{elementarySubsetDescriptor_injective} and
\code{separatingParameter_exists} proposed in
Section~\ref{rss:sec:formalization}. What
Lemma~\ref{rss:lem:descriptor} and Theorem~\ref{rss:thm:repair} add is the
general $(n,k)$ setting, the one-parameter compression $h_A(u)$, and an
explicit hitting bound $B_{n,k}$ valid in every characteristic.
\end{writenote}

\begin{writenote}[Sibling report, added 29 September 2026 (batches 40--41)]
The sibling report \code{sextic-block-resolvent-separators} in this
collection (batch~41, manuscript~03, written without knowledge of this
report) bounds the parameter search built on those sextic descriptors.
It compresses the pair- and triple-partition descriptors of
\repo{SexticSeparatingInvariants.lean} on the curve $(u,v)=(t,2t^2)$ to a
quadratic $h_P(t)$ and a linear $k_T(t)$ (its $h_P$ is not the $h_A(u)$ of
\eqref{rss:eq:descriptor}), counts the collisions by orbits of root
permutations, and proves in writing that one of $1,\ldots,196$ separates
both families for every separable sextic in characteristic zero, sharply
for arbitrary sets of complex test values. It is thus a sextic-partition
counterpart of Theorem~\ref{rss:thm:repair}, obtained by the kind of orbit
count that Research question~14.3 (Section~\ref{rss:sec:questions}) asks
for. It concerns neither septic triples nor positive characteristic, so it
answers no question of this report. It is not formalized: the Lean search
it bounds is still unbounded.
\end{writenote}

For septics and triples,
\[
 N=35,\qquad B_{7,3}=2\binom{35}{2}=1190.
\]
Thus 1191 \emph{distinct} parameters always suffice whenever the base field
has at least that many elements. Over $\Q$ one may use $0,1,\ldots,1190$.
In characteristic $29$, those same integer numerals would repeat and would
not constitute a valid hitting set. Over $\F_{29}(t)$ one can instead use
$1,t,t^2,\ldots,t^{1190}$.

Over a finite field with fewer elements, the theorem does not assert the
existence of a separating parameter in that field. Passing to a sufficiently
large finite extension, or keeping $u$ transcendental, supplies distinct
parameters; the associated change of base field must then be recorded when
interpreting Galois factor patterns.

\subsection{How to compute the repaired product without knowing the roots}

Regard $x_1,\ldots,x_n$ first as indeterminates. The product
\eqref{rss:eq:repaired} is a symmetric polynomial in them, with coefficients
in $\Z[u,Y]$. The fundamental theorem of symmetric polynomials expresses
it as a polynomial in the elementary symmetric functions of the full root
set. Substitution of the coefficients of the original monic polynomial
therefore computes $H_u$ over its coefficient field. This is an algebraic
construction; it does not require numerical roots or a pre-existing
radical solution.

A deterministic specialization procedure is consequently available:
compute $H_u$ for each chosen parameter, test
$\gcd(H_u,H'_u)=1$, and stop at the first success. Theorem~\ref{rss:thm:repair}
proves termination within the stated number of distinct candidates. It is
a correctness and termination bound, not a claim of good asymptotic
complexity. Expanding the universal symmetric polynomial can be expensive.

\subsection{Power-sum variants and their characteristic restriction}

When $1,\ldots,k$ are invertible in $K$, Newton identities show that the
vector
\[
 \left(\sum_{i\in A}x_i,\sum_{i\in A}x_i^2,\ldots,
 \sum_{i\in A}x_i^k\right)
\]
also distinguishes subsets. Then one may use
\[
 \sum_{i\in A}\bigl(x_i+ux_i^2+\cdots+u^{k-1}x_i^k\bigr).
\]
This has the appealing form of an ordinary subset sum after a common
polynomial transformation of the roots. The elementary-symmetric version
\eqref{rss:eq:descriptor}, however, avoids division by $2,\ldots,k$ and works
in every characteristic. These two versions should not be silently
identified in small characteristic.

\subsection{Repairing the characteristic-29 example explicitly}

For $X^7-t$ over $\F_{29}(t)$, use $u=1$. Then
\[
 h_A(1)=c_1(A)\alpha+c_2(A)\alpha^2+c_3(A)\alpha^3.
\]
The $35$ triples have distinct coefficient vectors
$(c_1,c_2,c_3)$ by Lemma~\ref{rss:lem:descriptor}, applied to the seven distinct
powers of $\zeta$. Since $1,\alpha,\ldots,\alpha^6$ are linearly
independent over $\F_{29}(t)$, all $35$ displayed values are distinct.
Thus one nonzero parameter already repairs this particular example.

The five now-distinct irreducible factors have the following compact exact
specification. For each of the five rows $(c_1,c_2,c_3)$ in the table in
Theorem~\ref{rss:thm:29}, take
\begin{equation}
 Q_A(Y)=\Res_Z\bigl(Z^7-t,
           Y-c_1Z-c_2Z^2-c_3Z^3\bigr).
 \label{rss:eq:repairfactors}
\end{equation}
Each $Q_A$ is monic of degree seven, and
$H_1=\prod_A Q_A$. They are distinct and irreducible by
Lemma~\ref{rss:lem:orbitfactor}. This is a coefficient-level description of
the repaired factors, not merely a statement that some repair exists.

For completeness, the affine solvable septic groups give the following
orbit patterns, now valid for any separating equivariant descriptor:
\begin{center}
\begin{tabular}{@{}cl@{}}
\toprule
Group & Factor degrees\\
\midrule
$C_7$ & $7,7,7,7,7$\\
$C_7\rtimes C_2$ & $7,7,7,14$\\
$C_7\rtimes C_3$ & $7,7,21$\\
$C_7\rtimes C_6$ & $14,21$\\
\bottomrule
\end{tabular}
\end{center}
These counts follow directly by acting on the 35 triples with
$i\mapsto ai+b$, where the allowed multipliers form the indicated subgroup
of $\F_7^\times$. The verifier also checks these four finite actions.

\section{One-radical Fourier--Kummer reconstruction}
\label{rss:sec:atlas}

We now address a different failure mode: reconstructing roots from modes
without incompatible radical branches. Assume throughout this section that
$\charac E\neq p$, that $\zeta\in E$ is a primitive $p$th root of unity,
and that $\sigma\in\operatorname{Aut}(E)$ has order $p$ and fixes $\zeta$.
Let $B=E^{\langle\sigma\rangle}$ and suppose that
\[
 \sigma(x_i)=x_{i+1},\qquad i\pmod p.
\]
We use the Fourier convention
\begin{equation}
 u_j=\sum_{i=0}^{p-1}\zeta^{-ij}x_i,
 \qquad
 x_i=\frac1p\sum_{j=0}^{p-1}\zeta^{ij}u_j.
 \label{rss:eq:fourier}
\end{equation}
This differs by an index sign from the positive-exponent forward convention
in the inspected Lean module. Substituting its character by its inverse
aligns the conventions exactly; no mathematical change is involved.

The finite geometric-sum identity gives
\[
 \sigma(u_j)=\zeta^j u_j,
 \qquad u_0\in B.
\]
If the $x_i$ are not all equal, at least one $u_j$ with $j\neq0$ is nonzero.
This is the only existence condition needed for a pivot chart.

For $j\in\{1,\ldots,p-1\}$ define $e_j(k)$ to be the unique integer in
$\{1,\ldots,p-1\}$ satisfying
\[
 j e_j(k)\equiv k\pmod p,\qquad 1\leq k\leq p-1.
\]

\begin{theorem}[Coherent one-radical chart]\label{rss:thm:chart}
Assume $u_j\neq0$. Put
\begin{equation}
 a_j=u_j^p,\qquad
 c_{j,k}=\frac{u_k}{u_j^{e_j(k)}}\quad(1\leq k\leq p-1).
 \label{rss:eq:ratios}
\end{equation}
Then $a_j\in B^\times$, $c_{j,k}\in B$, and $c_{j,j}=1$.
For \emph{any} root $\rho$ of $T^p-a_j$ in an algebraic closure of $E$,
define
\begin{equation}
 y_i=\frac1p\left(u_0+
       \sum_{k=1}^{p-1}\zeta^{ik}c_{j,k}\rho^{e_j(k)}\right).
 \label{rss:eq:reconstruct}
\end{equation}
There exists $r\in\F_p$ such that
\begin{equation}
 y_i=x_{i+rj^{-1}}\quad\text{for every }i\in\F_p.
 \label{rss:eq:branchshift}
\end{equation}
In particular, all branches give the same unordered root set, with no
additional independent $p$th-root choices.
\end{theorem}
\begin{proof}
The covariance formula gives
\[
 \sigma(a_j)=a_j,\qquad
 \sigma(c_{j,k})=\zeta^{k-je_j(k)}c_{j,k}=c_{j,k}.
\]
The nonzero denominator is part of the hypothesis, not a convention about
division by zero. Since $E$ already contains $u_j$ and all $p$th roots of
unity, every root of $T^p-a_j$ is $\rho=\zeta^r u_j$ for some $r$.
Substitution in \eqref{rss:eq:reconstruct} yields
\[
 c_{j,k}\rho^{e_j(k)}
 =\zeta^{r e_j(k)}u_k
 =\zeta^{rj^{-1}k}u_k.
\]
Fourier inversion now gives \eqref{rss:eq:branchshift} simultaneously for all
$i$. Thus coherence is derived, not assumed separately for each mode.
\end{proof}

\begin{corollary}[Cyclic extension generated by a pivot]\label{rss:cor:generator}
Under the hypotheses of Theorem~\ref{rss:thm:chart},
$[B(u_j):B]=p$. If $E/B$ is the cyclic extension of degree $p$, then
$E=B(u_j)$ and $T^p-a_j$ is the minimal polynomial of $u_j$ over $B$.
\end{corollary}
\begin{proof}
The $p$ values $\sigma^r(u_j)=\zeta^{rj}u_j$ are distinct. Thus its
minimal polynomial over $B$ has at least $p$ roots, while $T^p-a_j$ gives
an upper bound of $p$ on its degree. The remaining assertion follows by
the degree formula.
\end{proof}

This is a relative minimality statement: the cyclic degree-$p$ step really
is nontrivial, and one $p$th radical suffices over $B$. It is \emph{not} a
claim that an entire radical expression over the original coefficient field
uses the smallest possible number of radicals. The repository's warnings
about stronger optimality assertions remain relevant.

\begin{writenote}[An abstract counterpart in Rocq]
\repo{Algebra/PolynomialFormulas/Coq/LazardCyclicKummerGenerator.v} proves
\code{cyclic_kummer_generator} (l.~32): if $F/E$ is cyclic Galois of degree
$n$ and $E$ contains a primitive $n$th root of unity, there is a nonzero
$x\in F$ with $x^n\in E$ and $F=E(x)$ \cite{repoKummer}. It is the abstract
existence form of Corollary~\ref{rss:cor:generator}; the corollary adds that
any nonzero Fourier mode $u_j$ is such a generator.
\end{writenote}

\subsection{Transition maps on overlaps}

Suppose $u_j,u_h\neq0$, and set $e=e_j(h)$. Then
\[
 u_h=c_{j,h}u_j^e,\qquad
 a_h=c_{j,h}^{\,p}a_j^e.
\]
A branch $\rho_j^p=a_j$ determines the compatible branch
\begin{equation}
 \rho_h=c_{j,h}\rho_j^e,
 \qquad \rho_h^p=a_h.
 \label{rss:eq:transition}
\end{equation}
It gives the same reconstructed ordered tuple when the chart coordinate
conventions are matched. More explicitly, if
$\rho_j=\zeta^r u_j$, then
$\rho_h=\zeta^{re}u_h$, and the shift in the $h$-chart is
$re\,h^{-1}=rj^{-1}$ modulo $p$.

For the invariant coefficients, let $f=e_h(k)$ and $g=e_j(k)$. The integer
\[
 m=\frac{ef-g}{p}
\]
is well defined. Direct substitution gives the exact transition rule
\begin{equation}
 c_{h,k}=c_{j,k}c_{j,h}^{-f}a_j^{-m}.
 \label{rss:eq:transitioncoeff}
\end{equation}
This exhibits an actual atlas, with compatible local coordinates, rather
than unrelated formulas whose radical branches must be guessed anew.

\subsection{What the chart theorem does not compute}

The input to Theorem~\ref{rss:thm:chart} includes the invariant data
$u_0,a_j,c_{j,k}$ over $B$. Constructing these efficiently from the original
polynomial coefficients is a separate invariant-theoretic task. The theorem
does not hide that task inside a claim of an executable coefficient-only
solver. Its immediate use in ProveIt is as a small correctness theorem
between an invariant-recovery stage and an already formalized Fourier
reconstruction stage.

\section{Solvable prime-degree groups and the coefficient-field tower}
\label{rss:sec:groups}

We record the precise field-theoretic setting in which the cyclic chart
arises from a solvable prime-degree equation. The group-theoretic argument
is included to make the dependence explicit.

\begin{lemma}[Solvable transitive groups of prime degree]\label{rss:lem:affine}
A solvable transitive subgroup $G\leq S_p$ is conjugate to
$C_p\rtimes C_d\leq\AGL_1(\F_p)$ for some $d\mid p-1$.
\end{lemma}
\begin{proof}
A transitive action of prime degree is primitive: every block size divides
$p$. Choose a minimal nontrivial normal subgroup $N$ of $G$. Solvability
implies that $N$ is elementary abelian: its commutator subgroup is
characteristic in $N$ and normal in $G$, so minimality and solvability make
it trivial. The primary subgroups of the resulting finite abelian group
are characteristic, forcing a single prime; multiplication by that prime
has characteristic image and cannot be surjective, so this image is zero.
Its orbits form a block system,
so $N$ is transitive; otherwise every orbit would be a singleton, contrary
to faithfulness and $N\neq1$.

An abelian transitive permutation group is regular. Indeed, an element
fixing one point fixes every point by commuting with a transitive group.
Thus $|N|=p$ and $N\cong C_p$. Numbering the points by $\F_p$ identifies
$N$ with the translations. Its normalizer in $S_p$ is the affine group:
a normalizing permutation is determined by the image of $0$ and the
nonzero multiplier by which it conjugates a generator of translations.
Hence $G$ lies in that normalizer. Its multiplier image is a subgroup of
$\F_p^\times$, therefore cyclic of order $d\mid p-1$. Since all
translations are present, subtracting the translation part of each affine
element gives the asserted semidirect product.
\end{proof}

Now let $f\in K[X]$ be separable irreducible of degree $p$ in characteristic
zero, with solvable Galois group $G=C_p\rtimes C_d$, and let $L$ be its
splitting field. Write
\[
 B_0=L^{C_p},\qquad E=L(\zeta),\qquad B=B_0(\zeta).
\]
Then $B_0/K$ is cyclic of degree $d$, and $B/K$ is an abelian Galois
extension, as the compositum of $B_0/K$ with a cyclotomic extension.
Moreover,
\[
 [B:K]\leq d(p-1)\leq(p-1)^2.
\]
Because $[B:B_0]$ divides $p-1$, it is relatively prime to
$[L:B_0]=p$. Thus $L\cap B=B_0$ and
\[
 [E:B]=p.
\]
The cyclic generator on $L/B_0$ extends to $E$ fixing $B$, so all the
hypotheses of the Fourier--Kummer chart theorem are available after
numbering the roots by this generator.

\begin{corollary}[Field-theoretic solver decomposition]\label{rss:cor:tower}
In characteristic zero, a solvable irreducible prime-degree equation admits
the following decomposition: construct the abelian field $B$ of degree at
most $(p-1)^2$ over $K$, select a nonzero Fourier mode, and perform the
single cyclic $p$th-radical reconstruction of Theorem~\ref{rss:thm:chart}.
For $p=5$ the displayed bound is $16$; for $p=7$ it is $36$.
\end{corollary}

These are field-degree bounds, not optimal radical-count bounds. The
abelian stage itself is solvable by radicals by the classical Galois
construction after adjoining suitable roots of unity. A fast implementation
must exploit the smaller actual Galois group and its invariants rather than
blindly realize the worst-case degree bound.

In positive characteristic, the same cyclic $p$-chart remains valid when
$\charac K\neq p$ and the stated cyclic extension and root-of-unity data
are supplied. The equivalence between solvable groups and ordinary radical
towers needs its own characteristic hypotheses and is not inferred merely
from group solvability here.

\section{Every nonempty Fourier support really occurs}
\label{rss:sec:support}

A formula that works for generic tuples need not survive vanishing Fourier
modes. The following construction shows that no nonempty pattern can be
excluded solely from rational irreducibility or from maximal solvable
Galois group.

\begin{theorem}[Full support realization]\label{rss:thm:support}
Let $p$ be an odd prime, let
$\varnothing\neq S\subseteq\{1,\ldots,p-1\}$, and choose nonzero rational
numbers $c_k$ for $k\in S$. Let $\alpha=2^{1/p}$ and let $\zeta$ be a
primitive $p$th root of unity. Put
\begin{equation}
 P(Z)=\sum_{k\in S}c_k Z^k,\qquad
 x_i=P(\alpha\zeta^i),\quad 0\leq i\leq p-1.
 \label{rss:eq:supportfamily}
\end{equation}
The $x_i$ are the roots of an irreducible rational polynomial of degree $p$.
Its splitting field has Galois group $\AGL_1(\F_p)$. Relative to the
indicated cyclic numbering, its Fourier modes are
\begin{equation}
 u_0=0,\qquad
 u_k=\begin{cases}p c_k\alpha^k,&k\in S,\\0,&k\notin S.
 \end{cases}
 \label{rss:eq:supportmodes}
\end{equation}
Thus all $2^{p-1}-1$ nonempty supports occur.
\end{theorem}
\begin{proof}
The polynomial $Z^p-2$ is Eisenstein at $2$, so
$[\Q(\alpha):\Q]=p$. Linear independence of
$1,\alpha,\ldots,\alpha^{p-1}$ shows that $P(\alpha)\notin\Q$.
Since $p$ is prime, the tower formula forces
\[
 \Q(P(\alpha))=\Q(\alpha).
\]
Therefore $P(\alpha)$ has degree $p$, and its images under the $p$
embeddings of $\Q(\alpha)$ are exactly the displayed distinct $x_i$.
The normal closures of $\Q(P(\alpha))$ and $\Q(\alpha)$ coincide.

For completeness, $[\Q(\zeta):\Q]=p-1$ and $\alpha\notin\Q(\zeta)$:
otherwise the degree $p$ would divide $p-1$. Over $\Q(\zeta)$, adjoining
$\alpha$ gives the splitting field of $Z^p-2$; its Galois group embeds in
$\mu_p$, and nontriviality forces degree $p$. Hence the full normal closure
has degree $p(p-1)$ and its natural affine permutation action is all of
$\AGL_1(\F_p)$.

Finally, substitute \eqref{rss:eq:supportfamily} into \eqref{rss:eq:fourier}.
The geometric sum
$\sum_i\zeta^{i(k-j)}$ equals $p$ when $k=j$ and zero otherwise,
giving \eqref{rss:eq:supportmodes}.
\end{proof}

\begin{corollary}[Sharpness of the coordinate-pivot cover]\label{rss:cor:charts}
For a fixed primitive root of unity and a fixed cyclic numbering, none of
the $p-1$ charts $u_j\neq0$ can be omitted if the chart family is required
to cover every irreducible rational degree-$p$ example with full affine
Galois group.
\end{corollary}
\begin{proof}
Use Theorem~\ref{rss:thm:support} with $S=\{j\}$. Only the $j$th chart is
available.
\end{proof}

The qualifiers matter. This is a minimality statement for this specific
coordinate-pivot cover, not a lower bound for every conceivable nonlinear
cover. Relabelling the cyclic generator permutes the nonzero Fourier
indices; choosing such a relabelling adaptively is another form of chart
selection, not a contradiction.

\begin{writenote}[Nonlinear covers, added 29 September 2026 (batch~52)]
Part~II supplies that lower bound for covers of the universal cyclic
quotient $\Spec B$ of Section~\ref{rss:sec:global}: every cover by regular
Kummer charts, with arbitrary open sets and arbitrary nonlinear character
sections, has at least $p-1$ members (Proposition~\ref{ka:prop:atlaslower}
and Corollary~\ref{ka:cor:repo}). It concerns Zariski covers of that
universal quotient, not the coverage of individual rational examples stated
in the corollary above, so it complements the corollary rather than
replacing it.
\end{writenote}

\subsection{Concrete rational examples}

Taking every $c_k=1$, elimination of $Z$ from
$Z^p-2=0$ and $Y-P(Z)=0$ gives the following monic polynomials:
\begin{align*}
 p=5,\ S=\{1,2\}:&\quad Y^5-10Y^2-10Y-6,\\
 p=5,\ S=\{2,3\}:&\quad Y^5-10Y^3+20Y-12,\\
 p=5,\ S=\{1,2,3,4\}:&\quad Y^5-20Y^3-60Y^2-70Y-30,\\
 p=7,\ S=\{1,2\}:&\quad Y^7-14Y^3-28Y^2-14Y-6,\\
 p=7,\ S=\{2,3\}:&\quad Y^7-14Y^4-28Y^2+28Y-12,\\
 p=7,\ S=\{1,\ldots,6\}:&\quad
 Y^7-42Y^5-210Y^4-490Y^3\\
 &\hspace{1.8em}{}-630Y^2-434Y-126.
\end{align*}
Their irreducibility and Galois groups follow from
Theorem~\ref{rss:thm:support}, not from numerical root matching. The exact
resultants are also regenerated by the verifier.

For example, in the first septic case, $u_1=7\alpha$ and
$u_2=7\alpha^2$, while all other nonconstant modes vanish. In the $j=1$
chart,
\[
 a_1=7^7\cdot2,\qquad c_{1,1}=1,\qquad c_{1,2}=\frac17.
\]
The formula becomes
\[
 y_i=\frac17\left(\zeta^i\rho+
                       \frac{\zeta^{2i}}7\rho^2\right),
 \qquad \rho^7=7^7\cdot2.
\]
Choosing $\rho=7\alpha\zeta^r$ recovers
$x_{i+r}=\alpha\zeta^{i+r}+\alpha^2\zeta^{2(i+r)}$ exactly.

\section{Rigidity of the exceptional septic configurations}
\label{rss:sec:rigidity}

The bad-prime calculation can say more than where counterexamples exist.
In characteristics $29$ and $2$, the small number of surviving cyclotomic
characters forces the entire root configuration into a special form.

\subsection{Characteristic 29: precisely the translated binomial locus}

\begin{theorem}[Complete characteristic-29 collision classification]
\label{rss:thm:29rigidity}
Let $K$ be any field of characteristic $29$, and let $f\in K[X]$ be monic,
separable, and irreducible of degree seven. The following are equivalent:
\begin{enumerate}[label=(\roman*)]
\item Two different triples of roots of $f$ have the same sum.
\item For some $\eta,b\in K$ with $b\neq0$,
\[
 f(X)=(X-\eta)^7-b.
\]
\end{enumerate}
Whenever these conditions hold, the Galois group of $f$ is $C_7$, and
\begin{equation}
 R_{7,3,f}(Y)=
 \bigl((Y-3\eta)^7-12b\bigr)
 \bigl((Y-3\eta)^{28}-b^4\bigr).
 \label{rss:eq:29rigidfactor}
\end{equation}
In particular, every noncyclic irreducible septic over a field of
characteristic $29$ has separating triple sums.
\end{theorem}
\begin{proof}
The primitive seventh root $\zeta=23$ lies in the prime field, hence in
$K$. Choose a Galois seven-cycle $\sigma$ and number the roots by it.
Suppose a collision gives $q(\sigma)x_0=0$, with
$q=q_{A,B}$ as in \eqref{rss:eq:q}. The nonzero-character part of $q$ must
have a zero: otherwise the B\'ezout argument in
Theorem~\ref{rss:thm:bad} would rule out the collision.

Table~\ref{rss:tab:triples} shows that a pair with a resultant divisible by
$29$ is affine-equivalent to $A=012$, $B=346$. For the representative
\[
 q_*(T)=1+T+T^2-T^3-T^4-T^6,
\]
its values at the six nontrivial seventh roots are
\[
 \bigl(q_*(\zeta),q_*(\zeta^2),\ldots,q_*(\zeta^6)\bigr)
   =(0,14,18,25,17,20)\quad\text{in }\F_{29}.
\]
There is exactly one zero. An affine change of the subset labels
multiplies values by a nonzero root of unity and permutes the nontrivial
characters, so every potentially colliding pair has exactly one such zero.
Call its index $j$.

Project the relation $q(\sigma)x_0=0$ to the $r$th character by applying
$\sum_{i=0}^6\zeta^{-ir}\sigma^i$. Since $\sigma$ fixes $\zeta$, the
result is
\[
 q(\zeta^r)u_r=0.
\]
Thus every nonconstant Fourier mode except possibly $u_j$ vanishes.
At least one is nonzero, so $u_j\neq0$. Fourier inversion gives
\[
 x_i=\eta+\beta\zeta^{ij},\qquad
 \eta=\frac{u_0}{7},\quad\beta=\frac{u_j}{7}.
\]
The trace shows $\eta\in K$. The full root product is
$(X-\eta)^7-\beta^7$, and its coefficients show that
$b=\beta^7\in K$. This proves (i)$\Rightarrow$(ii).

Conversely, all roots of a polynomial in (ii) are
$\eta+\beta\zeta^i$. The characteristic-$29$ equality
\eqref{rss:eq:collision29} persists after adding $\eta$ to each root, since
both sides contain three roots. Hence there is a triple-sum collision.
The original polynomial is irreducible, and $K$ contains $\mu_7$, so
$K(\beta)$ is its cyclic degree-seven splitting field. Finally, translating
each root adds $3\eta$ to each triple sum. The five-orbit calculation in
Theorem~\ref{rss:thm:29}, with $t$ replaced by $b$, gives
\eqref{rss:eq:29rigidfactor}.
\end{proof}

The conclusion is stronger than adding $29$ to a blanket list of excluded
characteristics. An exact solver in that characteristic can first test
whether its input becomes a pure binomial after removing the trace. If it
does, the equation has an immediate Kummer solution and the repeated-factor
resolvent is understood completely. If it does not, the ordinary triple-sum
resolvent separates. Cyclicity alone is not sufficient for a collision;
the specific primitive element must give the translated binomial form.

\subsection{Characteristic two: affine Fano configurations}

A polynomial $L(Z)=\sum_i a_iZ^{2^i}$ in characteristic two is
\emph{additive}: $L(v+w)=L(v)+L(w)$. Its roots, when distinct, form a
vector space over $\F_2$. This standard structure is exactly the one forced
by a septic subset-sum collision.

\begin{theorem}[Complete characteristic-two collision classification]
\label{rss:thm:2rigidity}
Let $K$ have characteristic two, and let $f\in K[X]$ be monic, separable,
and irreducible of degree seven. The following are equivalent:
\begin{enumerate}[label=(\roman*)]
\item Two different pairs of roots have equal sums.
\item Two different triples of roots have equal sums.
\item For some $\eta,a,b,d\in K$, with $d\neq0$,
\begin{equation}
 f(X)=F(X-\eta),\qquad F(Z)=Z^7+aZ^3+bZ+d.
 \label{rss:eq:additivequotient}
\end{equation}
\end{enumerate}
In this case the roots are $\eta+(W\setminus\{0\})$, where $W$ is a
three-dimensional $\F_2$-subspace of the splitting field. The center
$\eta$ is the trace of a root, and
\begin{equation}
 R_{7,2,f}(Y)=F(Y)^3,\qquad
 R_{7,3,f}(Y)=(Y-\eta)^7f(Y)^4.
 \label{rss:eq:Fanoresolvents}
\end{equation}
Its Galois group embeds in $\operatorname{GL}_3(\F_2)$.
\end{theorem}
\begin{proof}
Adjoin a primitive seventh root $\zeta$. This extends the base field by
degree one or three, relatively prime to seven, so $f$ remains
irreducible: for a root field $M/K$ of degree seven, the degree of
$M K(\zeta)/K$ is divisible by both $7$ and $[K(\zeta):K]$ and is at most
their product, forcing equality. Its splitting-field group over
$K(\zeta)$ consequently contains a seven-cycle $\sigma$ fixing $\zeta$.

Assume (i) or (ii). The collision yields a nonzero polynomial
$q\in\F_2[T]$ with $\deg q\leq6$, $q(1)=0$, and
$q(\sigma)x_0=0$. In characteristic two,
\[
 \Phi_7(T)=(T^3+T+1)(T^3+T^2+1),
\]
with both cubic factors irreducible. The gcd of $q$ and $\Phi_7$ cannot
be one, by the safe-direction proof of Theorem~\ref{rss:thm:bad}. It cannot
be all of $\Phi_7$ either: degree at most six would give $q=\Phi_7$,
contradicting $q(1)=0$ and $\Phi_7(1)=1$. Therefore it is exactly one of
the two cubic factors.

The primitive roots in either factor have indices
$\{1,2,4\}$ or $\{3,5,6\}$. Reversing the cyclic numbering interchanges
these sets. Character projection as in Theorem~\ref{rss:thm:29rigidity}
therefore shows that, after such a reversal if necessary, the roots have
form
\[
 x_i=\eta+v_1\zeta^i+v_2\zeta^{2i}+v_4\zeta^{4i},
 \qquad \eta=\frac{u_0}{7}=u_0\in K.
\]
On the eight-element subfield generated by $\zeta$, the map
\[
 \psi(z)=v_1z+v_2z^2+v_4z^4
\]
is $\F_2$-linear. It is injective: a nonzero kernel element $h$ would
make $\psi(z)=\psi(z+h)$ for distinct nonzero $z,z+h$, contradicting
distinctness of the seven roots. Hence $W=\psi(\F_8)$ has dimension three,
and the roots are exactly $\eta+(W\setminus\{0\})$.

For a finite $\F_2$-vector space $W$ in a field, its monic root polynomial
$L_W(Z)=\prod_{w\in W}(Z-w)$ is additive. One elementary proof is induction
on the dimension. If $W=V\oplus\F_2 h$, then
\[
 L_W(Z)=L_V(Z)L_V(Z-h)
       =L_V(Z)^2-L_V(h)L_V(Z),
\]
which is additive whenever $L_V$ is. Since $|W|=8$ here,
\[
 L_W(Z)=Z^8+aZ^4+bZ^2+dZ.
\]
But $L_W(Z)=Zf(Z+\eta)$ belongs to $K[Z]$, so its displayed coefficients
lie in $K$. Its distinct roots force $d\neq0$. This proves (iii).

Conversely, under (iii), $ZF(Z)$ is an additive polynomial of degree eight
with derivative $d\neq0$. Its eight roots form a three-dimensional
$\F_2$-space $W$, and the nonzero roots are exactly the roots of $F$.
For each nonzero $w\in W$, exactly three unordered pairs of distinct
nonzero elements sum to $w$. Exactly seven unordered triples sum to zero,
and exactly four sum to each nonzero $w$. These counts follow respectively
by choosing one summand outside $\{0,w\}$, and by the seven Fano lines
and the remaining $28$ triples. After translation by $\eta$, pair sums
are unchanged and triple sums acquire $\eta$. This proves both collisions
and the identities \eqref{rss:eq:Fanoresolvents}.

Finally, every Galois automorphism fixes $\eta$ and permutes $W$ while
preserving addition. Its action on $W$ is faithful because the roots
generate the splitting field. Thus the Galois group embeds in
$\operatorname{GL}_3(\F_2)$.
\end{proof}

\begin{corollary}[A nonsolvable colliding configuration]
There are separable irreducible septics in characteristic two with
colliding pair and triple sums and with Galois group of order $168$,
isomorphic to $\operatorname{GL}_3(\F_2)$.
\end{corollary}
\begin{proof}
Let $z_1,z_2,z_3$ be algebraically independent over $\F_2$, put
$L=\F_2(z_1,z_2,z_3)$, and let $G=\operatorname{GL}_3(\F_2)$ act by
invertible linear substitutions. Set $K=L^G$. The fixed-field theorem
for a finite group of field automorphisms gives
$\Gal(L/K)=G$. The seven nonzero linear forms
$v_1z_1+v_2z_2+v_3z_3$ with $v\in\F_2^3\setminus\{0\}$ are distinct,
form one $G$-orbit, and generate $L$. Their monic root polynomial is
therefore separable irreducible of degree seven over $K$, with splitting
field $L$, and its nonzero root set is an additive three-space minus zero.
Theorem~\ref{rss:thm:2rigidity} applies.

The group has order $(8-1)(8-2)(8-4)=168$. It is nonsolvable: it is
generated by elementary transvections, and for distinct $i,j,k$ each
transvection $I+E_{ij}$ is the commutator of $I+E_{ik}$ and $I+E_{kj}$.
Thus the nontrivial group is perfect and cannot be solvable.
\end{proof}

The two exceptional characteristics therefore have markedly different
geometry. In characteristic $29$, a collision forces a single
nonconstant Fourier character and a cyclic binomial equation. In
characteristic two, a three-character Frobenius orbit survives and produces
an affine Fano configuration, which can support a nonsolvable Galois group.
The exceptional characteristic seven remains a separate unipotent problem,
since its cyclic action cannot be diagonalized by seventh roots of unity.

\section{A global obstruction: the universal cyclic cover is not one Kummer equation}
\label{rss:sec:global}

The local pivot formulas work, but could a more ingenious single regular
radical coordinate work on the entire distinct-root locus? The answer is
negative in a precise algebraic sense. We prove it, rather than infer it
from the zeros of a few familiar formulas.

\subsection{The universal ordered configuration}

Let $k$ be algebraically closed of characteristic zero and contain a
primitive $p$th root of unity. Put
\begin{equation}
 A=k[x_0,\ldots,x_{p-1},\Delta^{-1}],\qquad
 \Delta=\prod_{i<j}(x_i-x_j)^2.
 \label{rss:eq:configuration}
\end{equation}
Thus $\Spec A$ parametrizes ordered tuples of pairwise distinct roots.
Let $C_p=\langle\sigma\rangle$ act by cyclically permuting the variables,
and write $B=A^{C_p}$.

No nonidentity cyclic shift fixes a geometric point: since $p$ is prime,
such a fixed point would have all coordinates equal. The quotient
\[
 \Spec A\longrightarrow\Spec B
\]
is therefore a finite \'etale $C_p$-torsor. The term \emph{torsor} means that
a geometric fiber has a free transitive $C_p$-action; it is the algebraic
version of a covering with that deck group.

A global regular Kummer presentation would be a $C_p$-equivariant
isomorphism
\begin{equation}
 A\cong B[T]/(T^p-a),\qquad a\in B^\times,
 \label{rss:eq:globalKummer}
\end{equation}
where a generator acts by $T\mapsto\zeta^jT$ for some $j\neq0$.
The unit condition expresses the absence of ramification on the already
separable locus. It is not an arbitrary strengthening: a Kummer equation
with a zero radicand would fail to be \'etale at that zero.

The general Kummer framework distinguishes torsors given by a unit from
those encoded by a line bundle whose $p$th tensor power is trivial
\cite{StacksKummer}. We now compute the obstruction in this particular
configuration explicitly.

\subsection{The complete unit group}

For $1\leq a\leq(p-1)/2$ and $j\in\F_p$, define the oriented difference
\[
 \delta_{a,j}=x_j-x_{j+a}.
\]
Every unordered pair occurs exactly once among these factors, up to a
constant sign. Since the polynomial ring is a unique factorization domain,
localizing at $\Delta$ gives
\begin{equation}
 A^\times=
 \left\{c\prod_{a=1}^{(p-1)/2}\prod_{j\in\F_p}
          \delta_{a,j}^{n_{a,j}}:
       c\in k^\times,\ n_{a,j}\in\Z\right\}.
 \label{rss:eq:units}
\end{equation}
The representation is unique once the chosen orientations are fixed.
The cyclic action simply sends $\delta_{a,j}$ to $\delta_{a,j+1}$, without
an additional sign. In multiplicative-module notation,
\begin{equation}
 A^\times\cong k^\times\times
                \Z[C_p]^{(p-1)/2}.
 \label{rss:eq:unitmodule}
\end{equation}

\begin{theorem}[No nontrivial eigenunit]\label{rss:thm:eigenunit}
If $u\in A^\times$ and $\sigma(u)=\zeta^r u$, then $r=0$ in $\F_p$.
Consequently the universal cyclic cover has no presentation
\eqref{rss:eq:globalKummer}.
\end{theorem}
\begin{proof}
Write $u$ as in \eqref{rss:eq:units}. Equality up to a nonzero constant does
not change its exponents along the irreducible difference factors.
Therefore $n_{a,j}=n_{a,j-1}$ for all $a,j$. For each $a$, the exponent is
constant around its full cyclic orbit. The product
$\prod_j\delta_{a,j}$ is itself invariant, and the constant $c$ is fixed.
Thus $\sigma(u)=u$. Since $u\neq0$, this forces $\zeta^r=1$.

In \eqref{rss:eq:globalKummer}, the image of $T$ would be a unit, because
$T^p=a\in B^\times$, and would have a nontrivial eigencharacter. This is
impossible by the first assertion. Even an initially unequivariant
$B$-algebra presentation of this form would yield such an eigenunit:
$\sigma(T)^p=T^p$ forces $\sigma(T)/T\in\mu_p$, and this character
must be nontrivial because $T$ generates the extension.
\end{proof}

This rules out a single everywhere-regular Kummer coordinate on the
cyclic quotient, even though rational Kummer coordinates exist over its
function field. It does not rule out piecewise formulas, formulas on a
dense open subset, or a nontrivial line-bundle formulation. Nor is it a
claim about an unrestricted language of nested algebraic expressions with
arbitrary case distinctions.

\subsection{The Picard group of the cyclic quotient}

Recall that $\Pic B$ is the abelian group of isomorphism classes of
invertible $B$-modules, with tensor product as its operation. Such modules
are the algebraic counterparts of line bundles. Since $A$ is a localization
of a polynomial unique factorization domain, every invertible $A$-module
is free of rank one; in particular, $\Pic A=0$.

\begin{lemma}[Descent calculation]\label{rss:lem:descent}
For the above torsor,
\[
 \Pic B\cong H^1(C_p,A^\times).
\]
\end{lemma}
\begin{proof}
An invertible $B$-module becomes free of rank one over $A$.
Choose a generator after this base change. Transport of the generator
under $g\in C_p$ is multiplication by a unit $a_g\in A^\times$.
Compatibility of the action is exactly the cocycle equation
$a_{gh}=a_g\,g(a_h)$. Changing the chosen generator changes the cocycle
by a coboundary.

Conversely, a unit-valued cocycle gives a semilinear action on the free
rank-one $A$-module. Faithfully flat descent along the finite \'etale torsor
recovers an invertible $B$-module. These constructions are inverse and
respect tensor products. The use of $\Pic A=0$ is precisely what permits
one generator after base change for every class.
\end{proof}

\begin{theorem}[Picard group and the intrinsic chart obstruction]\label{rss:thm:picard}
For odd prime $p$,
\begin{equation}
 \Pic\bigl(k[x_0,\ldots,x_{p-1},\Delta^{-1}]^{C_p}\bigr)
 \cong\Z/p\Z.
 \label{rss:eq:picard}
\end{equation}
The nonzero Fourier-character eigenspaces of $A$ are invertible
$B$-modules representing the nonzero powers of a generator of this group.
\end{theorem}
\begin{proof}
By \eqref{rss:eq:unitmodule} and Lemma~\ref{rss:lem:descent}, it suffices to
compute first cohomology of the factors. The action on $k^\times$ is
trivial, so a cocycle is just a character. Thus
\[
 H^1(C_p,k^\times)=\operatorname{Hom}(C_p,k^\times)=\mu_p\cong\Z/p\Z.
\]
For the additive regular module $M=\Z[C_p]$, a cocycle is determined by
$v\in M$ satisfying
\[
 (1+\sigma+\cdots+\sigma^{p-1})v=0.
\]
In coordinates this says that the sum of the $p$ coefficients of $v$ is
zero. Every such vector is a cyclic difference $(\sigma-1)w$: solve for
the coefficients of $w$ by successive partial sums, whose consistency
around the cycle is exactly the zero-sum condition. Hence
$H^1(C_p,\Z[C_p])=0$. Equation~\eqref{rss:eq:picard} follows.

For each $r\in\F_p$, let
\[
 L_r=\{a\in A:\sigma(a)=\zeta^r a\}.
\]
A $C_p$-torsor becomes the regular permutation representation after a
faithfully flat base change. Its character pieces therefore have rank
one, and multiplication gives isomorphisms
$L_r\otimes_B L_s\cong L_{r+s}$. Descent identifies these modules with
the corresponding character cocycles, up to the harmless choice of which
character is called a generator. Thus $L_r$ has class $r$ times that
generator. For $r\neq0$, the class is nonzero and has order $p$.
\end{proof}

This also explains why $p$th powers can be invariant without their chosen
$p$th roots fitting into one global regular coordinate. The obstruction is
not a failure of Kummer theory; it is exactly the line-bundle alternative
allowed by Kummer theory.

\subsection{The Fourier charts trivialize the obstruction locally}

Define universal modes $u_j$ by \eqref{rss:eq:fourier} in $A$, and let
$a_j=u_j^p\in B$. The basic opens
\[
 D(a_j)\subseteq\Spec B,\qquad j=1,\ldots,p-1,
\]
cover $\Spec B$. Indeed, a point missing all of them lifts to a geometric
point of $\Spec A$ with every nonconstant Fourier mode zero; inversion
would make all its coordinates equal, contradicting $\Delta\neq0$.

On $D(a_j)$, $u_j$ is a unit and is a generator of $L_j$. Every other
character component can be divided by the corresponding power of $u_j$
to become invariant. Consequently
\begin{equation}
 A[a_j^{-1}]\cong
 B[a_j^{-1}][T]/(T^p-a_j),\qquad T\longmapsto u_j.
 \label{rss:eq:localKummer}
\end{equation}
To see that there are no missing elements, decompose an element into its
$p$ character components by averaging with powers of $\zeta$; each
component is an invariant multiple of one of $1,u_j,\ldots,u_j^{p-1}$.
These powers are independent by their distinct characters, and the monic
relation is exactly $u_j^p-a_j$.

The overlap formulas \eqref{rss:eq:transition}--\eqref{rss:eq:transitioncoeff}
are the concrete local gluing data for this nontrivial torsor. Theorem~
\ref{rss:thm:eigenunit} says that no single chart of the same global regular
kind can replace the cover. Corollary~\ref{rss:cor:charts} makes the narrower
coordinate-cover bound sharp even on rational examples.

\section{A practical proof architecture for ProveIt}
\label{rss:sec:formalization}

The article suggests a modular extension rather than a monolithic new
formula. Existing Fourier inversion can remain unchanged. Each added
module should expose the assumptions needed by its consumers and avoid a
statement whose hypotheses already contain the desired factorization.

\begin{writenote}[Status of this section]
This section is a plan. None of it is formalized, and no Lean or Rocq code
was delivered with the manuscript or is shipped with this report. Two facts
about the tree should be weighed before acting on it. First, descriptor
separation already exists for sextics in
\repo{Algebra/PolynomialFormulas/Lean/PolynomialFormulas/SexticSeparatingInvariants.lean}
(note after Theorem~\ref{rss:thm:repair}); the proposed
\code{elementarySubsetDescriptor_injective} and
\code{separatingParameter_exists} would generalize it rather than start from
nothing. Second, the arbitrary-degree Fourier core has a Rocq twin,
\repo{Algebra/PolynomialFormulas/Coq/LazardGeneralFourier.v} (note in
Section~\ref{rss:sec:boundary}), so
the atlas cluster could be built on either prover. The plan proposes Lean
only.
\end{writenote}

\subsection{Separating the four certificates}

A coefficient-level solver should keep four logically different objects:
\begin{enumerate}[label=(\arabic*)]
\item \textbf{Root and Galois data:} a separable irreducible polynomial, its
splitting field, and a justified permutation action on its roots.
\item \textbf{Separation data:} an injective equivariant invariant on the
relevant finite set, or a proved squarefreeness test for its orbit product.
\item \textbf{Invariant recovery data:} exact equalities connecting computed
coefficients to $u_0,a_j,c_{j,k}$ in a specified intermediate field.
\item \textbf{Branch data:} one nonzero pivot and one equation
$\rho^p=a_j$; coherence of all reconstructed modes then follows from
Theorem~\ref{rss:thm:chart}.
\end{enumerate}
The second item is not implied by the first. The fourth need not contain
$p-1$ independent radical choices. These distinctions address two different
ways an otherwise plausible solver can be unsound.

\subsection{Suggested declaration-level milestones}

The following are proposed names and statement boundaries, not purported
existing Lean identifiers.
\begin{center}
\small
\begin{longtable}{@{}L{0.40\textwidth}L{0.54\textwidth}@{}}
\toprule
Proposed declaration & Essential conclusion\\
\midrule\endfirsthead
\toprule Proposed declaration & Essential conclusion\\\midrule\endhead
\code{subsetSum_injective_of_resultant_ne_zero} &
A root $p$-cycle and nonzero cyclotomic resultants imply injectivity of
$k$-subset sums.\\
\code{badCharacteristic_binomial_witness} &
A common primitive cyclotomic root gives an Eisenstein binomial
counterexample over a rational function field.\\
\code{septicTriple_badPrimes} &
Finite determinant/orbit certificate yields exactly $\{2,7,29\}$.\\
\code{elementarySubsetDescriptor_injective} &
Subset elementary symmetric vectors distinguish subsets in every
characteristic.\\
\code{separatingParameter_exists} &
A union-of-root-sets argument proves the finite hitting bound.\\
\code{fourierPivot_ratios_fixed} &
The pivot ratios and pivot power lie in the fixed field.\\
\code{fourierPivot_allBranches_shift} &
Every branch of the single $p$th radical gives a cyclic shift of the
original root tuple.\\
\code{fourierSupport_pureExtension_realization} &
The rational family realizes each prescribed nonempty support.\\
\code{cyclicConfiguration_noEigenunit} &
The explicit unit group has no nontrivial character unit.\\
\bottomrule
\end{longtable}
\end{center}

The first cluster needs finite group actions, polynomial resultants,
splitting fields, and finite sets. The atlas cluster needs only field
algebra, finite sums, a primitive character, and a nonzero pivot. It is
therefore a natural small next target, close to the inspected repository
module. The global Picard calculation is conceptually stronger but needs
substantially more algebraic geometry and descent infrastructure; it should
not block the elementary solver improvements.

\subsection{Eliminating hidden circularity}

A theorem saying that a returned vector is correct \emph{provided} it
satisfies the full root factorization is not an invariant-recovery theorem.
Likewise, testing numerical roots against a target polynomial is not an
exact construction of the intermediate-field coefficients. A useful audit
should trace every exported root theorem backward to coefficient data,
nonzero tests, and explicit field operations.

For the finite classification, a checker can accept integer matrices,
verify their determinants by a fraction-free algorithm, verify affine orbit
coverage, and infer the prime set. It should not trust a table solely
because a computer algebra system printed it. The bundled auditor is an
ordinary executable model of that boundary; a verified Lean or Rocq checker
would be a separate deliverable.

For the branch theorem, no polynomial coefficient expansion is necessary.
The short proof works for every odd prime at once. Keeping this theorem
generic avoids large quintic- or septic-specific ring-normalization goals
and makes exceptional support cases explicit rather than leaving them to a
generic nonvanishing assumption.

\section{Further research questions and topics}
\label{rss:sec:questions}

The questions below are proposed follow-on research problems. Their status
as previously published open problems has not been asserted. Each has a
specific target and a natural first test.

\begin{question}[Bad characteristics in higher prime degrees]
Determine $\mathcal B_{p,k}$ for $p=11,13$ and the middle subset sizes.
Can the affine-orbit computation be replaced by a structural description
of the prime divisors of cyclotomic norms of balanced $\{0,\pm1\}$ vectors?
\end{question}
Theorem~\ref{rss:thm:bad} makes this a finite exact problem with guaranteed
counterexamples for every reported bad prime. A useful first deliverable is
an independently checked table, followed by explanations of unusually
large prime divisors. Complementation, cancellation of common subset
elements, and cyclotomic automorphisms should reduce the search.

\begin{question}[One bad prime, different Galois groups]
For higher $p$ and $\ell\in\mathcal B_{p,k}$, which transitive
Galois groups can occur for colliding examples? Can the characteristic-$2$
and characteristic-$29$ rigidity classifications for septics be extended
to a complete classification in characteristic seven?
\end{question}
Theorems~\ref{rss:thm:29rigidity} and~\ref{rss:thm:2rigidity} already answer two
of the septic cases: characteristic $29$ forces a cyclic binomial
configuration, whereas characteristic two permits a nonsolvable group of
order $168$. For higher primes, the number and arrangement of common
cyclotomic roots should impose analogous Fourier-support constraints.
Characteristic $p$ instead requires a unipotent cyclic-operator analysis.

\begin{question}[Smaller deterministic separating families]
For $(n,k)=(7,3)$, replace the bound 1191 by a substantially smaller uniform
number for a specified family of separating invariants. Can one exploit
the Galois action to count entire collision orbits at once?
\end{question}
The present bound simply adds the degrees of all pair differences, often
counting the same exceptional parameter many times. It is not claimed to
be sharp. A good improvement must still handle finite characteristic and
must distinguish the existence of one good parameter from a fixed short
hitting set over the original field.

\begin{question}[An exact characteristic audit for matching predicates]
For the additional root relations used to match invariant factors in
septic formulas, determine the corresponding balanced or unbalanced
coefficient vectors and their cyclotomic obstructions. Is there one finite
set of excluded characteristics that certifies every matching step?
\end{question}
This is more precise than assuming that the prime divisors of $p(p-1)$
are the only possible problems. The characteristic-$29$ calculation shows
that auxiliary invariants can introduce other primes. Unbalanced
relations require separately treating their constant-character component.

\begin{question}[Minimal nonlinear covers of the cyclic configuration]
\label{ka:q:rss-minimal-covers}
What is the smallest number of regular Kummer charts needed to cover
$\Spec B$ when arbitrary nontrivial character sections are allowed?
\end{question}
We prove that one global chart is impossible and that the standard Fourier
coordinate cover needs all $p-1$ of its members. These facts do not settle
the optimum among nonlinear choices. This question links radical formulas
to the number of generators of the nontrivial invertible modules $L_j$.

\begin{writenote}[Answered in Part~II, 29 September 2026 (batch~52)]
The minimum is $d(p-1)$ charts for $\Conf_p(\A^d)/C_p$, hence $p-1$ for the
question as posed ($d=1$): four for quintics and six for septics
(Corollary~\ref{ka:cor:repo}). It comes from the integral Chow ring
\[
 \CH^*\bigl(\Conf_p(\A^d)/C_p\bigr)=\Z[t]/\bigl(pt,\,t^{d(p-1)}\bigr),
 \qquad t=c_1(\LL)
\]
(Theorem~\ref{ka:thm:main}). A cover by $m$ regular Kummer charts
trivializes the character line on every chart and so forces $t^m=0$
(Lemma~\ref{ka:lem:covernilpotence}), while $t^{d(p-1)-1}\neq0$; the $d(p-1)$
Fourier charts attain the bound (Proposition~\ref{ka:prop:fourieratlas}). The
charts may be arbitrary Zariski opens with different nontrivial characters.
Each nontrivial eigenmodule $L_j$ needs exactly $p-1$ generators as an
invertible $B$-module, which settles the link named in the last sentence
above. Part~II works over any field of characteristic different from $p$
containing $\zeta$, so it covers this section's algebraically closed field of
characteristic zero. Part~II is not formalized.
\end{writenote}

\begin{question}[Other permutation subgroups and their unit characters]
Compute $\Pic(A^H)$ and its radical-coordinate obstructions for other
subgroups $H\leq S_n$ acting freely on ordered distinct-root
configurations, beginning with the affine prime-degree groups.
\end{question}
The descent method still reduces the problem to $H^1(H,A^\times)$.
However, permutations can reverse orientations of difference factors, so
the clean cyclic splitting \eqref{rss:eq:unitmodule} can acquire character
corrections. Those signs must be derived rather than discarded.

\begin{question}[Coefficient-efficient invariant ratios]
For quintics and septics, construct small exact expressions for the
chart data $a_j,c_{j,k}$ in a computationally convenient presentation of
the abelian intermediate field, and compare them with existing large
invariant formulas.
\end{question}
The atlas theorem guarantees that only one $p$th radical is necessary once
these invariants are known. The unresolved algorithmic issue is computing
them without already knowing an ordered root tuple. Modular reconstruction,
trace pairings, and carefully chosen resolvents are plausible approaches.

\begin{question}[Verified finite certificates and branch reconstruction]
Can the 29 norm determinants, the affine orbit partitions, and the generic
one-radical chart theorem be formally verified with a small trusted kernel
boundary and without native computation assumptions?
\end{question}
The provided certificates are deliberately small. A realistic first
milestone is the septic triple classification and the one-radical theorem,
followed by a coefficient-level example and a full audit of its dependency
graph. This is a formalization target, not a claim that such a proof is
already included.

\begin{question}[Numerically stable exact chart selection]
Among nonzero Fourier charts, which selection criterion minimizes
intermediate coefficient growth and numerical conditioning while remaining
compatible with exact certification?
\end{question}
The algebraic theorem only distinguishes zero from nonzero. A practical
hybrid solver could use approximate information to propose a pivot and
exact arithmetic to certify it. The mathematical output must remain
independent of the heuristic and preserve the transition identities.

\section{Conclusion}

Two kinds of exceptional behavior have different causes. Subset-sum
collisions are arithmetic: they are controlled exactly, away from the
prime degree itself, by cyclotomic norm divisibility. The exceptional
characteristic $29$ for septic triple sums is a concrete manifestation,
with a separable irreducible cyclic witness and a complete repeated-factor
resolvent. The collision locus is even rigid: characteristic $29$ forces
a translated pure binomial, while characteristic two forces an affine
Fano configuration and an additive-polynomial quotient. A separating
vector compressed to a scalar supplies a uniform
repair with an explicit finite search bound.

Exceptional radical coordinates are geometric as well as arithmetic.
Every nonempty Fourier support occurs even for rational examples with full
affine Galois group. A nonzero pivot gives a coherent one-radical local
formula, but the universal cyclic cover carries a nontrivial order-$p$
line-bundle class and admits no single global regular Kummer generator.
The right object is therefore an atlas with certified transitions, not a
formula whose denominators are assumed never to vanish.

For ProveIt, the immediate gain is a sharper proof architecture: separate
root separability from resolvent separation, separate invariant recovery
from Fourier inversion, and derive branch coherence from one shared
radical. The results here supply exact mathematical statements and small
certificates for those interfaces, while leaving coefficient optimization
and kernel formalization as explicit next tasks.

\appendix
\section{Reproducibility and the finite proof boundary}
\label{rss:app:repro}

\subsection{Files in the archive}

The archive contains this source and compiled PDF, a README, two checking
programs, and generated certificate files. The principal program is
\code{verify_results.py}; it uses exact SymPy integer and rational
arithmetic, not floating-point roots. The independent program
\code{audit_certificates.py} uses only the Python standard library to check
the orbit and determinant certificates.

Typical commands are
\begin{lstlisting}
python -m pip install sympy
python verify_results.py --out certificates
python audit_certificates.py --cert certificates/cyclotomic_norm_certificates.json
pdflatex -interaction=nonstopmode -halt-on-error article.tex
pdflatex -interaction=nonstopmode -halt-on-error article.tex
pdflatex -interaction=nonstopmode -halt-on-error article.tex
\end{lstlisting}
The source uses common LaTeX packages and has no external figure or
bibliography-file dependency.

\begin{writenote}[Shipped names]
The commands above use the delivery layout. In this report the scripts are
\code{code/verify_results.py} and \code{code/audit_certificates.py}, the
certificate files are in \code{data/} (not \code{certificates/}), and the
delivered PDF is not shipped. Run from the report directory, writing fresh
outputs to a scratch directory so that the shipped records are never
overwritten:
\begin{lstlisting}
uv run --no-project --with sympy==1.14.0 python code/verify_results.py --out <tmpdir>
py code/audit_certificates.py --cert data/cyclotomic_norm_certificates.json
\end{lstlisting}
The auditor is read-only without \code{--out} and prints its report, which
equals \code{data/independent_audit_report.json}. At the intake the fresh
outputs of the verifier equalled the files in \code{data/} after removing
carriage returns; the \code{"python"} field of
\code{verification_report.json} records the interpreter and may differ.
\end{writenote}

\subsection{Recorded checks}

The executed full verification run used Python 3.13.5 and SymPy 1.14.0.
The report \code{certificates/verification_report.json}\footnote{[write]
Shipped in this report as \code{data/verification_report.json}.} records a passing
result for:
\begin{center}
\begin{tabular}{@{}lr@{}}
\toprule
Check & Count\\
\midrule
Exact pairwise cyclotomic resultants & 850\\
Independent orbit-representative integer determinants & 29\\
Finite-field counterexample families & 4\\
Distinct repaired characteristic-$29$ coefficient vectors & 35\\
Nonempty Fourier supports, degrees five and seven & 78\\
Exact Fourier-mode identities & 516\\
Nonzero pivot charts & 224\\
Reconstructed root entries over all branches in those charts & 10,208\\
Exact worked minimal polynomials & 6\\
Solvable affine septic orbit patterns & 4\\
Finite-field character-root patterns for septic rigidity & 1,400\\
\bottomrule
\end{tabular}
\end{center}

The Fourier test algebra is
\[
 \Q[\alpha,\zeta]/(\alpha^p-2,\Phi_p(\zeta)).
\]
Its elements are stored with exact rational coefficients in the basis
$\alpha^a\zeta^b$, $0\leq a<p$, $0\leq b<p-1$.
The field degree calculation in Theorem~\ref{rss:thm:support} justifies this
basis. The tests run through every nonempty support with coefficient one,
every available pivot, all $p$ choices of the radical branch, and all $p$
reconstructed roots. General rational coefficients are covered by the
proof, not exhaustively by this finite test.

The checks do not prove that the cited repository modules compile in a
fresh environment, and no Lean or Rocq process was run for this article.
Nor do these checks verify the Picard calculation: that part is the written
descent proof in Section~\ref{rss:sec:global}. Keeping these boundaries explicit
prevents an exact computation from being mistaken for an unrelated formal
verification claim.

\section{A compact dependency map}

The central results have the following logical dependencies:
\begin{align*}
 &\text{Cauchy theorem + cyclic operator B\'ezout identity}\\
 &\hspace{2em}\Longrightarrow\text{safe resultant criterion},\\[3pt]
 &\text{cyclotomic common root + Eisenstein binomial}\\
 &\hspace{2em}\Longrightarrow\text{necessity at every non-$p$ bad prime},\\[3pt]
 &\text{finite norm tables + exact criterion + Artin--Schreier witness}\\
 &\hspace{2em}\Longrightarrow\text{complete small-degree classification},\\[3pt]
 &\text{septic norm table + characteristic-$29$ character projection}\\
 &\hspace{2em}\Longrightarrow\text{translated-binomial rigidity},\\[3pt]
 &\text{two cubic cyclotomic factors in characteristic two}\\
 &\hspace{2em}\Longrightarrow\text{affine Fano and additive-quotient rigidity},\\[3pt]
 &\text{subset polynomial identity + finite root bound}\\
 &\hspace{2em}\Longrightarrow\text{uniform separating repair},\\[3pt]
 &\text{Fourier covariance + one nonzero mode + inversion}\\
 &\hspace{2em}\Longrightarrow\text{all-branch coherent reconstruction},\\[3pt]
 &\text{prime-degree pure extension + geometric sums}\\
 &\hspace{2em}\Longrightarrow\text{all support patterns},\\[3pt]
 &\text{factorial configuration ring + explicit units + descent}\\
 &\hspace{2em}\Longrightarrow\Pic B\cong\Z/p\Z\text{ and no global eigenunit}.
\end{align*}
There is no dependence of the arithmetic classification on the geometric
Picard theorem, and no dependence of the Fourier chart theorem on the
finite norm table. The three strands can be independently reviewed and
formalized.


%% ====================================================================
%% Batch-52 addition (manuscript 09, "Optimal Kummer Atlases for Prime
%% Cyclic Configurations"). All labels of Part II carry the prefix ka:.
%% Sections continue Part I's arabic numbering (16 onwards) after its
%% appendices A-B; the manuscript's appendix is Appendix C.
%% ====================================================================
\clearpage
\setcounter{section}{15}
\renewcommand{\thesection}{\arabic{section}}
\makeatletter\gdef\theHsection{\arabic{section}}\gdef\Hy@chapapp{section}\makeatother
\phantomsection\part*{Part II.\quad Optimal Kummer atlases for prime cyclic configurations}
\addcontentsline{toc}{part}{Part II.\texorpdfstring{\quad}{ } Optimal Kummer atlases for prime cyclic configurations}

\noindent\emph{Sections~1--15 and Appendices~A--B form
the original report (Part~I), unchanged except for dated \textbf{[write]}
notes pointing to this addition. Part~II was added on 29 September 2026 in
batch~52 of ProveIt's incoming-report intake, from one manuscript written
against Part~I as printed above. It answers Part~I's Research
question~\ref{ka:q:rss-minimal-covers}: covering the cyclic quotient
$\Spec B$ of the distinct-root configuration by regular Kummer charts, with
arbitrary nonlinear character sections allowed, needs exactly $p-1$ charts.
More generally, $\Conf_p(\A^d)/C_p$ has integral Chow ring
$\Z[t]/(pt,t^{d(p-1)})$ and needs exactly $d(p-1)$ charts.
Section~\ref{ka:sec:write} records the provenance, fixes the notation and
states the formal status of the addition; Sections~\ref{ka:sec:introduction}--%
\ref{ka:sec:conclusion} and Appendix~\ref{ka:app:provenance} are the
manuscript.}

\begin{center}
{\small\sffamily\color{accent} PROVEIT RESEARCH COMPANION\par}
\vspace{0.3cm}
{\LARGE\bfseries\color{ink} Optimal Kummer Atlases for\\[2mm]
 Prime Cyclic Configurations\par}
\vspace{0.25cm}
{\large Exact nonlinear chart complexity, torsion Chow classes,\\
 and prime-degree radical reconstruction\par}
\vspace{0.3cm}
{Research prepared for Vladimir Reshetnikov\par}
{29 September 2026\par}
\vspace{0.2cm}
{\small AI-assisted research manuscript; written proofs, not Lean/Rocq formalization.\par}
\end{center}

\begin{quote}\small
\textbf{Abstract (manuscript~09).}
The ProveIt report \emph{Specialization-Safe Radical Solvers} asks whether
arbitrary nonlinear character sections can cover the universal cyclic
root configuration with fewer regular Kummer charts than the standard
$p-1$ Fourier charts. We answer this question: the exact minimum is
$p-1$, even when the covering opens are arbitrary and the nontrivial
character may vary between charts.

More generally, let $p$ be prime, let $d\geq1$, and let $k$ have
characteristic different from $p$ and contain a primitive $p$th root of
unity. For the free cyclic quotient
$X_{p,d}=\Conf_p(\A_k^d)/C_p$, write $\LL$ for a generating character
line bundle and $N=d(p-1)$. We prove
\[
 \CH^*(X_{p,d})=\Z[t]/(pt,t^N),\qquad t=c_1(\LL).
\]
The exact regular Kummer-chart count and the exact number of global
sections generating any nontrivial character line bundle are both $N$.
The proof identifies the contribution of the full diagonal and proves
that deleting all proper partial diagonals introduces no further Chow
relations. It uses integral cycle classes; rational coefficients would
erase the obstruction.

Further consequences include sharp noncompression for nonlinear
character covariants, explicit nonzero torsion cycles represented by
pure Fourier-mode loci, and invariance under allowing arbitrary partial
collisions while excluding the full diagonal. In a different measure of
complexity, the finite algebra of the cyclic cover needs exactly $d$
global algebra generators. Thus in the original one-dimensional problem
the cover is globally monogenic as an ordinary polynomial algebra, even
though its regular Kummer atlas needs $p-1$ charts. Exact finite checks,
a proof-dependency audit, and concrete further research problems accompany
the proofs. The resolution of the named repository question is explicit;
priority over the entire literature is not asserted.
\end{quote}

\noindent\textbf{Main answer in the original setting.}
For odd prime degree $p$, nonlinear regular Kummer charts cannot improve
on the Fourier cover. The optimum is $p-1$: four charts for quintics and
six for septics. This is a statement about a universal cyclic quotient,
not a count of radicals needed for one polynomial.

\section{[write] Provenance, conventions and status of Part~II}
\label{ka:sec:write}

\subsection{One manuscript, one addition}
\label{ka:sec:provenance}

Part~II is batch-52 manuscript~09, \emph{Optimal Kummer Atlases for Prime
Cyclic Configurations: Exact nonlinear chart complexity, torsion Chow
classes, and prime-degree radical reconstruction}, an AI-assisted research
manuscript prepared for Vladimir Reshetnikov and dated 29 September 2026
(22 pages, A4, as delivered). It arrived as the archive
\code{ProveIt_Optimal_Kummer_Atlases.zip} (inner directory
\code{Optimal_Kummer_Atlases/}, main file \code{optimal_kummer_atlases.tex})
in ProveIt commit \code{458ccfb80} and was placed in \code{a701d9098}. It was
written against ProveIt commit
\code{1ee53d57de253d16cdaad79d1f54bbd95d76d682}, at which Part~I's
\code{article.tex} had the Git blob \code{f8303298d0}. That blob was
unchanged at the placement commit, so every statement the manuscript makes
about ``the report'' refers to Part~I as printed above, before the dated
notes of this write. Its secondary repository source,
\code{Algebra/PolynomialFormulas/LazardQuinticFormalization.tex} (blob
\code{5281cbf260}), is also unchanged. The manuscript's checker, build script
and three JSON records are shipped beside Part~I's files with the prefix
\code{02-kummer-atlases-} (Part~I's unprefixed files count as~01); its
source, README and PDF are not shipped and survive in the arrival commit.
The manuscript did not see any other batch-52 delivery.

The manuscript contributes, in the order printed here: the question, the
main theorem and its scope (Section~\ref{ka:sec:introduction}); character
lines and the equivalence between Kummer charts and trivializations
(Section~\ref{ka:sec:torsors}); the Chow ring of the quotient with only the
full diagonal removed (Section~\ref{ka:sec:punctured}); invariance under
deleting the partial diagonals (Section~\ref{ka:sec:deletion}); the exact
chart count (Section~\ref{ka:sec:minimum}); eigenline generators and
nonlinear covariant noncompression (Section~\ref{ka:sec:generators}); local
reconstruction with one shared radical in every dimension~$d$
(Section~\ref{ka:sec:reconstruction}); the exact ordinary
algebra-generator number~$d$ and configuration noncompression
(Section~\ref{ka:sec:algebragenerators}); invariance under allowed partial
collisions (Section~\ref{ka:sec:robustness}); explicit torsion cycles and
small cases (Section~\ref{ka:sec:cycles}); a proof audit, its finite checks
and a formalization decomposition (Section~\ref{ka:sec:audit}); eight
research questions (Section~\ref{ka:sec:questions}); its conclusion
(Section~\ref{ka:sec:conclusion}); and its provenance and literature
comparison (Appendix~\ref{ka:app:provenance}). Every result, proof,
example, remark, table, limitation and question of the manuscript is
printed.

Where the merge had to choose:
\begin{itemize}
\item \textbf{Numbering.} The manuscript's Section~$n$ is Section~$n+16$
here and its Appendix~A is Appendix~C; theorem-like numbers follow the
sections (the manuscript's Theorem~1.1 is Theorem~\ref{ka:thm:main}). Equation
numbers continue Part~I's: the manuscript's equation~($m$) is printed as
equation~($m+30$).
\item \textbf{Front matter.} The manuscript's title block, abstract and
``Main answer'' paragraph open Part~II; its table of contents is replaced by
the report's.
\item \textbf{Overlap with Part~I.} Three results of Part~II restate or
reprove Part~I in dimension one, and all are kept. The local reconstruction
(Section~\ref{ka:sec:reconstruction}) is Part~I's Theorem~\ref{rss:thm:chart}
and branch transition for every~$d$. The group $\CH^1=\Pic\simeq\Z/p\Z$ for
$d=1$ is Part~I's Theorem~\ref{rss:thm:picard}, proved here by a different
route. The Fourier charts and their covering property
(Proposition~\ref{ka:prop:fourieratlas}) extend Part~I's
\eqref{rss:eq:localKummer} to every~$d$. Dated notes at those places say so.
\item \textbf{Bibliography.} Both sources cite the same Stacks Project
section on Kummer theory; it is printed once, with both annotations. The
manuscript's other five entries keep their text under keys prefixed
\code{ka:}. Its entry for ``the report'' is Part~I at the pin.
\item \textbf{Part~I} is unchanged except for the title-page provenance
line, a sentence in Section~\ref{rss:sec:provenance}, an item in
Section~\ref{rss:sec:nonclaims}, a label on Research
question~\ref{ka:q:rss-minimal-covers}, dated notes after
Corollary~\ref{rss:cor:charts} and after that question, a table-of-contents
line for Part~I, and the merged Stacks Project entry.
\end{itemize}

\subsection{Conventions and notation}
\label{ka:sec:notation}

No symbol of the manuscript was renamed and no normalization was changed.
Part~II uses Part~I's Fourier convention: $\sigma(x_{a,i})=x_{a,i+1}$ and
$u_{a,j}=\sum_i\zeta^{-ij}x_{a,i}$ \eqref{ka:eq:fourier}, so for $d=1$ the
mode $u_{1,j}$ is Part~I's $u_j$ \eqref{rss:eq:fourier}, with
$\sigma(u_j)=\zeta^ju_j$ in both. The character line $\LL_j$ is the sheaf of
$\zeta^j$-eigenfunctions, so for $d=1$ its module of global sections is
Part~I's $L_j$ (Theorem~\ref{rss:thm:picard}). A \emph{regular Kummer chart}
(Section~\ref{ka:sec:introduction}) has a unit radicand and a nontrivial
character, as in Part~I's \eqref{rss:eq:globalKummer} and
\eqref{rss:eq:localKummer}, but its open set is arbitrary, not necessarily
principal or affine. Several letters carry more than one meaning inside
Part~II or between the two parts:

\begin{center}
\small
\begin{tabularx}{\textwidth}{@{}L{0.15\textwidth}>{\raggedright\arraybackslash}X>{\raggedright\arraybackslash}X@{}}
\toprule
Symbol & Meaning in Part~II & Not to be confused with\\
\midrule
$A$, $B$, $k$ &
For $d=1$: $A=k[x_0,\ldots,x_{p-1},\Delta^{-1}]$, $B=A^{C_p}$
\eqref{ka:eq:original}; $k$ any field with $\charac k\neq p$ containing
$\zeta$ \eqref{ka:eq:field} &
The same $A$, $B$ in Section~\ref{rss:sec:global}, where $k$ is
algebraically closed of characteristic zero; $A$, $B$ as index subsets in
Sections~\ref{rss:sec:subset}--\ref{rss:sec:repair}.\\
$X$, $X_{p,d}$; $Y$ &
The quotient $U_{p,d}/C_p$ (for $d=1$, $X=\Spec B$); $Y=U_0/C_p$, only the
full diagonal removed &
Part~I's polynomial variable $X$ and resolvent variable $Y$.\\
$t$ &
$c_1(\LL)\in\CH^1(X)$ &
The transcendental $t$ of $X^7-t$ over $\F_{29}(t)$
(Section~\ref{rss:sec:witnesses}).\\
$N$ &
$d(p-1)$ &
Part~I's norm $N_{A,B}$, $N=\binom nk$ and a minimal normal subgroup.\\
$F$ &
The full diagonal $\{z_0=\cdots=z_{p-1}\}$ &
The finite fields $\F_q$.\\
$a$; $a_{a,j}$ &
A coordinate index $1\le a\le d$; the radicand $a_{a,j}=u_{a,j}^p$ of the
chart $D_{a,j}$; a unit radicand $a$ in $T^p-a$; a B\'ezout integer &
Part~I's $a_j=u_j^p$ (for $d=1$, $a_{1,j}=a_j$).\\
$c_{a',k}$; $c$ &
$u_{a',k}/u^{e_j(k)}$ for a fixed pivot $u=u_{b,j}$, first index a
coordinate \eqref{ka:eq:ratios}; $c$ is also a coordinate index and a ratio
in the pivot change, and a Wilson coefficient in Section~\ref{ka:sec:punctured} &
Part~I's $c_{j,k}$, first index the pivot. True: for $d=1$ Part~II's
$c_{1,k}$ is Part~I's $c_{j,k}$. False: it is Part~I's $c_{1,k}$.\\
$\atlas$, $\gen$, $\alggen$ &
Least number of regular Kummer charts; least number of generating global
sections; least number of global algebra generators &
Part~I's $\alpha=2^{1/p}$ (Section~\ref{rss:sec:support},
Appendix~\ref{rss:app:repro}).\\
$q$ &
The quotient map $U\to X$ &
Part~I's $q_{A,B}$ \eqref{rss:eq:q}; field sizes $q$ in both checkers.\\
$E$, $E_M$, $\EE_W$ &
A vector bundle; $\A^M\setminus\{0\}$; the bundle of the representation $W$ &
Part~I's field $E$ and descriptor $E(A)$.\\
$W$, $W_\pi$ &
An open set; the representation \eqref{ka:eq:normalrepresentation};
$V_\pi\setminus F$ &
Part~I's $L_W$ (Section~\ref{rss:sec:rigidity}).\\
$Z$ &
$Z_M$, the cycles $Z_S$, and $W/C_p$ in Theorem~\ref{ka:thm:sandwich} &
$\Z$; Part~I's elimination variable $Z$.\\
$C$, $C_b$ &
The deleted locus $Y\setminus X$ and its filtration &
The group $C_p$.\\
$h$ &
The hyperplane class of $\Pj^{M-1}$; a pivot index in \eqref{ka:eq:transition} &
Part~I's $h$ in \eqref{rss:eq:transition}, the same role.\\
$S$ &
A smooth variety; strata $S_\pi$; an index set in Proposition~\ref{ka:prop:cycles} &
Part~I's Fourier support $S$ (Theorem~\ref{rss:thm:support}); Lazard's $S_k$.\\
$f$, $M$ &
Covariants $f_i$, a morphism $f$, and $f(T)=\prod(T-x_i)$; a line bundle
$M$ and the dimension $M$ of $E_M$ &
Part~I's input polynomial $f$ and matrices $M_q$.\\
$\CH$ &
Integral Chow groups throughout &
Rational Chow groups, which vanish in positive degree here
(Section~\ref{ka:sec:deletion}).\\
\bottomrule
\end{tabularx}
\end{center}

Two tempting false readings are ruled out in the text itself and are
repeated here. First, a nonzero class in $\Pic B$ alone gives a lower bound
of two charts, not $p-1$; the higher powers $t^2,\ldots,t^{p-2}$ are needed.
Second, a cover by $p-1$ charts does \emph{not} mean that $p-1$ independent
$p$th radicals are needed at one input: each chart uses one radical
(Section~\ref{ka:sec:reconstruction}).

\subsection{Formal status and what is not claimed}
\label{ka:sec:nonclaims}

\begin{itemize}
\item \textbf{Not formalized.} No statement of Part~II is formalized in Lean
or Rocq, and no \code{ka:} label has a Lean or Rocq counterpart. The
repository has no Chow-group development. The nearest project declaration
is the Rocq lemma \code{cyclic_kummer_generator}
(\code{Algebra/PolynomialFormulas/Coq/LazardCyclicKummerGenerator.v}, l.~32):
a cyclic Galois extension containing the needed roots of unity has one Kummer
generator. It is a field-level existence statement; Part~II does not rely on
it, and it says nothing about regular charts on the configuration space.
Section~\ref{ka:sec:audit}'s formalization decomposition is a plan of
proposed milestones, not existing declarations, and no Lean or Rocq code was
delivered or is shipped.
\item \textbf{Evidence.} The geometric results rest on written proofs that
use standard integral intersection theory
\cite{ka:Fulton,ka:Totaro,ka:EdidinGraham}. The finite checker verifies
partition counts, Euler coefficients and Fourier identities only; it does not
verify Chow groups, localization, or the atlas lower bound. Its Fourier tests
are exhaustive over $\F_7$ and $\F_{11}$ and seeded samples over $\F_{29}$,
and its branch checks cover 128 inputs per field.
\item \textbf{The source's own non-claims} stand as printed: the Chow-theoretic
machinery is established work, and Totaro's treatment of cyclic products is
close background; no exhaustive priority review was made. Part~II does not
solve the generic quintic over its symmetric coefficient field, does not
compute coefficient-level invariants, does not minimize nested radicals for
one input, and gives no lower bound for programs with case distinctions.
It concerns regular Kummer presentations on Zariski opens of the specified
universal cyclic quotient. Theorem~\ref{ka:thm:sandwich} concerns
intermediate opens containing all distinct configurations, not further
deletions, and the algebra-generator count is not asserted for those larger
opens. The algebraic chart count is not claimed to equal holomorphic or
continuous counts. The research questions of Section~\ref{ka:sec:questions}
are not asserted to be known open problems.
\item \textbf{Review.} AI-assisted and unrefereed. At the write the proofs
were read; the checker was rerun on a copy (Appendix~\ref{ka:app:provenance});
and a separate script recomputed the proper-partition orbit counts from
Stirling numbers, the Wilson residues $((p-1)!)^d\equiv(-1)^d$, the integral
identity behind \eqref{ka:eq:regularbundle} for $p\le13$, the quintic
coefficient $2\cdot3\cdot4\equiv-1\pmod5$, the tuple and branch counts of
the tables in Section~\ref{ka:sec:audit}, and the pivot-change shift
$reh^{-1}\equiv rj^{-1}$. This is not an independent proof review.
\end{itemize}

\section{The question, the answer, and the scope}
\label{ka:sec:introduction}

\subsection{The repository problem}
The starting point is Research Question~14.5, entitled
\emph{Minimal nonlinear covers of the cyclic configuration}, in the
ProveIt research report \emph{Specialization-Safe Radical Solvers}
\cite{ka:ProveItReport}. Its setting is
\begin{equation}
 A=k[x_0,\ldots,x_{p-1},\Delta^{-1}],\qquad
 \Delta=\prod_{i<j}(x_i-x_j)^2,\qquad B=A^{C_p},
 \label{ka:eq:original}
\end{equation}
where $p$ is odd prime and a generator $\sigma$ cyclically permutes the
variables. The report works over an algebraically closed field of
characteristic zero for this geometric question. It establishes a
nontrivial character line bundle, rules out a single global regular
Kummer coordinate, and shows that none of the $p-1$ standard Fourier
coordinate charts can simply be omitted. It leaves open the optimum
when nonlinear character sections are allowed.

The distinction is genuine. Showing that every member of one particular
cover is necessary does not show that a different cover cannot be
smaller. Nor does a nonzero class in $\Pic(B)$ alone give a lower bound
larger than two. We determine the higher self-intersections of that class
and use them to rule out \emph{all} smaller covers.

The report is associated with the repository's polynomial-formula
project and its audit of solvable-quintic reconstruction
\cite{ka:ProveItLazard}. Our argument does not assume the correctness of any
unformalized claim there: the Fourier construction, the line-bundle
interpretation, and the needed obstruction are proved below.

\begin{writenote}[The report cited here is Part~I]
The report \cite{ka:ProveItReport} is Part~I of this report, cited at the
pinned blob, which is the text of Part~I printed above apart from the dated
\textbf{[write]} notes of batch~52. Its Research Question~14.5 is Research
question~\ref{ka:q:rss-minimal-covers}; the Sections~8, 10 and~12 named in
the bibliography entry are Sections~\ref{rss:sec:atlas},
\ref{rss:sec:support} and~\ref{rss:sec:global}. The nontrivial character
line bundle is Theorem~\ref{rss:thm:picard}, the impossibility of one global
regular Kummer coordinate is Theorem~\ref{rss:thm:eigenunit}, and the
irredundancy of the $p-1$ Fourier coordinate charts is
Corollary~\ref{rss:cor:charts}.
\end{writenote}

\subsection{Notation and the main theorem}
Fix throughout a prime $p$, an integer $d\geq1$, and a field $k$ satisfying
\begin{equation}
 \charac(k)\neq p,\qquad \zeta\in k,\qquad \ord(\zeta)=p.
 \label{ka:eq:field}
\end{equation}
Put
\begin{align}
 V_{p,d}&=(\A_k^d)^p,\notag\\
 U_{p,d}&=\Conf_p(\A_k^d)
   =\{(z_0,\ldots,z_{p-1}):z_i\neq z_j\text{ for }i\neq j\},\notag\\
 X_{p,d}&=U_{p,d}/C_p,\qquad N=d(p-1).
 \label{ka:eq:spaces}
\end{align}
All inequalities in the configuration definition are scheme-theoretic
open conditions, not restrictions to $k$-rational points. The action is
cyclic permutation of the labels. The quotient map is denoted $q$.
Let $\LL_j$ be its weight-$j$ character line bundle, so
$\LL_j\otimes\LL_h\simeq\LL_{j+h}$ and $\LL_0=\OO_X$.
We write $\LL=\LL_1$ and $t=c_1(\LL)$.

A \emph{regular Kummer chart} is an open subset $W\subset X$ on which
$q_*\OO_U$ has an equivariant presentation
\[
 (q_*\OO_U)|_W\simeq\OO_W[T]/(T^p-a),\qquad
 a\in\Gamma(W,\OO_W^\times),\quad \sigma(T)=\zeta^jT,
\]
for some $j\in\{1,\ldots,p-1\}$. Opens need not be principal or affine.
The number of charts in a smallest such finite cover is $\atlas(X)$.
For a globally generated line bundle $M$, let $\gen(M)$ be the least
number of global sections generating it at every point.

\Needspace{14\baselineskip}
\begin{theorem}[Exact ring and exact nonlinear chart complexity]
\label{ka:thm:main}
Under \eqref{ka:eq:field}, there is an isomorphism of graded integral rings
\begin{equation}
 \CH^*(X_{p,d})\simeq\Z[t]/(pt,t^{d(p-1)}),\qquad \deg t=1,
 \label{ka:eq:mainring}
\end{equation}
carrying $t$ to $c_1(\LL)$. Moreover,
\begin{equation}
 \atlas(X_{p,d})=\gen(\LL_j)=d(p-1)
 \qquad (1\leq j\leq p-1).
 \label{ka:eq:maincounts}
\end{equation}
The chart lower bound permits arbitrary nonlinear regular character
sections, arbitrary covering opens, and different nontrivial characters
on different opens.
\end{theorem}

The ring calculation is proved in Sections~\ref{ka:sec:punctured} and
\ref{ka:sec:deletion}; the chart and generator assertions are proved in
Sections~\ref{ka:sec:minimum} and \ref{ka:sec:generators}.
When $N>1$, the class $t^r$ has exact order $p$ for every
$1\leq r<N$. When $N=1$, necessarily $(p,d)=(2,1)$, the Chow ring is
$\Z$ and a single chart works. Including this boundary case prevents an
incorrect blanket assertion that every character line bundle here is
nontrivial.

\begin{corollary}[Answer to the repository question]
\label{ka:cor:repo}
In \eqref{ka:eq:original}, the minimum number of regular Kummer charts
covering $\Spec B$ is $p-1$. Each nontrivial character eigenspace of $A$
requires exactly $p-1$ generators as an invertible $B$-module.
\end{corollary}
\begin{proof}
For $d=1$, the ordered configuration is affine with coordinate ring $A$,
and the finite quotient is $\Spec B$. Apply Theorem~\ref{ka:thm:main} and the
equivalence between vector bundles on an affine scheme and finite
projective modules.
\end{proof}

\subsection{What is and is not claimed}
The mathematical contribution is a complete written answer to the named
repository question, together with the higher-dimensional ring,
noncompression, deletion-invariance, and generator statements proved
here. The intersection-theoretic machinery is established work of
Fulton, Totaro, and Edidin--Graham; our use of it is not a new foundation
for Chow theory \cite{ka:Totaro,ka:EdidinGraham}. Totaro's treatment of cyclic
products is particularly close background. A literature search and
comparison with those sources do not constitute an exhaustive priority
review, and no claim of first discovery of every intermediate formula is
made.

The article does not solve a generic quintic over its symmetric
coefficient field. It does not compute coefficient-level invariants,
minimize the number of nested radicals for an individual input, or give
a lower bound for arbitrary programs with case distinctions. It treats
\emph{regular Kummer presentations on Zariski opens of a specified
universal cyclic quotient}. All main assertions have written proofs;
none is asserted to have been checked by Lean or Rocq. The accompanying
Python checks test finite combinatorics and Fourier identities only.

\section{Cyclic torsors, eigenlines, and regular Kummer charts}
\label{ka:sec:torsors}

\subsection{The free locus and its quotient}
Let $F\subset V=V_{p,d}$ be the full diagonal:
\[
 F=\{z_0=z_1=\cdots=z_{p-1}\}\simeq\A_k^d.
\]
If a nonidentity element of $C_p$ fixes a geometric point of $V$, it
generates the entire group, and all the $p$ coordinates of that point
are equal. Thus the action is free on
\[
 U_0=V\setminus F,
\]
and in particular on $U=U_{p,d}$. The quotients
$Y=U_0/C_p$ and $X=U/C_p$ are smooth quasiprojective $k$-schemes. The
quotient maps are finite \'etale $C_p$-torsors of degree $p$. These
quotients can be constructed as opens of the finite affine quotient
$V/C_p$; freeness makes the restrictions torsors. Smoothness descends
along the \'etale quotient maps.

Here a torsor means that every geometric fiber has a free transitive
$C_p$-action. It does \emph{not} mean that this finite cover is trivial
in the Zariski topology.

\subsection{The character algebra}
Because $p$ is invertible and $k$ contains $\zeta$, the projectors
\begin{equation}
 e_j=\frac1p\sum_{r=0}^{p-1}\zeta^{-jr}\sigma^r
 \qquad (j\in\F_p)
 \label{ka:eq:projectors}
\end{equation}
act on $q_*\OO_U$ and yield
\begin{equation}
 q_*\OO_U=\bigoplus_{j\in\F_p}\LL_j,
 \qquad \LL_j=\{f:\sigma(f)=\zeta^j f\}.
 \label{ka:eq:eigenalgebra}
\end{equation}
After a faithfully flat base change trivializing the torsor, this is the
character decomposition of the algebra of functions on $C_p$. Each
summand has rank one, and multiplication gives isomorphisms
\begin{equation}
 \LL_j\otimes\LL_h\xrightarrow{\sim}\LL_{j+h}.
 \label{ka:eq:lineproduct}
\end{equation}
These properties descend. In particular, $\LL^{\otimes p}\simeq\OO_X$,
$\LL_j\simeq\LL^{\otimes j}$, and $p\,t=0$.
The same constructions apply on $Y$ and restrict to those on $X$.

\begin{lemma}[A Kummer chart is a trivialization of a character line]
\label{ka:lem:kummer}
For a nonempty open $W\subset X$, the following are equivalent:
\begin{enumerate}[label=(\roman*)]
\item $W$ is a regular Kummer chart.
\item $\LL_j|_W$ is trivial for some $j\neq0$ in $\F_p$.
\item $\LL|_W$ is trivial.
\end{enumerate}
If these conditions hold, every nontrivial character line is trivial on
$W$.
\end{lemma}
\begin{proof}
In a Kummer presentation, $T$ generates the relevant eigenline, proving
(i)$\Rightarrow$(ii). If $j$ is nonzero modulo $p$, choose integers $a,b$
with $aj+bp=1$. Using tensor powers and duals,
\[
 \LL\simeq\LL_j^{\otimes a}\otimes
                    (\LL^{\otimes p})^{\otimes b}.
\]
Thus (ii) implies (iii), which also trivializes every $\LL_j$.

Conversely, suppose $s$ is a generator of $\LL_j|_W$, where $j\neq0$.
The multiplication isomorphisms show that $s^p=a$ is a unit of
$\OO_W$. The powers $1,s,\ldots,s^{p-1}$ generate all $p$ distinct
character pieces. They therefore give the asserted isomorphism
$\OO_W[T]/(T^p-a)\simeq(q_*\OO_U)|_W$.
\end{proof}

This is the concrete form of the standard Kummer description by an
invertible sheaf and a trivialization of its $p$th tensor power
\cite{StacksKummer}. Notice that a trivial eigenline does not require
$a$ to have a $p$th root on $W$. The finite torsor can remain nontrivial
on a perfectly valid Kummer chart.

\begin{remark}[Why the radicand must be a unit]
If $a$ vanishes at a point, the fiber $T^p-a$ has a multiple root over an
algebraic closure. It cannot present an \'etale cover there. Unit
radicands are consequently forced by the intended geometry rather than
added as an artificial restriction.
\end{remark}

\begin{remark}[Equivariance does not hide a simpler one-equation chart]
All nonempty opens of $X$ are integral. Suppose the same cover has a
$W$-algebra presentation $\OO_W[T]/(T^p-a)$ with $a$ a unit, without a
specified action on $T$. On the integral inverse image of $W$,
$\sigma(T)/T$ is a $p$th root of unity, hence one of the constant
$\zeta^j$. It cannot be $1$, since then $T$ would be invariant and could
not generate a degree-$p$ algebra. Thus such a presentation automatically
has a nontrivial character, and is included in our chart definition.
\end{remark}

\section{The punctured cyclic quotient and its Chow ring}
\label{ka:sec:punctured}

\subsection{The intersection-theoretic input}
We use integral Chow groups: $\CH_i(S)$ consists of $i$-dimensional
algebraic cycles modulo rational equivalence. For smooth pure-dimensional
$S$, put $\CH^r(S)=\CH_{\dim S-r}(S)$, with its intersection product.
The needed properties are the following standard ones
\cite{ka:Fulton,ka:Totaro,ka:EdidinGraham}: proper pushforward; pullback between smooth varieties; restriction to
opens; the localization sequence
\begin{equation}
 \CH_i(Z)\longrightarrow\CH_i(S)\longrightarrow
 \CH_i(S\setminus Z)\longrightarrow0;
 \label{ka:eq:localization}
\end{equation}
homotopy invariance for vector bundles; Chern classes and their projection
formula; and the projective-bundle calculation
$\CH^*(\Pj^{M-1})=\Z[h]/(h^M)$.

A zero section of a rank-$r$ vector bundle contributes its top Chern
class under the vector-bundle identification of Chow groups. Thus, if
$E\to S$ is a rank-$r$ vector bundle over a smooth base,
\begin{equation}
 \CH^*(\Tot(E)\setminus S)
   \simeq \CH^*(S)/(c_r(E)).
 \label{ka:eq:eulerquotient}
\end{equation}
Indeed, localization and homotopy invariance identify the image of the
zero-section pushforward with multiplication by $c_r(E)$. All these
operations are integral. In particular, no division by the degree of a
finite quotient is permitted.

The following elementary consequence will be central.
\begin{lemma}[Chow groups of an open affine space]
\label{ka:lem:openaffine}
If $W$ is a nonempty open subset of $\A_k^m$, then
$\CH_m(W)=\Z[W]$ and $\CH_i(W)=0$ for $i<m$.
\end{lemma}
\begin{proof}
Homotopy invariance gives the assertion for affine space. Restriction to
$W$ is surjective by localization, so all lower groups vanish. In top
dimension there is one irreducible component, and there are no
$(m+1)$-dimensional subvarieties on which a rational equivalence between
top-dimensional cycles could be supported.
\end{proof}

\subsection{A finite approximation to the character class}
We give an ordinary Chow-group proof of the first ring calculation;
one may equivalently phrase it in equivariant Chow theory. This makes
the integral coefficient issue and the stabilization range explicit.

Let $E_M=\A_k^M\setminus\{0\}$, with the scalar action
$v\mapsto\zeta v$, and let $Q_M=E_M/C_p$. The quotient $Q_M$ is the
complement of the zero section in $\OO_{\Pj^{M-1}}(-p)$. To see this,
regard $E_M$ as the punctured tautological line bundle
$\OO_{\Pj^{M-1}}(-1)^\times$ and take the $p$th tensor power of a
nonzero vector. The fibers of this map are precisely the $C_p$-orbits.
Consequently, \eqref{ka:eq:eulerquotient} gives
\begin{equation}
 \CH^*(Q_M)=\Z[h]/(h^M,ph).
 \label{ka:eq:approxclassifying}
\end{equation}
The character line of $E_M\to Q_M$ is the pullback of the tautological
line or its dual, depending on the convention for the generator. Its
first Chern class $t_M$ is therefore $h$ or $-h$. In particular,
\begin{equation}
 \CH^*(Q_M)=\Z[t_M]/(pt_M,t_M^M).
 \label{ka:eq:approxclassifying2}
\end{equation}
The choice of sign never changes any ring or nonvanishing assertion
below. On each space we consistently define the class by its character
line.

\subsection{Only the full diagonal removed}
Since $p$ is invertible, taking the average of the $p$ points splits the
permutation representation:
\begin{equation}
 V=F\oplus W,\qquad
 W\simeq\bigoplus_{j=1}^{p-1}(\chi^j)^{\oplus d},
 \qquad \dim W=N.
 \label{ka:eq:normalrepresentation}
\end{equation}
Here $F$ is the trivial representation of dimension $d$, and $\chi$
is a generating character. Thus
$U_0\simeq\A^d\times(W\setminus\{0\})$.

\begin{proposition}[Chow ring before deleting partial diagonals]
\label{ka:prop:punctured}
For $Y=(V\setminus F)/C_p$, with its character class $t_Y$,
\begin{equation}
 \CH^*(Y)=\Z[t_Y]/(pt_Y,t_Y^N).
 \label{ka:eq:puncturedring}
\end{equation}
\end{proposition}
\begin{proof}
Form the mixed quotient
\[
 Z_M=(U_0\times E_M)/C_p.
\]
There are two ways to view it. First, it is the complement of the zero
section in the rank-$M$ vector bundle
$(U_0\times\A^M)/C_p\to Y$. The zero section has codimension $M$,
so homotopy invariance and localization give
\begin{equation}
 \CH^r(Y)\xrightarrow{\sim}\CH^r(Z_M)
 \qquad\text{whenever }r<M.
 \label{ka:eq:stablerange}
\end{equation}
The character line on $Y$ pulls back to the character line of the
common torsor $U_0\times E_M\to Z_M$.

Second, project $Z_M$ to $Q_M$. The decomposition
\eqref{ka:eq:normalrepresentation} identifies $Z_M$ with a trivial
$\A^d$-factor times the complement of the zero section of the associated
rank-$N$ vector bundle $\EE_W\to Q_M$. Its line summands have first
Chern classes $j t_M$, with $1\leq j\leq p-1$, each repeated $d$ times.
A dual convention merely permutes this multiset. Hence
\begin{equation}
 c_N(\EE_W)=\left((p-1)!\right)^d t_M^N.
 \label{ka:eq:euler}
\end{equation}
By \eqref{ka:eq:eulerquotient} and \eqref{ka:eq:approxclassifying2},
\begin{equation}
 \CH^*(Z_M)
 =\Z[t_M]/\left(pt_M,t_M^M,((p-1)!)^d t_M^N\right).
 \label{ka:eq:mixedring}
\end{equation}
The coefficient $((p-1)!)^d$ is relatively prime to $p$. If $c$ denotes
this coefficient, choose integers $a,b$ with $ac+bp=1$. Since $N\geq1$,
$pt_M^N=0$, and therefore $t_M^N=a c t_M^N$. Thus the ideal generated
by the Euler class is exactly $(t_M^N)$. Equivalently, Wilson's theorem
gives $c\equiv(-1)^d\pmod p$.

Choose $M>\max\{N,r\}$. Then \eqref{ka:eq:mixedring}, together with
\eqref{ka:eq:stablerange}, computes $\CH^r(Y)$ with generator $t_Y^r$.
The identifications preserve products and character classes. Since $M$
can be chosen arbitrarily large, they determine the entire graded ring
and give \eqref{ka:eq:puncturedring}.
\end{proof}

The stable-range argument is the standard finite-dimensional
classifying-space method of Totaro, also underlying the equivariant
construction of Edidin--Graham \cite{ka:Totaro,ka:EdidinGraham}. What remains
is not a formal consequence of the codimension of the full diagonal:
the configuration space removes \emph{many additional} diagonals, some
of low codimension. We must prove that they do not kill lower powers
of $t_Y$.

\section{Deleting partial diagonals creates no new Chow relations}
\label{ka:sec:deletion}

\subsection{Prime cyclic symmetry of the partition lattice}
A partition $\pi$ of the label set $\F_p$ specifies the linear subspace
\begin{equation}
 V_\pi=\{z_i=z_j\text{ whenever }i,j\text{ lie in one block of }\pi\}
       \simeq\A_k^{d b(\pi)},
 \label{ka:eq:partitionflat}
\end{equation}
where $b(\pi)$ is the number of blocks. Let $S_\pi$ be the stratum on
which $\pi$ is the \emph{exact} equality partition. Then
\[
 S_\pi\simeq\Conf_{b(\pi)}(\A_k^d),
\]
a nonempty open of $V_\pi$. The indiscrete partition corresponds to $F$;
the discrete partition corresponds to $U$.

\begin{lemma}[There are only two invariant partitions]
\label{ka:lem:partitions}
The only partitions of $\F_p$ fixed by translation by $1$ are the
discrete and indiscrete partitions. Every other partition has trivial
stabilizer under $C_p$ and lies in an orbit of size $p$.
\end{lemma}
\begin{proof}
A translation-invariant partition defines a translation-invariant
equivalence relation. Let $H$ be the equivalence class of $0$.
If $a,b\in H$, then $a\sim b$, and translation by $-b$ gives
$a-b\sim0$. Thus $H$ is an additive subgroup of $\F_p$. It is either
$\{0\}$ or all of $\F_p$, and the other blocks are its cosets. This
gives the two stated partitions. Finally, in a group of prime order,
a nontrivial stabilizer would be the entire group.
\end{proof}

For a proper, nontrivial partition, the strata in its orbit are disjoint
and form an induced $C_p$-scheme:
\begin{equation}
 \coprod_{g\in C_p}S_{g\pi}\simeq C_p\times S_\pi.
 \label{ka:eq:inducedstratum}
\end{equation}
Its quotient is therefore isomorphic to $S_\pi$ itself. In particular,
its ordinary Chow groups vanish below its top dimension, by
Lemma~\ref{ka:lem:openaffine}. This is where primality performs the decisive
work.

\subsection{The vanishing pushforward}
Let $q_0:U_0\to Y$ be the free quotient map. For a proper nontrivial
partition, set
\[
 W_\pi=V_\pi\setminus F\subset U_0.
\]
This is a closed subvariety of $U_0$. It is also a nonempty open of the
affine space $V_\pi$, and has codimension $d(p-b(\pi))>0$ in $U_0$.

\begin{lemma}[Transfer-zero for a partial diagonal]
\label{ka:lem:transferzero}
The proper map $W_\pi\to Y$ has zero pushforward of its fundamental
class:
\begin{equation}
 (q_0|_{W_\pi})_*[W_\pi]=0\quad\text{in }\CH_*(Y).
 \label{ka:eq:transferzero}
\end{equation}
\end{lemma}
\begin{proof}
Factor the map as the closed immersion
$i_\pi:W_\pi\hookrightarrow U_0$ followed by the finite map $q_0$.
The scheme $U_0$ is a nonempty open of the affine space $V$, so its
Chow groups below dimension $dp$ vanish. The class
$i_{\pi*}[W_\pi]$ lies in the strictly smaller dimension $db(\pi)$ and
is therefore zero. Proper pushforward by $q_0$ preserves this zero.
\end{proof}

No division by $p$ appears. The argument proves zero by factoring a
proper pushforward through a zero Chow group, not by multiplying a class
by $p$ and then cancelling $p$. That cancellation would be invalid for
the torsion classes we are trying to preserve.

\subsection{From the strata to the whole deleted locus}
Let
\[
 C=Y\setminus X=(V\setminus(U\cup F))/C_p.
\]
For $1\leq b\leq p-1$, the locus in $V$ consisting of tuples with at
most $b$ distinct points is a finite union of the flats $V_\pi$, hence
closed. Removing $F$ and taking the finite quotient provides a closed
filtration
\[
 \varnothing=C_1\subset C_2\subset\cdots\subset C_{p-1}=C.
\]
The difference $C_b\setminus C_{b-1}$ is the disjoint union, over cyclic
orbits of partitions with $b$ blocks, of the quotient strata in
\eqref{ka:eq:inducedstratum}. Each is isomorphic to $S_\pi$.

\begin{lemma}[Generators for the boundary Chow groups]
\label{ka:lem:boundarygenerators}
The groups $\CH_*(C)$ are generated, as abelian groups, by the images
of $[W_\pi]$ under $W_\pi\to C$, where $\pi$ runs over one representative
of every cyclic orbit with $2\leq b(\pi)\leq p-1$.
\end{lemma}
\begin{proof}
Apply localization successively to the closed filtration above. On a
stratum $S_\pi$, the only nonzero Chow group is generated by its
fundamental class. The map $W_\pi\to C$ is proper: it is the
restriction of the finite quotient map to a closed subvariety of $U_0$.
On the dense stratum $S_\pi$ its map to the quotient stratum is an
isomorphism, because the stabilizer of $\pi$ is trivial. Its fundamental
class thus lifts the generator of that quotient stratum with coefficient
one. Subtract these lifts from an arbitrary class. Localization leaves
a class on the preceding closed piece. Induction finishes the proof.

There can be relations between these generators; no independence claim
is needed. The important point is that every generator has a specified
proper lift whose pushforward can be evaluated.
\end{proof}

\begin{theorem}[Partial-diagonal deletion invariance]
\label{ka:thm:deletion}
Restriction induces an isomorphism of integral Chow rings
\begin{equation}
 \CH^*(Y)\xrightarrow{\sim}\CH^*(X_{p,d}).
 \label{ka:eq:deletioniso}
\end{equation}
In particular, \eqref{ka:eq:mainring} holds.
\end{theorem}
\begin{proof}
Localization gives the exact sequence
\[
 \CH_*(C)\xrightarrow{\iota_*}\CH_*(Y)
      \longrightarrow\CH_*(X)\longrightarrow0.
\]
By Lemma~\ref{ka:lem:boundarygenerators}, it suffices to evaluate $\iota_*$
on the proper lifts of the fundamental classes $[W_\pi]$.
Every one of these pushforwards is zero by
Lemma~\ref{ka:lem:transferzero}. Thus $\iota_*=0$, and restriction is an
isomorphism. The schemes $X$ and $Y$ are smooth of the same dimension;
restriction respects their intersection products. Their character lines
also correspond under restriction. Proposition~\ref{ka:prop:punctured}
therefore gives the claimed ring with precisely the stated generator.
\end{proof}

\subsection{What the calculation says in elementary terms}
The full diagonal has normal representation consisting of all nontrivial
characters, each $d$ times. Its Euler class gives the first relation
$t^N=0$. The remaining equality strata have no cyclic stabilizer, and
their classes disappear through ordinary affine-space rational
equivalences before they reach the quotient. They therefore contribute
no smaller relation.

It follows that
\begin{equation}
 \CH^r(X_{p,d})=
 \begin{cases}
  \Z,&r=0,\\
  \Z/p\Z,&1\leq r<N,\\
  0,&r\geq N.
 \end{cases}
 \label{ka:eq:gradedgroups}
\end{equation}
In particular $\CH^*(X)\otimes\Q=\Q$. A calculation with rational
Chow groups would have lost \emph{all} of the positive-degree evidence
for the atlas lower bound. When $N>1$, smoothness gives
$\Pic(X)=\CH^1(X)\simeq\Z/p\Z$, recovering the original report's
one-dimensional Picard calculation while also determining its higher
products.

\begin{writenote}[Second route to Part~I's Picard group]
Part~I's calculation is Theorem~\ref{rss:thm:picard}: for odd $p$ and $k$
algebraically closed of characteristic zero, by descent from the explicit
unit group (Lemma~\ref{rss:lem:descent}). For $d=1$ and odd $p$ the degree-one
part of \eqref{ka:eq:gradedgroups} proves the same group by a different
route (localization, the scalar approximation and the Euler class), over any
field of characteristic different from $p$ that contains $\zeta$. Both proofs
are printed. At $(p,d)=(2,1)$, $N=1$ and $\CH^1=0$; this case lies outside
Part~I's hypothesis that $p$ is odd.
\end{writenote}

\section{The exact number of regular Kummer charts}
\label{ka:sec:minimum}

\subsection{A cover-nilpotence lemma}
The lower bound works for arbitrary covering opens, not merely opens
cut out by global sections. We isolate the precise fact needed.

\begin{lemma}[A trivializing cover kills a Chern power]
\label{ka:lem:covernilpotence}
Let $S$ be a smooth pure-dimensional variety and let $M$ be a line
bundle. If $S=W_1\cup\cdots\cup W_m$ and $M|_{W_i}$ is trivial for
each $i$, then
\begin{equation}
 c_1(M)^m=0\quad\text{in }\CH^m(S).
 \label{ka:eq:covernilpotence}
\end{equation}
\end{lemma}
\begin{proof}
Let $D_i=S\setminus W_i$ and let $n=\dim S$.
We construct a cycle class supported on
$D_1\cap\cdots\cap D_r$ whose pushforward is
$c_1(M)^r\cap[S]$. Start with $[S]$ when $r=0$.
Suppose $\beta$ on $D_1\cap\cdots\cap D_{r-1}$ has the required
pushforward. Apply the first Chern class of the restriction of $M$ to
$\beta$. On the intersection with $W_r$, this line bundle is trivial,
so the resulting class restricts to zero. Localization lifts it to a
class on $D_1\cap\cdots\cap D_r$. The projection formula shows that
its pushforward to $S$ is $c_1(M)^r\cap[S]$.

For $r=m$, the support is empty because the $W_i$ cover $S$.
The resulting class is zero. Identifying homological and cohomological
Chow groups on the smooth variety $S$ proves the assertion.
\end{proof}

This proof uses Chern-class operations on the possibly singular closed
intersections $D_1\cap\cdots\cap D_r$; it does not assume that the
boundary divisors meet properly, or even that those closed sets are
divisors.

\begin{proposition}[Universal lower bound for the atlas]
\label{ka:prop:atlaslower}
Every regular Kummer-chart cover of $X_{p,d}$ has at least $N$ members.
\end{proposition}
\begin{proof}
By Lemma~\ref{ka:lem:kummer}, every chart trivializes $\LL$, regardless
of its chosen nontrivial character. A cover with $m<N$ members would
force $t^m=0$ by Lemma~\ref{ka:lem:covernilpotence}, contrary to
\eqref{ka:eq:gradedgroups}. For $N=1$, a nonempty scheme cannot be
covered by zero opens, which is the required lower bound.
\end{proof}

\subsection{Fourier charts attain the bound}
Write $x_{a,i}$ for the $a$th coordinate of $z_i$, where
$1\leq a\leq d$ and $i\in\F_p$. Choose the action on functions so that
$\sigma(x_{a,i})=x_{a,i+1}$, and define
\begin{equation}
 u_{a,j}=\sum_{i=0}^{p-1}\zeta^{-ij}x_{a,i}
 \qquad (j\in\F_p).
 \label{ka:eq:fourier}
\end{equation}
A change of index and the finite geometric-series identity give
\begin{equation}
 \sigma(u_{a,j})=\zeta^j u_{a,j},\qquad
 x_{a,i}=\frac1p\sum_{j=0}^{p-1}\zeta^{ij}u_{a,j}.
 \label{ka:eq:fourierinverse}
\end{equation}
Thus $u_{a,j}$ is a global section of $\LL_j$. Let
\[
 a_{a,j}=u_{a,j}^p\in\Gamma(X,\OO_X),\qquad
 D_{a,j}=\{a_{a,j}\neq0\}\subset X,
 \quad 1\leq j\leq p-1.
\]
There are exactly $d(p-1)=N$ such opens.

\begin{proposition}[Optimal Fourier atlas]
\label{ka:prop:fourieratlas}
The $D_{a,j}$ cover $X$ and are regular Kummer charts, with
\begin{equation}
 (q_*\OO_U)|_{D_{a,j}}
 \simeq\OO_{D_{a,j}}[T]/(T^p-a_{a,j}),
 \qquad T\longmapsto u_{a,j}.
 \label{ka:eq:fourierchart}
\end{equation}
\end{proposition}
\begin{proof}
If all nonconstant modes vanish at a geometric point, Fourier inversion
makes each coordinate $x_{a,i}$ independent of $i$. All $p$ points are
then equal, which is impossible in $U$. Hence the displayed opens cover.
On $D_{a,j}$, $u_{a,j}^p$ is a unit. Consequently $u_{a,j}$ is a unit
on the inverse image, and it generates the line bundle $\LL_j$.
Lemma~\ref{ka:lem:kummer} gives \eqref{ka:eq:fourierchart}.
\end{proof}

Propositions~\ref{ka:prop:atlaslower} and \ref{ka:prop:fourieratlas} prove
$\atlas(X)=N$. Since the lower bound already covers all open sets and
all nonlinear character sections, no further argument about special
Fourier supports is needed for global optimality.

\section{Minimal generators and nonlinear covariant noncompression}
\label{ka:sec:generators}

\subsection{Every nontrivial eigenline needs exactly \texorpdfstring{$N$}{N} sections}
Fix $r\in\F_p^\times$. For $j\in\F_p^\times$, let $e_j(r)$ be the
unique integer in $\{1,\ldots,p-1\}$ with
\begin{equation}
 j e_j(r)\equiv r\pmod p.
 \label{ka:eq:exponents}
\end{equation}
Then
\begin{equation}
 s_{a,j}^{(r)}=u_{a,j}^{e_j(r)}\in\Gamma(X,\LL_r),
 \qquad 1\leq a\leq d,\quad1\leq j\leq p-1.
 \label{ka:eq:fixedweightsections}
\end{equation}
Their common zero locus is the same as that of all the nonconstant
Fourier modes, hence is empty on $X$. They give a surjection
$\OO_X^N\twoheadrightarrow\LL_r$.

Conversely, $m$ global generators of $\LL_r$ produce a cover by their
nonvanishing opens, on each of which $\LL_r$ and hence $\LL$ is trivial.
Proposition~\ref{ka:prop:atlaslower} gives $m\geq N$. This proves
$\gen(\LL_r)=N$ and completes Theorem~\ref{ka:thm:main}.

There is also a useful direct Chern-class proof of this lower bound.
If $\OO_X^m\twoheadrightarrow\LL_r$ has kernel $K$, then $K$ is a
vector bundle of rank $m-1$ and
\[
 c(K)=\frac{1}{1+r t}.
\]
Its degree-$m$ component is $(-r t)^m$ and must vanish. For $m<N$,
this contradicts \eqref{ka:eq:gradedgroups}, since $r^m$ is prime to $p$.

\begin{example}[Four quintic generators of one fixed character]
For $p=5$, $d=1$, the weight-one eigenmodule is generated by
\[
 u_1,\qquad u_2^3,\qquad u_3^2,\qquad u_4^4.
\]
Indeed, the weights of all four expressions are $1$ modulo $5$.
No three elements of this invertible $B$-module can generate it,
regardless of their polynomial degrees or their allowed denominators in
$A$. Thus the result concerns module generation, not just the inability
to omit one entry from this particular list.
\end{example}

\subsection{Unavoidable common exceptional points}
\begin{theorem}[Sharp nonlinear character noncompression]
\label{ka:thm:noncompression}
Let $1\leq m<N$ and let
$f_1,\ldots,f_m\in\Gamma(U_{p,d},\OO_U)$ be regular functions with
\[
 \sigma(f_i)=\zeta^{j_i}f_i,\qquad j_i\in\F_p^\times.
\]
Then the $f_i$ have a common zero on $U_{p,d}$ after passage to an
algebraic closure of $k$. For $m=N$ there is a family without a common
zero, namely the nonconstant Fourier modes.
\end{theorem}
\begin{proof}
Each $f_i$ is a section of $\LL_{j_i}$ on $X$. If they had no common
geometric zero, their tuple would be a nowhere-vanishing section of
\[
 E=\bigoplus_{i=1}^m\LL_{j_i}.
\]
It would define a trivial line subbundle of $E$, with a locally free
quotient of rank $m-1$. The top Chern class of $E$ would therefore be
zero. But
\[
 c_m(E)=(j_1\cdots j_m)t^m\neq0,
\]
by \eqref{ka:eq:gradedgroups}. The contradiction proves the first assertion.
The Fourier family proves sharpness.
\end{proof}

Equivalently, there is no equivariant regular map from $U_{p,d}$ to the
punctured affine representation with fewer than $N$ nontrivial character
coordinates. The target coordinate weights may be chosen arbitrarily.
In particular, increasing the nonlinear degree of the covariants does
not remove this obstruction. In the affine one-dimensional case,
regular functions already include expressions with poles along the
excluded discriminant; allowing those denominators does not change the
result.

This is a more precise statement than saying that familiar radicals
sometimes vanish. Any list that is too short has a common exceptional
configuration, even when the list is designed by a completely different
algebraic construction.

\section{Explicit local reconstruction and coherent radical branches}
\label{ka:sec:reconstruction}

The chart optimum does not by itself explain how roots are recovered.
We record the actual local algebra, keeping the denominator and branch
conditions visible. For $d=1$, this is the Fourier--Kummer construction
already developed in \cite{ka:ProveItReport}; its role here is to provide
the upper bound and to connect the obstruction with radical solvers.

\begin{writenote}[Relation to Part~I]
For $d=1$, Proposition~\ref{ka:prop:branches} below is
Theorem~\ref{rss:thm:chart}, with \eqref{ka:eq:allbranchshift} equal to
\eqref{rss:eq:branchshift}, and \eqref{ka:eq:transition} is Part~I's branch
transition \eqref{rss:eq:transition}; Part~I also gives the coefficient rule
\eqref{rss:eq:transitioncoeff}, which is not restated here. They are printed
because they are stated for every $d$. Index caution: in \eqref{ka:eq:ratios}
the first index $a'$ of $c_{a',k}$ is a \emph{coordinate} index and the pivot
$u_{b,j}$ is fixed, whereas the first index of Part~I's $c_{j,k}$ is the
\emph{pivot}. True reading for $d=1$: Part~II's $c_{1,k}$ with pivot $u_{1,j}$
is Part~I's $c_{j,k}$. False reading: Part~II's $c_{1,k}$ is Part~I's
$c_{1,k}$ (true only for the pivot $j=1$). In the subsection on changing the
pivot, $c$ is at once a coordinate index ($u_{c,h}$) and a ratio.
Section~\ref{ka:sec:notation} lists the other overloaded letters.
\end{writenote}

Choose a pivot $u_{b,j}\neq0$, and abbreviate
\[
 u=u_{b,j},\qquad a=u^p,\qquad e_j(k)\in\{1,\ldots,p-1\},
 \quad j e_j(k)\equiv k\pmod p.
\]
On the chart $D(a)$ define
\begin{equation}
 c_{a',k}=\frac{u_{a',k}}{u^{e_j(k)}}
 \qquad (1\leq a'\leq d,\ 1\leq k\leq p-1).
 \label{ka:eq:ratios}
\end{equation}
Every denominator is a unit on this chart. The covariance formula gives
\[
 \sigma(c_{a',k})=\zeta^{k-j e_j(k)}c_{a',k}=c_{a',k}.
\]
Thus all the $c_{a',k}$ and $u_{a',0}$ are regular invariant functions
there. Also $c_{b,j}=1$.

\begin{proposition}[One shared radical reconstructs the whole tuple]
\label{ka:prop:branches}
At a geometric point of $D(a)$, choose any root $\rho$ of $T^p-a$ and put
\begin{equation}
 y_{a',i}=\frac1p\left(u_{a',0}+
       \sum_{k=1}^{p-1}\zeta^{ik}c_{a',k}\rho^{e_j(k)}\right).
 \label{ka:eq:reconstruct}
\end{equation}
There is an $r\in\F_p$ such that simultaneously for every $a'$ and $i$,
\begin{equation}
 y_{a',i}=x_{a',i+rj^{-1}}.
 \label{ka:eq:allbranchshift}
\end{equation}
In particular, all branches produce the same unordered configuration.
\end{proposition}
\begin{proof}
Since $u\neq0$, all the roots of $T^p-u^p$ are
$\rho=\zeta^r u$, with $r\in\F_p$. Substitution gives
\[
 c_{a',k}\rho^{e_j(k)}
 =\zeta^{r e_j(k)}u_{a',k}
 =\zeta^{rj^{-1}k}u_{a',k}.
\]
Fourier inversion proves \eqref{ka:eq:allbranchshift}. The same $r$ works
for all coordinates because there was only one radical choice.
\end{proof}

\subsection{Changing the pivot}
Suppose a second pivot $u_{c,h}$ is also nonzero. On the overlap set
\[
 e=e_j(h),\qquad c=\frac{u_{c,h}}{u_{b,j}^{e}},
 \qquad a_{c,h}=c^p a_{b,j}^{e}.
\]
A branch $\rho_{b,j}$ determines the compatible branch
\begin{equation}
 \rho_{c,h}=c\rho_{b,j}^{e}.
 \label{ka:eq:transition}
\end{equation}
If $\rho_{b,j}=\zeta^r u_{b,j}$, then
$\rho_{c,h}=\zeta^{re}u_{c,h}$. Its reconstructed shift is
$re h^{-1}=rj^{-1}$ modulo $p$, so the ordered reconstructed tuple is
unchanged. This supplies exact transition data rather than an
independent choice of a new radical on each chart.

\subsection{What the local data do not supply}
The formula assumes the invariant functions in \eqref{ka:eq:ratios} are
already available in the relevant cyclic quotient or intermediate field.
It is not an algorithm for computing them from symmetric polynomial
coefficients. Indeed, for $d=1$ the field $\operatorname{Frac}(B)$ is
larger than the symmetric coefficient field
$k(x_0,\ldots,x_{p-1})^{S_p}$. The extension with group $C_p$ appears
only after this change of base. Nothing here bypasses the Galois
obstruction to solving a generic degree-$p$ polynomial by radicals over
the original symmetric coefficient field.

The theorem counts charts on that specified cyclic base. On a single
chart only one $p$th radical is used. Having to cover the entire base
with $p-1$ charts does not mean that $p-1$ independent $p$th radicals
are needed at any individual input.

\section{Ordinary algebra generators: a different exact complexity}
\label{ka:sec:algebragenerators}

The number of Kummer charts is not the number of generators of the
finite algebra as an ordinary algebra. We can determine the latter
number as well, and the two answers are markedly different.

\begin{definition}
Let $\alggen(q_*\OO_U)$ be the least number of global sections
$y_1,\ldots,y_s$ that generate $q_*\OO_U$ as an $\OO_X$-algebra sheaf.
Generation is local on the base, but the chosen sections themselves are
global. These generators are not required to have character weights.
\end{definition}

\begin{theorem}[No regular equivariant configuration compression]
\label{ka:thm:configurationcompression}
If $d>s\geq1$, there is no $C_p$-equivariant regular morphism
\begin{equation}
 U_{p,d}\longrightarrow U_{p,s}.
 \label{ka:eq:compression}
\end{equation}
If $d\leq s$, such a morphism exists by adjoining zero coordinates to
each point.
\end{theorem}
\begin{proof}
An equivariant morphism between the free spaces would induce a morphism
$f:X_{p,d}\to X_{p,s}$ pulling back the target torsor, and hence its
character line, to the source torsor and character line. Therefore
$f^*(t_{p,s})=t_{p,d}$. Pullback on the Chow rings of the smooth varieties
would imply
\[
 t_{p,d}^{s(p-1)}=f^*(t_{p,s}^{s(p-1)})=0.
\]
But $0<s(p-1)<d(p-1)$, so the left side is nonzero by
Theorem~\ref{ka:thm:main}. This is a contradiction. The coordinate inclusion
when $d\leq s$ visibly preserves distinctness and cyclic permutations.
\end{proof}

\begin{theorem}[Exact global algebra-generator number]
\label{ka:thm:alggen}
For the universal prime cyclic configuration cover,
\begin{equation}
 \alggen(q_*\OO_U)=d.
 \label{ka:eq:alggen}
\end{equation}
\end{theorem}
\begin{proof}
For the upper bound, take $x_{1,0},\ldots,x_{d,0}$, the coordinates
of the point with label $0$. On a geometric fiber of $q$, their
simultaneous values are precisely the $p$ distinct points
$z_0,\ldots,z_{p-1}$. Polynomial functions in these $d$ coordinates
separate those points and generate the full product algebra of functions
on the fiber, by the Chinese remainder theorem for their distinct
maximal ideals.

To pass from fibers to sheaf generation, note that each multiplication
operator by $x_{a,0}$ has a monic characteristic polynomial of degree
$p$. By Cayley--Hamilton, the monomials
$\prod_a x_{a,0}^{e_a}$ with $0\leq e_a<p$ span the subalgebra generated
by the chosen sections. They define a morphism from a finite free sheaf
to $q_*\OO_U$ that is surjective on every geometric fiber. Its coherent
cokernel is zero by Nakayama's lemma. Thus $d$ generators suffice.

For the lower bound, suppose $s$ global functions $y_1,\ldots,y_s$ on
$U$ generate the algebra. On each geometric fiber their simultaneous
values must distinguish all $p$ points: otherwise every polynomial in
them would take equal values at two of the factors of the product
algebra. The functions
\[
 (\sigma^i y_1,\ldots,\sigma^i y_s),\qquad i\in\F_p,
\]
therefore define an equivariant regular morphism
$U_{p,d}\to U_{p,s}$. Theorem~\ref{ka:thm:configurationcompression} forces
$s\geq d$. Zero generators cannot generate a rank-$p>1$ algebra, so
this also covers the remaining possibility $s=0$.
\end{proof}

\subsection{The one-dimensional cover is globally monogenic}
For $d=1$, define the invariant polynomial
\begin{equation}
 f(T)=\prod_{i=0}^{p-1}(T-x_i)\in B[T].
 \label{ka:eq:ordinarypoly}
\end{equation}
There is an isomorphism
\begin{equation}
 B[T]/(f(T))\xrightarrow{\sim} A,
 \qquad T\longmapsto x_0.
 \label{ka:eq:monogenic}
\end{equation}
For a direct proof, base change along the faithfully flat map $B\to A$.
All pairwise differences $x_i-x_j$ are units in $A$, so the Chinese
remainder theorem identifies the source after base change with $A^p$.
The torsor identity identifies $A\otimes_B A$ with the same product,
and the map evaluates $T$ at the $p$ cyclic translates of $x_0$.
It is an isomorphism after this base change, hence before it.

Thus the original cover has one global \emph{ordinary} polynomial
coordinate while requiring $p-1$ regular \emph{Kummer} charts. The
obstruction is to the special equation $T^p-a$ with an eigen-coordinate,
not to every possible global primitive element.

\subsection{The regular-representation bundle detects the difference}
There is also an informative Chern calculation. Let $E=q_*\OO_U$.
Since $E=\bigoplus_{j=0}^{p-1}\LL_j$, the identity
$\prod_{j=1}^{p-1}(1+jz)=1-z^{p-1}$ over $\F_p$ gives, integrally
in our Chow ring,
\begin{equation}
 c(E)=1-t^{p-1}.
 \label{ka:eq:regularbundle}
\end{equation}
The intermediate coefficients vanish because they are divisible by $p$
and every positive power of $t$ is $p$-torsion. For $d\geq2$,
$t^{p-1}\neq0$, so $E$ is not a free vector bundle. In particular,
it cannot have a global monic one-variable polynomial presentation.
For $d=1$, that Chern class vanishes, consistently with the explicit
free basis $1,x_0,\ldots,x_0^{p-1}$ in \eqref{ka:eq:monogenic}.
This Chern observation alone proves non-monogenicity for $d\geq2$;
Theorem~\ref{ka:thm:alggen} supplies the exact stronger value $d$.

\section{Allowing partial collisions does not improve the atlas}
\label{ka:sec:robustness}

The deletion theorem has a useful consequence that goes beyond the
pairwise-distinct configuration space.

\begin{theorem}[Intermediate-open invariance]
\label{ka:thm:sandwich}
Let $W\subset V_{p,d}$ be any $C_p$-invariant open satisfying
\begin{equation}
 U_{p,d}\subseteq W\subseteq V_{p,d}\setminus F.
 \label{ka:eq:sandwich}
\end{equation}
Then, for $Z=W/C_p$ and its generating character line $\LL_Z$,
\begin{equation}
 \CH^*(Z)=\Z[t_Z]/(pt_Z,t_Z^N),\qquad
 \atlas(Z)=\gen((\LL_Z)_j)=N\quad(j\neq0).
 \label{ka:eq:sandwichconclusions}
\end{equation}
\end{theorem}
\begin{proof}
There are open immersions $X\subseteq Z\subseteq Y$. Both restriction
maps in
\[
 \CH^*(Y)\longrightarrow\CH^*(Z)\longrightarrow\CH^*(X)
\]
are surjective by localization. Their composite is an isomorphism by
Theorem~\ref{ka:thm:deletion}. The first map must therefore be injective
as well, and then the second is an isomorphism. Character classes
restrict compatibly, giving the ring formula.

The cover-nilpotence argument gives the lower bound $N$ for Kummer
charts. The same $N$ Fourier modes have no common zero anywhere outside
$F$, so they supply the upper bound on $Z$. Their powers of a fixed
nonzero weight similarly give the exact eigenline generator number.
\end{proof}

\begin{corollary}[Cyclic graph configurations]
\label{ka:cor:graphs}
Let $\Gamma$ be any nonempty simple graph on $\F_p$ invariant under
translation. Impose only the inequalities $z_i\neq z_j$ for edges
$\{i,j\}$ of $\Gamma$. The resulting free cyclic quotient has the same
Chow ring and the same optimal Kummer-chart number $d(p-1)$.
\end{corollary}
\begin{proof}
The open in question contains the pairwise-distinct configuration and
excludes the full diagonal, because the graph has an edge. Apply
Theorem~\ref{ka:thm:sandwich}.
\end{proof}

For example, imposing only consecutive cyclic inequalities
$z_i\neq z_{i+1}$ already has the same obstruction as imposing every
pairwise inequality. Thus adding a large family of allowed partial
collisions does not create a smaller global Kummer atlas.

The direction of the hypothesis is essential. If one deletes additional
points from $U$, for instance by keeping only $D(u_{a,j}^p)$, a single
Kummer chart does work. Our theorem concerns intermediate opens
\emph{containing all distinct configurations}; it does not assert
invariance under arbitrary further deletions. Likewise, the ordinary
algebra-generator conclusion \eqref{ka:eq:alggen} is not asserted for these
larger opens, since the point-coordinate values may now repeat.

\section{Explicit torsion cycles and small-degree examples}
\label{ka:sec:cycles}

\subsection{Geometric representatives of every nonzero power}
The obstruction is not merely a formal symbol in a presentation. Its
powers admit very concrete cycle representatives.

\begin{proposition}[Fourier-coordinate cycles]
\label{ka:prop:cycles}
Let $S$ be a subset of the $N$ index pairs $(a,j)$ with $j\neq0$, and
write $r=|S|<N$. The equations $u_{a,j}=0$ for $(a,j)\in S$ define a
nonempty smooth integral closed subvariety $Z_S$ of $X_{p,d}$ of
codimension $r$. Its Chow class is
\begin{equation}
 [Z_S]=\left(\prod_{(a,j)\in S}j\right)t^r.
 \label{ka:eq:cycles}
\end{equation}
For $1\leq r<N$ it is a generator of $\CH^r(X)\simeq\Z/p\Z$.
\end{proposition}
\begin{proof}
The Fourier transform is an invertible linear change of coordinates on
$V$. Thus these equations define a linear subspace of codimension $r$.
Its intersection with $U$ is a smooth integral open subset. It is
nonempty: retain any one of the nonconstant modes not in $S$, set it
to a nonzero value, and set the other nonconstant modes to zero.
Fourier inversion gives $p$ distinct points, since a nonzero character
of $C_p$ takes $p$ distinct values.

The intersection descends to a smooth integral closed subvariety of
$X$, and regular-embedding codimension can be checked after the \'etale
cover $U\to X$. Its equations are the components of a regular section
of $\bigoplus_{(a,j)\in S}\LL_j$. Its fundamental class is the top
Chern class of this bundle, which is \eqref{ka:eq:cycles}. The coefficient
is nonzero modulo $p$, proving the last assertion.
\end{proof}

\subsection{A pure-mode locus represents the sharp obstruction}
Take $S$ to contain every nonconstant index except one, $(b,h)$.
The resulting locus has all nonconstant modes zero except
$u_{b,h}\neq0$. Before quotienting it is $\A^d\times\mathbb G_m$,
and the cyclic group scales the second factor with a nontrivial weight.
Its quotient is again $\A^d\times\mathbb G_m$, with coordinate
$u_{b,h}^p$ on the multiplicative factor. For $N>1$, this very simple
subvariety represents a nonzero multiple of $t^{N-1}$, of exact order
$p$ in the ambient Chow group.

When $d=1$, its root configurations are
\[
 x_i=\eta+\beta\zeta^{hi},\qquad \beta\neq0,
\]
with root polynomial
\begin{equation}
 \prod_i(T-x_i)=(T-\eta)^p-\beta^p.
 \label{ka:eq:binomiallocus}
\end{equation}
Thus the highest nonzero obstruction power is represented by a locus of
translated binomial configurations. The cover itself is easy to write
on that locus; its nonzero \emph{ambient cycle class} is what obstructs
covering the whole space with too few charts. There is no contradiction
between these two statements.

\subsection{Numerical values and the exceptional smallest case}
\begin{center}
\begin{tabular}{@{}ccclcc@{}}
\toprule
$p$ & $d$ & $N$ & Integral Chow ring & Kummer charts & Algebra generators\\
\midrule
2 & 1 & 1 & $\Z$ & 1 & 1\\
3 & 1 & 2 & $\Z[t]/(3t,t^2)$ & 2 & 1\\
5 & 1 & 4 & $\Z[t]/(5t,t^4)$ & 4 & 1\\
7 & 1 & 6 & $\Z[t]/(7t,t^6)$ & 6 & 1\\
3 & 2 & 4 & $\Z[t]/(3t,t^4)$ & 4 & 2\\
5 & 2 & 8 & $\Z[t]/(5t,t^8)$ & 8 & 2\\
\bottomrule
\end{tabular}
\end{center}
The last column concerns the distinct-point cover only, as in
Theorem~\ref{ka:thm:alggen}. The chart count equals the eigenline generator
count, not the ordinary algebra-generator count.

For $p=2,d=1$, the single nonconstant Fourier mode is $x_0-x_1$, which
is already a unit on the entire configuration space. It gives the global
Kummer coordinate predicted by $N=1$.

For a quintic, $t^3\neq0$ but $t^4=0$. The product coefficient for the
pure weight-one locus is $2\cdot3\cdot4=24\equiv-1\pmod5$, so that
locus represents $-t^3$. Three proposed nonlinear charts cannot cover
$X_{5,1}$, because their trivializations would force this class to vanish.
For a septic, the analogous obstruction is $t^5\neq0$; the exact optimum
is six, not merely a bound between two and six.

\section{Proof audit, reproducible checks, and formalization}
\label{ka:sec:audit}

\subsection{The logical dependency boundary}
The main result is not inferred from finite experiments. Its proof has
the following dependency structure.
\begin{center}
\small
\begin{tabularx}{\textwidth}{@{}>{\raggedright\arraybackslash}p{0.25\textwidth}Y>{\raggedright\arraybackslash}p{0.22\textwidth}@{}}
\toprule
Conclusion & Essential input & Evidence supplied\\
\midrule
Punctured-quotient Chow ring & Vector-bundle localization, scalar
$C_p$ approximation, Euler class & Written proof,
Proposition~\ref{ka:prop:punctured}\\
Partial-diagonal deletion invariance & Prime cyclic partition lemma,
ordinary Chow groups of open affine spaces, proper pushforward & Written
proof, Theorem~\ref{ka:thm:deletion}\\
Optimal arbitrary-open atlas & Nonzero powers of $c_1(\LL)$,
cover-nilpotence, Fourier cover & Written proof,
Section~\ref{ka:sec:minimum}\\
Eigenline and algebra generators & Explicit sections, torsor fibers,
configuration noncompression & Written proofs,
Sections~\ref{ka:sec:generators} and \ref{ka:sec:algebragenerators}\\
Finite orbit and Fourier identities & Integer and finite-field
arithmetic & Executed Python checks; not kernel verification\\
\bottomrule
\end{tabularx}
\end{center}
The classical Chow-theoretic properties are mathematical inputs, not
newly formalized infrastructure. The finite checker supplies no evidence
for localization, Chern-class functoriality, or the nonvanishing of a Chow
class. Those assertions are established by the written arguments using
standard intersection theory \cite{ka:Fulton,ka:Totaro,ka:EdidinGraham}.

Several tempting but invalid shortcuts have been avoided. We do not
cancel $p$ in a torsion group. We do not infer the nonlinear optimum
from coordinate-cover irredundancy. We do not identify an ordinary
monogenic presentation with a Kummer presentation. We do not confuse
trivializing an eigenline with splitting the torsor. We do not apply the
prime partition lemma in composite degree.

\subsection{Exact checks included in the package}
The finite checker, \code{code/verify_finite_checks.py}, uses only
Python's standard library. It writes \code{data/finite_checks.json}.\footnote{[write]
Shipped as \code{code/02-kummer-atlases-verify_finite_checks.py} and
\code{data/02-kummer-atlases-finite_checks.json}. Its default output path is
\code{data/finite_checks.json} beside the script's parent directory; always
pass \code{--output} as in the note at the end of Appendix~\ref{ka:app:provenance}.}
The recorded run used Python~3.13.5
and completed successfully. It enumerates all set partitions for
$p=2,3,5,7$, checks the cyclic orbits and the two fixed partitions,
and compares the counts with independently computed Stirling numbers.
For the proper partitions, the numbers of cyclic orbits are:
\begin{center}
\begin{tabular}{@{}ccl@{}}
\toprule
$p$ & Total partitions & Proper-partition orbit counts by increasing block number\\
\midrule
2 & 2 & none\\
3 & 5 & $1$\\
5 & 52 & $3,5,2$\\
7 & 877 & $9,43,50,20,3$\\
\bottomrule
\end{tabular}
\end{center}
The program also checks that the proper partition
$\{\{0,2\},\{1,3\}\}$ is fixed by the cyclic action in degree four,
a concrete warning against dropping primality.

The Fourier tests are exact finite-field calculations, not floating-point
approximations. They check inversion, character covariance, the common
zero locus of the nonconstant modes, the fixed-character powers
\eqref{ka:eq:fixedweightsections}, invariant pivot ratios, and all branches
on a selected group of inputs.
\begin{center}
\begin{tabular}{@{}ccccr@{}}
\toprule
$p$ & Field & Input coverage & Tuple checks & Branch checks\\
\midrule
3 & $\F_7$ & exhaustive & $343$ & $384$\\
5 & $\F_{11}$ & exhaustive & $161051$ & $640$\\
7 & $\F_{29}$ & seeded sample & $4096$ & $896$\\
\bottomrule
\end{tabular}
\end{center}
For each row, branch checks use all $p$ branches on 128 nonconstant
inputs; they are not exhaustive branch checks on all tuples.
Additional tests examine 256 vector-valued tuples for each of the six
pairs $(p,d)$ with $p\in\{3,5,7\}$ and $d\in\{2,3\}$, as well as
Euler coefficients for six primes and four dimensions. The report
records precisely which tests were exhaustive and which were sampled.
The sampling seed is fixed for reproducibility.

The general partition lemma and the Fourier identities have already
been proved in the text. These tests are meant to catch normalization,
index, and implementation errors, especially the sign of the Fourier
transform and the branch-shift exponent. They do not replace those
proofs and are not described as formal certificates for the article's
geometric conclusions.

\subsection{A realistic formalization decomposition}
A proof-assistant implementation should separate the elementary and
intersection-theoretic layers. The following are proposed milestones,
not declarations claimed to exist in the repository.

The first layer is finite and algebraic: define the prime cyclic action
on equality partitions; prove Lemma~\ref{ka:lem:partitions}; verify the
Fourier character projectors, inversion, and the pivot branch shift;
and prove that all nonconstant modes vanish exactly on the full
diagonal. These statements need finite group actions, field algebra,
and finite sums, but no Chow theory.

The next layer is geometric: construct the free quotient, its character
line bundles, and the multiplication isomorphisms; prove the Kummer
chart equivalence; and formalize the implication from a finite
trivializing cover to vanishing of a Chern power. The last implication
can be isolated as a generic theorem, independent of configuration
spaces.

The final layer is intersection-theoretic: establish the scalar finite
approximation, its stable range, the Euler-class quotient, and the
zero-pushforward boundary argument. The ordinary Chow proof in
Section~\ref{ka:sec:deletion} is deliberately arranged to avoid building
an equivariant intersection theory for every singular partial-diagonal
union. A sufficient verified interface would consist of localization,
proper pushforward, homotopy invariance, and first/top Chern-class
operations with their required compatibilities.

No audit of the current coverage of these facilities in a particular
proof-assistant library was performed for this manuscript. A file that
postulates them would be a conditional formalization, not a completed
kernel proof. The delivered package contains no Lean or Rocq code.

\section{Further research questions and topics}
\label{ka:sec:questions}

The original nonlinear-cover question is settled in the model specified
here. The questions below concern genuine extensions or refinements of
that model. They are proposed research targets, not assertions that each
has already been posed as an open problem in the published literature.

\begin{question}[Composite-degree cyclic configurations]
For composite $n$, determine the integral Chow ring of
$\Conf_n(\A^d)/C_n$ and the optimal regular Kummer-chart count when
$\charac(k)\nmid n$ and $\mu_n\subset k$.
\end{question}
The action on distinct configurations remains free, but proper equality
partitions can have nontrivial stabilizers. Their quotient strata need
not have trivial positive-codimension Chow groups, so the deletion proof
changes essentially. The degree-four partition in the finite checks is
the first obstruction. An actionable first case is $n=4$, with separate
tracking of the order-two and primitive order-four character pieces.
The prime-degree formula should not be extrapolated without this work.

\begin{question}[Descent without a primitive root of unity]
What is the best replacement for an optimal character-line atlas over
a field that does not contain $\mu_p$? How much of the atlas and its
transition data descends from the cyclotomic extension?
\end{question}
The constant group $C_p$ and the group scheme $\mu_p$ must be
separated. A local binomial presentation is not automatically an
available model for every constant cyclic torsor over a nonsplitting
base. A precise formulation should use Galois descent of the character
pieces or their higher-rank sums, rather than silently reuse
Lemma~\ref{ka:lem:kummer}. The split-base result supplies a lower-bound
test after scalar extension but not a complete descended construction.

\begin{question}[Other permutation groups and solvable towers]
For a subgroup $H\leq S_n$, which products of character Chern classes
survive on $\Conf_n(\A^d)/H$, and what exact local radical-coordinate
complexity do they force along a solvable subgroup series?
\end{question}
The partition filtration still exists, but strata now contribute through
their stabilizers. The natural next targets are prime-degree affine
groups and small solvable imprimitive groups. This problem goes beyond
computing $\Pic$: higher products are needed for sharp cover counts.
A useful deliverable would identify the full image of each stabilizer's
cycle pushforward, with integral coefficients throughout.

\begin{question}[Coefficient-efficient optimal atlases]
Can the optimal $p-1$ charts for $d=1$ be presented with substantially
smaller exact invariant formulas, while preserving all specialization
and branch conditions?
\end{question}
The chart count cannot improve, but expression degree, bit complexity,
intermediate-field representation, and certificate size can. The
quintic and septic pivot ratios are concrete test cases. Any proposed
algorithm should compute its invariant data from the intended input
field rather than assume an ordered root tuple. The geometric lower
bound does not preclude significant algorithmic improvement here.

\begin{question}[Regular versus holomorphic or continuous chart covers]
For $k=\C$, compare the optimal regular Zariski Kummer atlas with
holomorphic and continuous trivializing covers of the associated
character line bundle.
\end{question}
The present proof uses algebraic Chow classes. To transfer it to another
category, one must understand their cycle-class images and the relevant
notion of an open cover. Equality of the algebraic and topological
numbers is not a consequence of this paper. Computing the image of
$t^r$ under realization, beginning with $p=5,d=1$, would clarify which
part of the obstruction is specifically algebraic.

\begin{question}[Beyond Chow groups: what changes under partial collisions?]
The inclusions in \eqref{ka:eq:sandwich} induce isomorphisms of Chow rings.
How do their higher Chow groups, algebraic cobordism groups, or motives
differ?
\end{question}
The removed strata are nonempty, so an isomorphism of ordinary Chow
rings should not be mistaken for an equivalence of all invariants.
The proof identifies explicit boundary strata and zero ordinary
pushforwards, providing a concrete starting point for studying higher
localization maps. A first test is the comparison between consecutive-edge
and complete-graph configurations in degree five.

\begin{question}[Explicit rational equivalences for the obstruction cycles]
Can the relations $p[Z_S]=0$ and the transfer-zero identities be realized
by small, explicit rational-equivalence chains suitable for a verified
cycle checker?
\end{question}
The current proof derives these relations through general Chow
operations. Making the chains explicit could reduce the trusted
geometric interface needed for the first concrete case. The pure-mode
surface representing $t^3$ for quintics is a particularly small target.
Such a chain must also be paired with a genuine nonvanishing argument;
exhibiting $p$-torsion alone does not prove exact order $p$.

\begin{question}[Global generators in other configuration families]
For graph configurations or other free invariant opens, determine the
minimum number of global ordinary algebra generators of the finite
cover, and compare it with the unchanged Kummer-chart count.
\end{question}
On the distinct-point space that number is exactly $d$. When partial
collisions are allowed, the coordinate functions at one label can fail
to distinguish an entire cyclic fiber. The Chow ring remains the same
under \eqref{ka:eq:sandwich}, but that fact alone supplies no matching
upper bound for ordinary algebra generation. This separates a new
constructive problem from the atlas problem already solved here.

\section{Conclusion}
\label{ka:sec:conclusion}
The sharp answer to the repository's nonlinear-cover problem is
$p-1$. Its explanation is the entire sequence of nonzero powers
\[
 t,t^2,\ldots,t^{p-2}
 \quad\text{in}\quad
 \CH^*(\Conf_p(\A^1)/C_p),
\]
not merely the first Picard-group obstruction. The higher-dimensional
formula is $d(p-1)$, and the proof remains valid for every split base
field of characteristic different from $p$.

The essential geometric step is that all proper partial-diagonal strata
have trivial cyclic stabilizer and push forward to zero through the
ordinary Chow groups of an open affine space. Consequently, removing
those strata preserves the characteristic classes of the punctured
regular representation. A general cover-nilpotence lemma converts those
classes into a sharp lower bound, matched by explicit Fourier charts.

The same calculation controls nonlinear covariants, exhibits concrete
torsion cycles, and determines the exact ordinary algebra-generator
number $d$. In particular, a one-dimensional cyclic cover can have one
global ordinary polynomial coordinate while needing $p-1$ regular
Kummer charts. Keeping these different notions of complexity separate
is essential for interpreting what an optimal radical atlas actually
proves.

% [write] The manuscript's Appendix A is printed as Appendix C (Part I has A and B).
\setcounter{section}{2}
\renewcommand{\thesection}{\Alph{section}}
\makeatletter\gdef\theHsection{\Alph{section}}\gdef\Hy@chapapp{\Hy@appendixstring}\makeatother
\section{Provenance, literature comparison, and reproducibility}
\label{ka:app:provenance}

\subsection{Repository snapshot}
The relevant repository text was checked at the full commit
\begin{center}
\code{1ee53d57de253d16cdaad79d1f54bbd95d76d682}.
\end{center}
The selected source is
\begin{quote}\small
\path{SetTheory/Cardinals/docs/reports/galois-theory-and-radicals/}\allowbreak
\path{specialization-safe-radical-solvers/article.tex}.
\end{quote}
Its inspected Git blob is
\code{f8303298d0ed04eace36343ec5208c9dbe127042}.
The question answered here is Section~14, Research Question~14.5 of
that report. Its existing results distinguish the nonexistence of one
global Kummer coordinate, its Picard-group calculation, and minimality
within the coordinate Fourier cover. Our arbitrary-open lower bound
addresses precisely the gap left by those results.

The repository was changing during inspection; the commit above is an
immutable anchor rather than a claim that a moving default branch will
retain the same contents. The adjacent report directory was inspected,
and targeted repository searches for the Kummer/Chow connection were
performed. This was not a mathematical audit of every research archive
or of all concurrent work in the repository.

\subsection{Relationship to the external literature}
Totaro's construction of $\CH^*(BG)$ and his analysis of cyclic products
supply close methodological precedents \cite{ka:Totaro}. The scalar
approximation and Euler-class quotient in
Proposition~\ref{ka:prop:punctured} belong to this established framework.
Edidin--Graham provide the systematic equivariant intersection theory
behind the same calculation \cite{ka:EdidinGraham}; ordinary intersection
theory supplies the localization and Chern-class operations
\cite{ka:Fulton}. The eigenline interpretation of a Kummer torsor is also
standard \cite{StacksKummer}.

The explicit contribution developed here is the combination of the
partial-diagonal zero-pushforward proof with the sharp arbitrary-open
atlas lower bound and its concrete consequences for the specified
repository question. The manuscript deliberately supplies a proof
rather than infer novelty from a search failure. A specialist literature
review could identify prior formulations or stronger existing results;
such a discovery would change the attribution, not the logical role of
the proof as an answer to the repository question.

\subsection{Package contents and rebuild}
The package contains the standalone LaTeX source and compiled PDF, this
manuscript's README, the finite checker, its JSON output, and a source
manifest. The bibliography is embedded in the \texttt{.tex} file; no
BibTeX database or external graphics are needed.

Run the finite checks with
\begin{quote}\ttfamily\small
python3 code/verify\_finite\_checks.py
\end{quote}
and compile with
\begin{quote}\ttfamily\small
pdflatex -interaction=nonstopmode -halt-on-error optimal\_kummer\_atlases.tex\\
pdflatex -interaction=nonstopmode -halt-on-error optimal\_kummer\_atlases.tex
\end{quote}
A third pass may be needed after changes to the table of contents. The
included \texttt{build.sh} uses \texttt{latexmk} when available and
otherwise runs three PDFLaTeX passes. No network access is needed to
run the checker or build the article once the TeX packages are installed.

\begin{writenote}[Shipped names and a safe rerun]
The paragraph above describes the delivered package. In this report the
checker is \code{code/02-kummer-atlases-verify_finite_checks.py}, its recorded
output \code{data/02-kummer-atlases-finite_checks.json}, the source manifest
\code{data/02-kummer-atlases-source_manifest.json}, the build record
\code{data/02-kummer-atlases-build_check.json} (which describes the delivered
22-page PDF, not this report's PDF), and the build script
\code{code/02-kummer-atlases-build.sh}. The manuscript's source, PDF and
README are not shipped; this report is built from \code{article.tex}.
The build script does not work in the shipped layout: it changes to its own
directory \code{code/} and then calls \code{code/verify_finite_checks.py} and
compiles \code{optimal_kummer_atlases.tex}, neither of which exists there.
In the delivered layout it reran the checker in place, overwriting the
recorded \code{data/finite_checks.json}. Run without \code{--output}, the
checker writes \code{data/finite_checks.json} relative to the report
directory, an unshipped file beside the shipped record; on Windows its
\code{write_text} emits CRLF line endings. Run from the report directory,
writing to a scratch directory:
\begin{lstlisting}
py code/02-kummer-atlases-verify_finite_checks.py --output <tmpdir>/finite_checks.json
\end{lstlisting}
At the write this exited with \code{PASS} in about 9~s, and after removing
carriage returns the fresh file equalled
\code{data/02-kummer-atlases-finite_checks.json} except for the field
\code{"python_version"} (3.14.4 at the write, 3.13.5 recorded).
\end{writenote}

\begin{thebibliography}{99}

\bibitem{Lazard2004}
Daniel Lazard,
\emph{Solving Quintics by Radicals},
in \emph{The Legacy of Niels Henrik Abel}, Springer, 2004, pp.~207--225.
DOI: \href{https://doi.org/10.1007/978-3-642-18908-1_6}{10.1007/978-3-642-18908-1\_6}.

\bibitem{FaggalLazard2014}
Mohammed A. Faggal and Daniel Lazard,
\emph{Solving Quintics and Septics by Radicals},
International Journal of Scientific and Research Publications,
4(12), December 2014, pp.~1--10.
\url{https://www.ijsrp.org/research-paper-1214/ijsrp-p3679.pdf}.
Especially Proposition~7 and Section~5.2. The conclusions drawn in
Section~\ref{rss:sec:published} concern separation and distinct factors, not an
unqualified refutation of radical solvability.

\bibitem{StacksKummer}
The Stacks Project Authors,
\emph{Kummer theory}, Section~59.28, Tag~03PK.
\url{https://stacks.math.columbia.edu/tag/03PK}.
Used for the standard line-bundle interpretation of Kummer torsors; the
unit and Picard-group calculations for the present configuration are
proved in this article.
[write] Part~II cites the same section, especially the description of
torsors by invertible sheaves with a trivialized tensor power (consulted
September 2026).

\bibitem{repoReadme}
Vladimir Reshetnikov and contributors,
\emph{ProveIt}, root README, inspected snapshot
\snapshot.
\url{https://github.com/VladimirReshetnikov/ProveIt/blob/e21766d04c2b8a9b2cdba0cd43e563024b3bd1b9/README.md}.

\bibitem{repoReport}
ProveIt project,
\emph{A Lean Transcription of the Corrected Lazard Formula for Solvable
Quintics}, report checkpoint dated 13 August 2026, inspected in the pinned
snapshot.
\url{https://github.com/VladimirReshetnikov/ProveIt/blob/e21766d04c2b8a9b2cdba0cd43e563024b3bd1b9/Algebra/PolynomialFormulas/LazardQuinticFormalization.tex}.

\bibitem{repoCrosswalk}
ProveIt project,
\emph{Lazard paper claim crosswalk}, inspected in the pinned snapshot.
\url{https://github.com/VladimirReshetnikov/ProveIt/blob/e21766d04c2b8a9b2cdba0cd43e563024b3bd1b9/Algebra/PolynomialFormulas/LazardPaperClaimCrosswalk.md}.

\bibitem{repoFourier}
ProveIt project,
\emph{Arbitrary-degree Fourier reconstruction in Lazard's convention},
\code{LazardGeneralFourier.lean}, inspected in the pinned snapshot.
\url{https://github.com/VladimirReshetnikov/ProveIt/blob/e21766d04c2b8a9b2cdba0cd43e563024b3bd1b9/Algebra/PolynomialFormulas/Lean/PolynomialFormulas/LazardGeneralFourier.lean}.

% [write] The four entries below were added when the manuscript was written
% into the collection.  The cited files are identical at the pin and at
% batch 36, so the pinned URLs are used.
\bibitem{repoFourierRocq}
ProveIt project, [write]
\code{Algebra/PolynomialFormulas/Coq/LazardGeneralFourier.v}, the Rocq
counterpart of \code{LazardGeneralFourier.lean}.
\url{https://github.com/VladimirReshetnikov/ProveIt/blob/e21766d04c2b8a9b2cdba0cd43e563024b3bd1b9/Algebra/PolynomialFormulas/Coq/LazardGeneralFourier.v}.

\bibitem{repoSextic}
ProveIt project, [write]
\code{Algebra/PolynomialFormulas/Lean/PolynomialFormulas/SexticSeparatingInvariants.lean}
and \code{Algebra/PolynomialFormulas/Decidability.md}.
\url{https://github.com/VladimirReshetnikov/ProveIt/blob/e21766d04c2b8a9b2cdba0cd43e563024b3bd1b9/Algebra/PolynomialFormulas/Lean/PolynomialFormulas/SexticSeparatingInvariants.lean}.

\bibitem{repoKummer}
ProveIt project, [write]
\code{Algebra/PolynomialFormulas/Coq/LazardCyclicKummerGenerator.v}.
\url{https://github.com/VladimirReshetnikov/ProveIt/blob/e21766d04c2b8a9b2cdba0cd43e563024b3bd1b9/Algebra/PolynomialFormulas/Coq/LazardCyclicKummerGenerator.v}.

\bibitem{repoCounterexamples}
ProveIt project, [write]
\code{LazardAlternatingResolventCounterexample} and
\code{LazardArtinSchreierRadicalCounterexample}, in
\code{Algebra/PolynomialFormulas/Lean/PolynomialFormulas/} (\code{.lean}) and
\code{Algebra/PolynomialFormulas/Coq/} (\code{.v}); status in
\code{Algebra/PolynomialFormulas/LazardPaperClaimCrosswalk.md}, rows at
l.~200 and~202.

% [write] The five entries below are Part II's (batch 52, manuscript 09), as
% delivered, with keys prefixed ka:.  Its Stacks Project entry is merged into
% StacksKummer above.
\bibitem{ka:ProveItReport}
ProveIt research-report collection,
\emph{Specialization-Safe Radical Solvers: Exact Bad Characteristics for
Subset-Sum Resolvents, a Fourier--Kummer Atlas, and a Global Obstruction}.
September 2026, batch 36, manuscript 04. In particular Sections 8, 10,
12, and Research Question 14.5. AI-assisted, unrefereed report.
Pinned source:\par
\url{https://github.com/VladimirReshetnikov/ProveIt/blob/1ee53d57de253d16cdaad79d1f54bbd95d76d682/SetTheory/Cardinals/docs/reports/galois-theory-and-radicals/specialization-safe-radical-solvers/article.tex}.

\bibitem{ka:ProveItLazard}
ProveIt project,
\emph{A Lean transcription of the corrected Lazard quintic formula}.
Repository audit report, inspected September 2026.\par
\url{https://github.com/VladimirReshetnikov/ProveIt/blob/1ee53d57de253d16cdaad79d1f54bbd95d76d682/Algebra/PolynomialFormulas/LazardQuinticFormalization.tex}.

\bibitem{ka:Totaro}
Burt Totaro,
\emph{The Chow ring of a classifying space}.
Preprint, 1998, arXiv:math/9802097. In particular Section 1 and
Sections 7--8 on cyclic products.\par
\url{https://arxiv.org/abs/math/9802097}.

\bibitem{ka:EdidinGraham}
Dan Edidin and William Graham,
\emph{Equivariant intersection theory}.
Preprint, 1996, arXiv:alg-geom/9603008, version 3;
revised version arXiv:alg-geom/9609018.\par
\url{https://arxiv.org/abs/alg-geom/9603008};
\url{https://arxiv.org/abs/alg-geom/9609018}.

\bibitem{ka:Fulton}
William Fulton,
\emph{Intersection Theory}, second edition.
Springer, 1998. DOI: 10.1007/978-1-4612-1700-8.\par
\url{https://doi.org/10.1007/978-1-4612-1700-8}.

\end{thebibliography}
\end{document}
